% =====================================================================
%  TEMA · MODO TABLERO — "clave del profe" a TODO EL ANCHO
%  Rompe a propósito el layout de libro (twoside angosto): página ancha
%  de una sola columna, banner moderno, imágenes libres al lado con el
%  texto rodeándolas (\figinline), gráficos flat TikZ/SVG y TODO RESUELTO.
%
%  Pensado para material que el PROFE resuelve en el tablero (taller/quiz
%  con solución). NO es el estilo Santillana del kit: es intencionalmente
%  más "revista/póster", más aire, más ancho.
%
%  INVOCACIÓN (tras % =====================================================================
%  DESIGN SYSTEM v5 — EDITORIAL ITSTZ
%  Filosofía: texto primero · un solo color de asignatura · sin sombras
%  · sin marcos pesados · tratamientos planos (filete lateral / regla).
%
%  Vocabulario semántico (esto es TODO lo disponible):
%    \section / \subsection      jerarquía
%    definicion{título}          concepto/regla clave (filete color)
%    ejemplo{título}             ejemplo resuelto (numerado, sin caja)
%    procedimiento{título}       pasos (filete neutro) + \paso{n}{texto}
%    nota[título]                aparte único gris (sabías/dato/video/margen)
%    alerta{título}              advertencia ámbar (usar rara vez)
%    practica{título}            ejercicios/hoja/ICFES (regla superior)
%    \figura{archivo}{pie}       imagen real o placeholder automático
%    \desbloquea{lista}          pie discreto de prerrequisitos
%
%  REGLA DE ORO: las secciones (sNN.tex) NO llevan \vspace, color ni
%  minipage de maquetación. Solo contenido semántico. El diseño vive aquí.
% =====================================================================
\documentclass[11pt, letterpaper, twoside]{article}

\usepackage{fontspec}
\usepackage{polyglossia}
\setdefaultlanguage{spanish}
\usepackage[inner=2.1cm, outer=3.7cm, top=2.6cm, bottom=2.3cm,
            marginparwidth=2.55cm, marginparsep=0.45cm,
            headheight=16pt, headsep=18pt, footskip=24pt]{geometry}
\usepackage{microtype}
\usepackage{multicol}
\usepackage{xcolor}
\usepackage{graphicx}
\usepackage{tikz}
\usetikzlibrary{calc, arrows.meta, positioning}
\usepackage{wrapfig}
\usepackage[most]{tcolorbox}
\usepackage{amsmath, amssymb}
\usepackage{siunitx}
\usepackage{fontawesome5}
\usepackage{fancyhdr}
\usepackage{titlesec}
\usepackage{enumitem}
\usepackage{etoolbox}
\usepackage{titletoc}
\usepackage{array}
\usepackage{booktabs}
\usepackage{colortbl}
\usepackage{nicematrix}   % rejilla completa (hvlines) para tablas de hojas

% --- TIPOGRAFÍA -------------------------------------------------------
% Cuerpo: Charter (serif cálida y muy legible). La matemática (Latin
% Modern, serif) combina con ella. Títulos/etiquetas/UI: Inter Display.
\setmainfont{Charter}[
  UprightFont    = {Charter},
  BoldFont       = {Charter Bold},
  ItalicFont     = {Charter Italic},
  BoldItalicFont = {Charter Bold Italic},
  Ligatures = TeX, Numbers = Lining ]
\setsansfont{Inter}[
  UprightFont = Inter-Regular, BoldFont = Inter-SemiBold,
  ItalicFont  = Inter-Italic, Ligatures = TeX, Numbers = Lining ]
\newfontfamily\InterLight{Inter-Light}[ItalicFont=Inter-LightItalic, Ligatures=TeX]
\newfontfamily\InterDisplay{Inter Display}[Ligatures=TeX]
\newfontfamily\InterDisplaySB{Inter Display SemiBold}[Ligatures=TeX]
\newfontfamily\InterDisplayThin{Inter Display Thin}[Ligatures=TeX]
\newfontfamily\InterDisplayBlk{Inter Display Black}[Ligatures=TeX]

% --- RITMO --------------------------------------------------------------
\linespread{1.22}
\setlength{\parindent}{0pt}
\setlength{\parskip}{0.5em}
\newlength{\bloqueskip}\setlength{\bloqueskip}{1.0em}   % aire entre bloques
\setlength{\columnsep}{1.6em}                            % medianil multicols

% =====================================================================
%  COLOR — UNA sola variable de asignatura + neutros. Sin arcoíris.
% =====================================================================
\definecolor{mat} {RGB}{ 30,  64, 120}   % azul profundo
\definecolor{cnat}{RGB}{ 16, 110,  80}   % esmeralda
\definecolor{csoc}{RGB}{165,  75,  45}   % terracota
\definecolor{eco} {RGB}{ 70,  55, 130}   % índigo
\definecolor{alertcolor}{RGB}{180, 95, 10}  % ámbar — solo advertencias

\colorlet{subjectcolor}{mat}             % por defecto
\newcommand{\setsubject}[1]{\colorlet{subjectcolor}{#1}}

% =====================================================================
%  JERARQUÍA TIPOGRÁFICA
% =====================================================================
\setcounter{secnumdepth}{0}
\titleformat{\section}
  {\LARGE\InterDisplaySB\color{subjectcolor}}{}{0pt}{}
  [\vspace{2pt}{\color{subjectcolor!30}\titlerule[0.8pt]}]
\titlespacing*{\section}{0pt}{3.2ex plus 0.6ex}{1.8ex plus 0.3ex}

\titleformat{\subsection}
  {\large\InterDisplaySB\color{black!82}}{}{0pt}{}
\titlespacing*{\subsection}{0pt}{2.0ex plus 0.4ex}{0.5ex}

% --- TABLA DE CONTENIDO -------------------------------------------------
\titlecontents{section}[0em]
  {\addvspace{0.45pc}\sffamily\bfseries\small\color{black!85}}
  {}{}{\titlerule*[0.35pc]{.}\contentspage}
\titlecontents{subsection}[1.4em]
  {\sffamily\footnotesize\color{black!60}}
  {}{}{\titlerule*[0.35pc]{.}\contentspage}

% =====================================================================
%  ENCABEZADO / PIE — editorial limpio (regla fina, sin franjas)
% =====================================================================
\newcommand{\hdrinstituto}{Instituto Técnico Santo Tomás · Zapatoca}
\newcommand{\hdrcapitulo}{}            % se fija en el documento maestro
\pagestyle{fancy}\fancyhf{}
\renewcommand{\headrulewidth}{0pt}
\fancyhead[LE]{\footnotesize\sffamily\color{black!55}\hdrinstituto}
\fancyhead[RO]{\footnotesize\sffamily\color{subjectcolor}\hdrcapitulo}
\fancyhead[LO,RE]{}
\renewcommand{\headrule}{\vspace{2pt}{\color{black!18}\hrule height 0.4pt}}

% --- PIE: número exterior (cuadro de color) + crédito interior + filete ---
\renewcommand{\footrulewidth}{0pt}
\renewcommand{\footrule}{{\color{black!15}\hrule height 0.4pt}\vspace{3pt}}
\newcommand{\pagenumbox}{%
  \tikz[baseline=-3.2pt]{\node[fill=subjectcolor, text=white, inner xsep=7pt,
     inner ysep=3pt, font=\sffamily\bfseries\footnotesize]{\thepage};}}
\fancyfoot[RO,LE]{\pagenumbox}
\fancyfoot[LO,RE]{\footnotesize\sffamily\color{black!45}%
  \hdrcapitulo\ \textcolor{black!22}{·}\ Prof.\ Juan Diego A.\ Prada R.}

% =====================================================================
%  HELPERS
% =====================================================================
% Etiqueta de bloque:  TAG · título   (color del tag configurable)
\newcommand{\blocklabel}[3]{% #1=color  #2=TAG  #3=título
  {\sffamily\bfseries\footnotesize\color{#1}\MakeUppercase{#2}}%
  \ifblank{#3}{}{\,{\color{black!35}\textbf{·}}\,%
    {\sffamily\bfseries\footnotesize\color{black!78}#3}}%
  \par\nopagebreak\vspace{0.4em}}

% --- CAJAS AL MARGEN — personajes históricos e hitos -------------------
% Van en el margen exterior ancho (\marginpar respeta la geometría twoside
% y flota para no chocar). Son contenido semántico: colócalas en el cuerpo,
% cerca del párrafo que comentan, ENTRE párrafos. NO usarlas dentro de
% multicols, tcolorbox, tablas, ni en la misma línea que un \section.
%   \personaje{Nombre}{años · rol}{descripción breve}
%   \hito{año}{qué pasó (breve)}
\setlength{\marginparpush}{8pt}
\newcommand{\margenbox}[1]{\marginpar{%
  \begingroup\raggedright\footnotesize\sffamily
  \setlength{\parskip}{0.18em}\linespread{1.05}\selectfont
  #1\par\endgroup}}
\newcommand{\personaje}[3]{\margenbox{%
  {\color{subjectcolor}\rule{\linewidth}{1.1pt}}\par
  {\InterDisplaySB\scriptsize\color{subjectcolor}#1}\par
  {\fontsize{6.8pt}{8pt}\selectfont\itshape\color{black!50}#2}\par
  \vspace{0.15em}{\scriptsize\color{black!72}#3}}}
\newcommand{\hito}[2]{\margenbox{%
  {\sffamily\bfseries\scriptsize\colorbox{subjectcolor}{\textcolor{white}{\,#1\,}}}\par
  \vspace{0.15em}{\scriptsize\color{black!72}#2}}}

% =====================================================================
%  ENTORNOS — todos planos: filete lateral, sin marco, sin sombra
% =====================================================================

% --- DEFINICIÓN / concepto clave ---------------------------------------
\newtcolorbox{definicion}[1]{
  enhanced, breakable, sharp corners, boxrule=0pt, frame hidden,
  colback=subjectcolor!6!white,
  borderline west={2.2pt}{0pt}{subjectcolor},
  left=5mm, right=4mm, top=3mm, bottom=3mm,
  before skip=\bloqueskip, after skip=\bloqueskip,
  before upper={\blocklabel{subjectcolor}{Definición}{#1}}
}

% --- PROCEDIMIENTO / pasos ---------------------------------------------
\newtcolorbox{procedimiento}[1]{
  enhanced, breakable, sharp corners, boxrule=0pt, frame hidden,
  colback=black!3!white,
  borderline west={2.2pt}{0pt}{black!45},
  left=5mm, right=4mm, top=3mm, bottom=3mm,
  before skip=\bloqueskip, after skip=\bloqueskip,
  before upper={\blocklabel{black!60}{Procedimiento}{#1}}
}
\newcommand{\paso}[2]{%
  \par\vspace{0.45em}\noindent
  \tikz[baseline=-0.6ex]\node[circle, fill=subjectcolor, text=white,
     inner sep=0pt, minimum size=1.45em,
     font=\sffamily\bfseries\footnotesize]{#1};\hspace{0.5em}#2\par}

% Glosa "¿por qué?" bajo un paso — narración en lenguaje cercano
\newcommand{\porque}[1]{\par\vspace{0.1em}\noindent\hspace{2.0em}%
  {\sffamily\footnotesize\itshape\color{subjectcolor}¿por qué?\,}%
  {\sffamily\footnotesize\itshape\color{black!55}#1}\par}

% Operación glosada: cuenta a la izquierda + porqué al margen derecho
% (estilo "comentario al lado del código"). Úsala dentro de ejemplo.
% \opg{ecuación}{porqué}   — primer arg en modo display (mat, cap0)
% \opgt{etiqueta}{porqué}  — primer arg en texto sans bold (física, ciencias)
\newcommand{\opg}[2]{\par\vspace{0.25em}\noindent
  \makebox[0.46\linewidth][l]{$\displaystyle #1$}%
  {\sffamily\footnotesize\itshape\color{black!52}#2}\par}
\newcommand{\opgt}[2]{\par\vspace{0.25em}\noindent
  \makebox[0.46\linewidth][l]{\sffamily\bfseries\color{subjectcolor!80}#1}%
  {\sffamily\footnotesize\itshape\color{black!52}#2}\par}

% --- NOTA / aparte único (sabías · dato · video · margen) --------------
\newtcolorbox{nota}[1][]{
  enhanced, breakable, sharp corners, boxrule=0pt, frame hidden,
  colback=black!4!white,
  borderline west={2.2pt}{0pt}{black!35},
  left=5mm, right=4mm, top=3mm, bottom=3mm,
  before skip=\bloqueskip, after skip=\bloqueskip,
  before upper={\ifblank{#1}{}{\blocklabel{black!55}{Nota}{#1}}}
}

% --- ALERTA / advertencia (única en ámbar, uso escaso) -----------------
\newtcolorbox{alerta}[1]{
  enhanced, breakable, sharp corners, boxrule=0pt, frame hidden,
  colback=alertcolor!7!white,
  borderline west={2.2pt}{0pt}{alertcolor},
  left=5mm, right=4mm, top=3mm, bottom=3mm,
  before skip=\bloqueskip, after skip=\bloqueskip,
  before upper={\blocklabel{alertcolor}{Atención}{#1}}
}
% Sub-ítem de alerta (antes \caso)
\newcommand{\caso}[2]{\par\vspace{0.35em}\noindent
  {\sffamily\bfseries\color{alertcolor!85!black}#1.}\enspace #2\par}

% --- CONTRASTE (la virtud y su exceso → pregunta que cuestiona) --------
%   \contraste{Concepto}{La virtud / lo valioso}{Su exceso / su riesgo}{Pregunta}
\newtcolorbox{contrastebox}{
  enhanced, breakable, sharp corners, boxrule=0pt, frame hidden,
  colback=subjectcolor!5!white,
  borderline west={2.2pt}{0pt}{subjectcolor},
  left=5mm, right=4mm, top=3mm, bottom=3mm,
  before skip=\bloqueskip, after skip=\bloqueskip,
}
\newcommand{\contraste}[4]{%
\begin{contrastebox}
\blocklabel{subjectcolor}{Cuestiona esto}{\faBalanceScale\ #1}
\noindent\begin{minipage}[t]{0.45\linewidth}
  {\sffamily\footnotesize\bfseries\color{subjectcolor!85!black}La virtud}\par
  \vspace{0.2em}#2
\end{minipage}\hfill
{\color{subjectcolor!35}\vrule width 0.4pt}\hfill
\begin{minipage}[t]{0.45\linewidth}
  {\sffamily\footnotesize\bfseries\color{alertcolor!85!black}Su exceso}\par
  \vspace{0.2em}#3
\end{minipage}
\par\vspace{0.6em}
{\color{subjectcolor!25}\titlerule[0.6pt]}\par\vspace{0.4em}
\noindent{\sffamily\itshape\color{subjectcolor}\faArrowRight\ #4}
\end{contrastebox}}

% --- EJEMPLO (sin caja: regla fina + etiqueta numerada) ----------------
\newcounter{ejemplo}[section]
\newenvironment{ejemplo}[1]
  {\par\addvspace{\bloqueskip}\refstepcounter{ejemplo}%
   \noindent{\color{subjectcolor!28}\rule{\linewidth}{0.5pt}}\par\vspace{0.4em}%
   \noindent{\sffamily\bfseries\color{subjectcolor}Ejemplo \theejemplo.}%
   \enspace{\sffamily\bfseries\color{black!78}#1}\par\vspace{0.3em}}
  {\par\addvspace{\bloqueskip}}

% --- PRÁCTICA (ejercicios / hoja / ICFES) ------------------------------
\newtcolorbox{practica}[1]{
  enhanced, breakable, sharp corners, boxrule=0pt, frame hidden,
  colback=white,
  borderline north={1.4pt}{0pt}{subjectcolor},
  left=0mm, right=0mm, top=3mm, bottom=2mm,
  before skip=\bloqueskip, after skip=\bloqueskip,
  before upper={\blocklabel{subjectcolor}{Práctica}{#1}\raggedright}
}

% Ejercicio numerado con nivel opcional [básico|intermedio|reto]
\newcounter{ejer}[section]
\newcommand{\ejer}[2][]{%
  \stepcounter{ejer}\par\vspace{0.5em}\noindent
  {\sffamily\bfseries\color{subjectcolor}\theejer.}\enspace #2%
  \ifblank{#1}{}{\ \textcolor{black!45}{\footnotesize[\textit{#1}]}}\par}

% Niveles = un solo color a tres intensidades (NO nuevos colores)
\newcommand{\nivel}[2]{\par\smallskip\noindent
  \textcolor{subjectcolor!#1}{\small\faCircle}\enspace
  {\sffamily\bfseries\footnotesize\color{black!70}#2}\par}
\newcommand{\basico}{\nivel{40}{Básico}}
\newcommand{\intermedio}{\nivel{70}{Intermedio}}
\newcommand{\reto}{\nivel{100}{Reto}}

% Selección múltiple tipo ICFES — opciones con burbuja (estilo Saber)
\newcounter{pregicfes}[section]
\newcommand{\bubbleletter}[1]{\tikz[baseline=(bc.base)]\node[circle,
  draw=subjectcolor!55, line width=0.6pt, inner sep=1.5pt,
  font=\sffamily\bfseries\scriptsize, text=subjectcolor!90!black](bc){#1};}
\newcommand{\icfes}[5]{%
  \stepcounter{pregicfes}\par\vspace{0.6em}\noindent
  {\sffamily\bfseries\color{subjectcolor}\thepregicfes.}\enspace #1\par
  \begin{enumerate}[label={\protect\bubbleletter{\Alph*}}, leftmargin=1.5cm,
                    itemsep=3pt, topsep=3pt]
    \item #2 \item #3 \item #4 \item #5
  \end{enumerate}}
% Pregunta abierta de argumentación
\newcommand{\pregabierta}[1]{\par\vspace{0.6em}\noindent
  {\sffamily\bfseries\color{subjectcolor}\faComment\ Argumenta:}\enspace
  \textit{#1}\par\lineaescritura\lineaescritura\lineaescritura}

% --- LÍNEAS DE RESPUESTA -----------------------------------------------
\newcommand{\linea}[1][7cm]{\underline{\hspace{#1}}}
\newcommand{\lineaescritura}{\par\vspace{0.9em}\noindent
  {\color{black!30}\rule{\linewidth}{0.3pt}}}

% =====================================================================
%  EVALUACIÓN TIPO ICFES — cierre de unidad (banner + 2 columnas)
%  Uso:  \begin{evaluacion} \begin{multicols}{2} \icfes... \end{multicols}
%        \pregabierta{...} \end{evaluacion}
% =====================================================================
\newtcolorbox{evaluacion}[1][cierre de sección]{
  enhanced, breakable, sharp corners,
  colback=white, colframe=subjectcolor, boxrule=1pt,
  colbacktitle=subjectcolor, coltitle=white,
  fonttitle=\InterDisplaySB\normalsize,
  title={\faClipboardCheck\ \ EVALUACIÓN TIPO ICFES\ \ %
         \textnormal{\sffamily\small· #1}},
  toptitle=2mm, bottomtitle=2mm, left=4mm, right=4mm, top=3.5mm, bottom=3mm,
  before={\clearpage},
  before upper={\small\itshape\color{black!55}%
    Selección múltiple. Marca una sola opción por pregunta.\par\medskip\raggedright}
}

% =====================================================================
%  HOJA DE TRABAJO — recortable: borde punteado, escala de grises,
%  en página aparte (clearpage antes; el contenido que sigue continúa).
% =====================================================================
\newtcolorbox{hoja}[1]{
  enhanced, breakable, sharp corners,
  colback=black!2!white, boxrule=0pt, frame hidden,
  borderline={0.8pt}{0pt}{black!50, dash pattern=on 3pt off 2.5pt},
  before={\clearpage}, after={\par},
  left=6mm, right=6mm, top=6mm, bottom=6mm,
  before upper={%
    {\sffamily\bfseries\footnotesize\color{black!60}\faCut\ \ HOJA DE TRABAJO}%
    {\color{black!35}\ \textbf{·}\ }{\sffamily\bfseries\footnotesize\color{black!80}#1}%
    \hfill{\sffamily\footnotesize\color{black!55}%
      Nombre:\ \underline{\hspace{4.2cm}}\quad Fecha:\ \underline{\hspace{2cm}}}\par
    \vspace{2mm}{\color{black!30}\hrule height 0.4pt}\vspace{3mm}\raggedright}
}

% =====================================================================
%  TABLAHOJA — tabla con TODAS las líneas (rejilla) para hojas de trabajo.
%  Celdas altas para llenar a mano. Centrada. El encabezado se marca con
%  \rowcolor en la primera fila.
%  Uso:  \begin{tablahoja}{ccc} enc & enc & enc \\ a & b & c \\ \end{tablahoja}
% =====================================================================
\NewDocumentEnvironment{tablahoja}{m}{%
  \par\smallskip\centering\small\setlength{\tabcolsep}{4pt}%
  \begin{NiceTabular}{#1}[hvlines, cell-space-top-limit=5pt,
     cell-space-bottom-limit=9pt, colortbl-like]%
}{%
  \end{NiceTabular}\par\smallskip
}

% =====================================================================
%  FIGURAS — imagen real o placeholder automático (\IfFileExists)
% =====================================================================
\newcounter{figura}[section]
\newcommand{\figura}[3][0.7\linewidth]{% [ancho]{archivo}{pie}
  \par\addvspace{\bloqueskip}\refstepcounter{figura}%
  \begin{center}%
  \IfFileExists{assets/imgs/#2}{%
    \includegraphics[width=#1]{assets/imgs/#2}}{%
    \begin{tikzpicture}
      \node[draw=black!30, dashed, line width=0.6pt, fill=black!3,
            text width=0.66\linewidth, minimum height=2.6cm,
            align=center, inner sep=8pt,
            text=black!55, font=\sffamily\footnotesize]
        {{\Large\faImage}\\[4pt]\textbf{[ imagen pendiente ]}\\#3};
    \end{tikzpicture}}%
  \par\vspace{0.4em}%
  {\footnotesize\sffamily\color{black!60}%
    \textbf{Figura \thefigura.}\enspace #3}%
  \end{center}\par\addvspace{\bloqueskip}}

% =====================================================================
%  APERTURA DE SECCIÓN — marco branded de serie + título + "qué verás"
%  El "marco" es CONSTANTE en todas las secciones (= identidad de marca):
%  número grande de serie, kicker, cuadro-firma, franja de acento y sello
%  del instituto. El arte (imagen real o placeholder) va DENTRO del marco.
%  Uso (1ª línea de cada sNN.tex):
%    \aperturaseccion{ruta/img.png}{Título}{Una o dos frases.}
%  Personaliza por documento (opcional):  \renewcommand{\seccionkicker}{...}
% =====================================================================
\newcounter{secportada}
\newcounter{unidadnum}                       % nº de unidad actual (lo fija \aperturacapitulo)
% Número visible de sección = Unidad.Sección (ej. 1.2). Si no hay unidad
% definida (unidadnum=0), cae al número simple de sección.
\newcommand{\numseccion}{\ifnum\value{unidadnum}>0 \theunidadnum.\thesecportada\else\thesecportada\fi}
\newcommand{\seccionkicker}{\hdrcapitulo}   % por defecto, lo del encabezado

% El escenario branded (contenedor constante del arte)
\newtcolorbox{secstage}{%
  enhanced, sharp corners, boxrule=0pt, frame hidden,
  colback=subjectcolor!6, halign=center, valign=center,
  height=5.0cm, left=2mm, right=2mm, top=10mm, bottom=8mm,
  overlay={%
    % cuadro-firma + kicker (arriba-izquierda)
    \fill[subjectcolor] ([shift={(4mm,-4mm)}]frame.north west) rectangle ++(2.6mm,-2.6mm);
    \node[anchor=north west, font=\InterDisplaySB\footnotesize, text=subjectcolor!85]
      at ([shift={(8.5mm,-3.4mm)}]frame.north west) {\seccionkicker};
    % número de serie grande y tenue (arriba-derecha) = marca
    \node[anchor=north east, font=\InterDisplayBlk, text=subjectcolor!18, inner sep=0pt]
      at ([shift={(-4mm,-1.5mm)}]frame.north east)
      {\fontsize{34}{34}\selectfont \numseccion};
    % franja de acento (abajo-izquierda)
    \fill[subjectcolor] ([shift={(4mm,4mm)}]frame.south west) rectangle ++(2.8cm,2.2mm);
    % sello del instituto (abajo-derecha)
    \node[anchor=south east, font=\sffamily\scriptsize, text=subjectcolor!55]
      at ([shift={(-4mm,3.4mm)}]frame.south east) {ITSTZ · Zapatoca};
  }}

\newcommand{\aperturaseccion}[3]{% {archivo}{título}{intro}
  \clearpage\refstepcounter{secportada}%
  \begin{secstage}
  \IfFileExists{assets/imgs/#1}{%
    \includegraphics[width=0.9\linewidth, height=3.2cm, keepaspectratio]{assets/imgs/#1}}{%
    \begin{tikzpicture}
      \node[draw=subjectcolor!35, dashed, line width=0.6pt, fill=white,
            minimum width=0.86\linewidth, minimum height=2.9cm,
            align=center, text=subjectcolor!50, font=\sffamily\small]
        {{\fontsize{22}{22}\selectfont\faImage}\\[5pt]
         \textbf{[ imagen pendiente ]}\\[1pt]
         \scriptsize\ttfamily assets/imgs/#1};
    \end{tikzpicture}}%
  \end{secstage}
  \vspace{0.5em}
  {\color{subjectcolor}\rule{3.2cm}{2.5pt}}\par\vspace{0.3em}
  \section{{\color{subjectcolor}\numseccion}\enspace #2}%
  {\InterLight\fontsize{12.5}{17}\selectfont\color{black!58} #3\par}
  \vspace{0.4em}}

% =====================================================================
%  DESTACADO — "lo esencial": triángulo lateral + degradado muy sutil
% =====================================================================
\newtcolorbox{destacado}[1]{
  enhanced, breakable, sharp corners, boxrule=0pt, frame hidden,
  interior style={left color=subjectcolor!16, right color=subjectcolor!2},
  borderline west={2.4pt}{0pt}{subjectcolor},
  left=7mm, right=5mm, top=3.5mm, bottom=3.5mm,
  before skip=\bloqueskip, after skip=\bloqueskip,
  overlay unbroken and first={%
    \fill[subjectcolor] ([yshift=-4.5mm]frame.north west)
        -- ++(5.5pt,2.75pt) -- ++(0,-5.5pt) -- cycle;},
  before upper={\blocklabel{subjectcolor}{Lo esencial}{#1}}
}

% --- Resaltado inline suave ---
\newtcbox{\resaltar}{on line, nobeforeafter, boxsep=0pt,
  left=2.5pt, right=2.5pt, top=1pt, bottom=1pt, arc=2.5pt,
  colback=subjectcolor!15, colframe=subjectcolor!15}

% =====================================================================
%  EXTENSIONES SANTILLANA (opt-in) — cajas temáticas para dar variedad
%  visual y romper la monotonía. NO alteran los entornos base; úsalas
%  donde aporten. Mantienen el color de asignatura (coherencia de marca).
% =====================================================================

% --- ¿SABÍAS QUE? — curiosidad breve · tarjeta redondeada con bombillo --
\newtcolorbox{sabias}{
  enhanced, breakable, arc=2.5mm, boxrule=0.8pt,
  colback=subjectcolor!6!white, colframe=subjectcolor!30,
  left=4mm, right=4mm, top=3mm, bottom=3mm,
  before skip=\bloqueskip, after skip=\bloqueskip,
  before upper={{\sffamily\bfseries\footnotesize\color{subjectcolor}%
    \faLightbulb\enspace\MakeUppercase{¿Sabías que?}}%
    \par\nopagebreak\vspace{0.45em}}
}

% --- CIENCIA Y SOCIEDAD — impacto/ética · banner de título -------------
\newtcolorbox{cienciasociedad}[1]{
  enhanced, breakable, sharp corners,
  colback=subjectcolor!5!white, colframe=subjectcolor, boxrule=0pt,
  colbacktitle=subjectcolor, coltitle=white,
  fonttitle=\sffamily\bfseries\footnotesize,
  title={\faGlobeAmericas\ \ \MakeUppercase{Ciencia y sociedad}%
         \ifblank{#1}{}{\ \ \textnormal{\sffamily\small· #1}}},
  toptitle=1.6mm, bottomtitle=1.6mm,
  left=4mm, right=4mm, top=3mm, bottom=3mm,
  before skip=\bloqueskip, after skip=\bloqueskip,
}

% --- MANOS A LA OBRA — práctica experimental · marco sólido con matraz --
\newtcolorbox{laboratorio}[1]{
  enhanced, breakable, sharp corners, arc=0pt,
  colback=white, colframe=subjectcolor!55, boxrule=1pt,
  left=4mm, right=4mm, top=3mm, bottom=3mm,
  before skip=\bloqueskip, after skip=\bloqueskip,
  before upper={{\sffamily\bfseries\footnotesize\color{subjectcolor}%
    \faFlask\enspace\MakeUppercase{Manos a la obra}}%
    \ifblank{#1}{}{\ {\color{black!35}\textbf{·}}\ %
      {\sffamily\bfseries\footnotesize\color{black!78}#1}}%
    \par\nopagebreak\vspace{0.45em}}
}
% Lista de materiales dentro de \begin{laboratorio}
\newcommand{\materiales}[1]{\par\noindent
  {\sffamily\bfseries\footnotesize\color{black!70}Materiales:}\ %
  {\footnotesize #1}\par\vspace{0.3em}}

% --- IMAGEN INLINE con texto alrededor (Santillana) -------------------
%   \figinline[r|l]{ancho}{archivo}{pie}   — colócala al inicio de un párrafo
\newcommand{\figinline}[4][r]{%
  \begin{wrapfigure}{#1}{#2}
  \centering\vspace{-0.5\baselineskip}
  \IfFileExists{assets/imgs/#3}{%
    \includegraphics[width=\linewidth]{assets/imgs/#3}}{%
    \fbox{\parbox[c][3cm][c]{0.92\linewidth}{\centering\sffamily\scriptsize
      \color{black!55}{\Large\faImage}\\[3pt][imagen pendiente]\\\ttfamily\tiny #3}}}%
  \par\vspace{2pt}{\footnotesize\sffamily\color{black!60}\textbf{Fig.}\enspace #4}%
  \vspace{-0.4\baselineskip}
  \end{wrapfigure}}

% --- APERTURA CON HERO LATERAL (imagen grande al lado del título) ------
%   \aperturahero{archivo}{título}{intro}  — alternativa a \aperturaseccion
\newcommand{\aperturahero}[3]{%
  \clearpage\refstepcounter{secportada}\stepcounter{section}%
  \addcontentsline{toc}{section}{\numseccion\enspace #2}%
  \noindent\begin{minipage}[c]{0.62\linewidth}
    {\InterDisplaySB\footnotesize\color{subjectcolor!85}\MakeUppercase{\seccionkicker}}\par
    \vspace{1.8mm}
    {\LARGE\InterDisplaySB\color{subjectcolor}%
      {\color{subjectcolor!35}\numseccion}\,\,#2}\par
    \vspace{2mm}{\color{subjectcolor!30}\rule{\linewidth}{0.8pt}}\par\vspace{2mm}
    {\InterLight\fontsize{12}{16}\selectfont\color{black!58}#3}
  \end{minipage}\hfill
  \begin{minipage}[c]{0.34\linewidth}\centering
    \IfFileExists{assets/imgs/#1}{%
      \includegraphics[width=\linewidth]{assets/imgs/#1}}{%
      \fbox{\parbox[c][3.4cm][c]{0.95\linewidth}{\centering\sffamily\scriptsize
        \color{black!55}{\Large\faImage}\\[3pt][imagen pendiente]\\\ttfamily\tiny #1}}}%
  \end{minipage}\par\vspace{1.3em}}

% --- MAPAS (conceptual y mental) — estilos TikZ con texto nítido -------
\tikzset{
  croot/.style={draw=subjectcolor, fill=subjectcolor, text=white, rounded corners=3pt,
    font=\sffamily\bfseries\footnotesize, align=center, inner sep=5pt, minimum height=9mm},
  cnode/.style={draw=subjectcolor!55, fill=subjectcolor!8, rounded corners=2.5pt,
    font=\sffamily\footnotesize, align=center, inner sep=4pt, text width=2.5cm},
  cleaf/.style={draw=subjectcolor!45, fill=white, rounded corners=2.5pt,
    font=\sffamily\footnotesize, align=center, inner sep=4pt},
  clink/.style={-{Stealth[length=5pt]}, subjectcolor!65, semithick},
  cline/.style={subjectcolor!55, semithick},
}
% Caja contenedora para un mapa (título + regla superior)
\newtcolorbox{mapabox}[2][Mapa conceptual]{
  enhanced, sharp corners, boxrule=0pt, frame hidden, colback=white,
  borderline north={1.4pt}{0pt}{subjectcolor},
  left=1mm, right=1mm, top=3mm, bottom=2mm,
  before skip=\bloqueskip, after skip=\bloqueskip,
  before upper={\blocklabel{subjectcolor}{#1}{#2}}
}

% --- Papel cuadriculado para graficar a mano (hojas de trabajo) --------
%   \papelcuadricula[filas]{ancho en cm}   (cada celda = 0.6 cm)
\newcommand{\papelcuadricula}[2][8]{%
  \par\smallskip\begin{center}
  \begin{tikzpicture}
    \draw[step=0.6cm, black!16, line width=0.3pt] (0,0) grid (#2,{#1*0.6});
    \draw[black!45, line width=0.7pt] (0,0) -- (#2,0);
    \draw[black!45, line width=0.7pt] (0,0) -- (0,{#1*0.6});
  \end{tikzpicture}
  \end{center}\par\smallskip}

% =====================================================================
%  APERTURA DE CAPÍTULO — título grande, número marca de agua, aire
% =====================================================================
\newcommand{\aperturacapitulo}[3]{% {número}{título}{subtítulo}
  \setcounter{unidadnum}{#1}\setcounter{secportada}{0}% fija unidad y reinicia secciones
  \thispagestyle{empty}%
  \begin{tikzpicture}[remember picture, overlay]
    \node[anchor=north east, text=subjectcolor!9,
          font=\InterDisplayThin]
      at ([xshift=-0.9cm, yshift=-2.2cm]current page.north east)
      {\fontsize{210}{210}\selectfont #1};
  \end{tikzpicture}%
  \begingroup\raggedright\null\vspace{3.8cm}\par
  {\InterDisplaySB\fontsize{12}{14}\selectfont\color{subjectcolor}%
    CAP\'ITULO\ #1}\par
  \vspace{0.18cm}{\color{black!16}\rule{4cm}{1pt}}\par
  \vspace{0.9cm}
  {\InterDisplaySB\fontsize{46}{50}\selectfont\color{black!88}#2}\par
  \vspace{0.5cm}{\color{subjectcolor}\rule{3.2cm}{2.5pt}}\par
  \vspace{0.55cm}
  {\InterLight\fontsize{15}{21}\selectfont\color{black!55}#3}\par
  \vfill
  {\sffamily\footnotesize\color{black!45}\hdrinstituto\quad
    \textcolor{black!22}{·}\quad Prof.\ Juan Diego A.\ Prada R.}\par
  \endgroup\clearpage}

% =====================================================================
%  DESBLOQUEA — pie discreto, no caja
% =====================================================================
\newcommand{\desbloquea}[1]{%
  \par\addvspace{\bloqueskip}%
  \noindent{\color{subjectcolor!25}\rule{\linewidth}{0.5pt}}\par\vspace{0.4em}%
  {\footnotesize\sffamily\color{black!60}%
    \textcolor{subjectcolor}{\faUnlock}\ \textbf{Dominar esto te desbloquea:}\ #1}\par}
; opcional tema de asignatura):
%     % =====================================================================
%  DESIGN SYSTEM v5 — EDITORIAL ITSTZ
%  Filosofía: texto primero · un solo color de asignatura · sin sombras
%  · sin marcos pesados · tratamientos planos (filete lateral / regla).
%
%  Vocabulario semántico (esto es TODO lo disponible):
%    \section / \subsection      jerarquía
%    definicion{título}          concepto/regla clave (filete color)
%    ejemplo{título}             ejemplo resuelto (numerado, sin caja)
%    procedimiento{título}       pasos (filete neutro) + \paso{n}{texto}
%    nota[título]                aparte único gris (sabías/dato/video/margen)
%    alerta{título}              advertencia ámbar (usar rara vez)
%    practica{título}            ejercicios/hoja/ICFES (regla superior)
%    \figura{archivo}{pie}       imagen real o placeholder automático
%    \desbloquea{lista}          pie discreto de prerrequisitos
%
%  REGLA DE ORO: las secciones (sNN.tex) NO llevan \vspace, color ni
%  minipage de maquetación. Solo contenido semántico. El diseño vive aquí.
% =====================================================================
\documentclass[11pt, letterpaper, twoside]{article}

\usepackage{fontspec}
\usepackage{polyglossia}
\setdefaultlanguage{spanish}
\usepackage[inner=2.1cm, outer=3.7cm, top=2.6cm, bottom=2.3cm,
            marginparwidth=2.55cm, marginparsep=0.45cm,
            headheight=16pt, headsep=18pt, footskip=24pt]{geometry}
\usepackage{microtype}
\usepackage{multicol}
\usepackage{xcolor}
\usepackage{graphicx}
\usepackage{tikz}
\usetikzlibrary{calc, arrows.meta, positioning}
\usepackage{wrapfig}
\usepackage[most]{tcolorbox}
\usepackage{amsmath, amssymb}
\usepackage{siunitx}
\usepackage{fontawesome5}
\usepackage{fancyhdr}
\usepackage{titlesec}
\usepackage{enumitem}
\usepackage{etoolbox}
\usepackage{titletoc}
\usepackage{array}
\usepackage{booktabs}
\usepackage{colortbl}
\usepackage{nicematrix}   % rejilla completa (hvlines) para tablas de hojas

% --- TIPOGRAFÍA -------------------------------------------------------
% Cuerpo: Charter (serif cálida y muy legible). La matemática (Latin
% Modern, serif) combina con ella. Títulos/etiquetas/UI: Inter Display.
\setmainfont{Charter}[
  UprightFont    = {Charter},
  BoldFont       = {Charter Bold},
  ItalicFont     = {Charter Italic},
  BoldItalicFont = {Charter Bold Italic},
  Ligatures = TeX, Numbers = Lining ]
\setsansfont{Inter}[
  UprightFont = Inter-Regular, BoldFont = Inter-SemiBold,
  ItalicFont  = Inter-Italic, Ligatures = TeX, Numbers = Lining ]
\newfontfamily\InterLight{Inter-Light}[ItalicFont=Inter-LightItalic, Ligatures=TeX]
\newfontfamily\InterDisplay{Inter Display}[Ligatures=TeX]
\newfontfamily\InterDisplaySB{Inter Display SemiBold}[Ligatures=TeX]
\newfontfamily\InterDisplayThin{Inter Display Thin}[Ligatures=TeX]
\newfontfamily\InterDisplayBlk{Inter Display Black}[Ligatures=TeX]

% --- RITMO --------------------------------------------------------------
\linespread{1.22}
\setlength{\parindent}{0pt}
\setlength{\parskip}{0.5em}
\newlength{\bloqueskip}\setlength{\bloqueskip}{1.0em}   % aire entre bloques
\setlength{\columnsep}{1.6em}                            % medianil multicols

% =====================================================================
%  COLOR — UNA sola variable de asignatura + neutros. Sin arcoíris.
% =====================================================================
\definecolor{mat} {RGB}{ 30,  64, 120}   % azul profundo
\definecolor{cnat}{RGB}{ 16, 110,  80}   % esmeralda
\definecolor{csoc}{RGB}{165,  75,  45}   % terracota
\definecolor{eco} {RGB}{ 70,  55, 130}   % índigo
\definecolor{alertcolor}{RGB}{180, 95, 10}  % ámbar — solo advertencias

\colorlet{subjectcolor}{mat}             % por defecto
\newcommand{\setsubject}[1]{\colorlet{subjectcolor}{#1}}

% =====================================================================
%  JERARQUÍA TIPOGRÁFICA
% =====================================================================
\setcounter{secnumdepth}{0}
\titleformat{\section}
  {\LARGE\InterDisplaySB\color{subjectcolor}}{}{0pt}{}
  [\vspace{2pt}{\color{subjectcolor!30}\titlerule[0.8pt]}]
\titlespacing*{\section}{0pt}{3.2ex plus 0.6ex}{1.8ex plus 0.3ex}

\titleformat{\subsection}
  {\large\InterDisplaySB\color{black!82}}{}{0pt}{}
\titlespacing*{\subsection}{0pt}{2.0ex plus 0.4ex}{0.5ex}

% --- TABLA DE CONTENIDO -------------------------------------------------
\titlecontents{section}[0em]
  {\addvspace{0.45pc}\sffamily\bfseries\small\color{black!85}}
  {}{}{\titlerule*[0.35pc]{.}\contentspage}
\titlecontents{subsection}[1.4em]
  {\sffamily\footnotesize\color{black!60}}
  {}{}{\titlerule*[0.35pc]{.}\contentspage}

% =====================================================================
%  ENCABEZADO / PIE — editorial limpio (regla fina, sin franjas)
% =====================================================================
\newcommand{\hdrinstituto}{Instituto Técnico Santo Tomás · Zapatoca}
\newcommand{\hdrcapitulo}{}            % se fija en el documento maestro
\pagestyle{fancy}\fancyhf{}
\renewcommand{\headrulewidth}{0pt}
\fancyhead[LE]{\footnotesize\sffamily\color{black!55}\hdrinstituto}
\fancyhead[RO]{\footnotesize\sffamily\color{subjectcolor}\hdrcapitulo}
\fancyhead[LO,RE]{}
\renewcommand{\headrule}{\vspace{2pt}{\color{black!18}\hrule height 0.4pt}}

% --- PIE: número exterior (cuadro de color) + crédito interior + filete ---
\renewcommand{\footrulewidth}{0pt}
\renewcommand{\footrule}{{\color{black!15}\hrule height 0.4pt}\vspace{3pt}}
\newcommand{\pagenumbox}{%
  \tikz[baseline=-3.2pt]{\node[fill=subjectcolor, text=white, inner xsep=7pt,
     inner ysep=3pt, font=\sffamily\bfseries\footnotesize]{\thepage};}}
\fancyfoot[RO,LE]{\pagenumbox}
\fancyfoot[LO,RE]{\footnotesize\sffamily\color{black!45}%
  \hdrcapitulo\ \textcolor{black!22}{·}\ Prof.\ Juan Diego A.\ Prada R.}

% =====================================================================
%  HELPERS
% =====================================================================
% Etiqueta de bloque:  TAG · título   (color del tag configurable)
\newcommand{\blocklabel}[3]{% #1=color  #2=TAG  #3=título
  {\sffamily\bfseries\footnotesize\color{#1}\MakeUppercase{#2}}%
  \ifblank{#3}{}{\,{\color{black!35}\textbf{·}}\,%
    {\sffamily\bfseries\footnotesize\color{black!78}#3}}%
  \par\nopagebreak\vspace{0.4em}}

% --- CAJAS AL MARGEN — personajes históricos e hitos -------------------
% Van en el margen exterior ancho (\marginpar respeta la geometría twoside
% y flota para no chocar). Son contenido semántico: colócalas en el cuerpo,
% cerca del párrafo que comentan, ENTRE párrafos. NO usarlas dentro de
% multicols, tcolorbox, tablas, ni en la misma línea que un \section.
%   \personaje{Nombre}{años · rol}{descripción breve}
%   \hito{año}{qué pasó (breve)}
\setlength{\marginparpush}{8pt}
\newcommand{\margenbox}[1]{\marginpar{%
  \begingroup\raggedright\footnotesize\sffamily
  \setlength{\parskip}{0.18em}\linespread{1.05}\selectfont
  #1\par\endgroup}}
\newcommand{\personaje}[3]{\margenbox{%
  {\color{subjectcolor}\rule{\linewidth}{1.1pt}}\par
  {\InterDisplaySB\scriptsize\color{subjectcolor}#1}\par
  {\fontsize{6.8pt}{8pt}\selectfont\itshape\color{black!50}#2}\par
  \vspace{0.15em}{\scriptsize\color{black!72}#3}}}
\newcommand{\hito}[2]{\margenbox{%
  {\sffamily\bfseries\scriptsize\colorbox{subjectcolor}{\textcolor{white}{\,#1\,}}}\par
  \vspace{0.15em}{\scriptsize\color{black!72}#2}}}

% =====================================================================
%  ENTORNOS — todos planos: filete lateral, sin marco, sin sombra
% =====================================================================

% --- DEFINICIÓN / concepto clave ---------------------------------------
\newtcolorbox{definicion}[1]{
  enhanced, breakable, sharp corners, boxrule=0pt, frame hidden,
  colback=subjectcolor!6!white,
  borderline west={2.2pt}{0pt}{subjectcolor},
  left=5mm, right=4mm, top=3mm, bottom=3mm,
  before skip=\bloqueskip, after skip=\bloqueskip,
  before upper={\blocklabel{subjectcolor}{Definición}{#1}}
}

% --- PROCEDIMIENTO / pasos ---------------------------------------------
\newtcolorbox{procedimiento}[1]{
  enhanced, breakable, sharp corners, boxrule=0pt, frame hidden,
  colback=black!3!white,
  borderline west={2.2pt}{0pt}{black!45},
  left=5mm, right=4mm, top=3mm, bottom=3mm,
  before skip=\bloqueskip, after skip=\bloqueskip,
  before upper={\blocklabel{black!60}{Procedimiento}{#1}}
}
\newcommand{\paso}[2]{%
  \par\vspace{0.45em}\noindent
  \tikz[baseline=-0.6ex]\node[circle, fill=subjectcolor, text=white,
     inner sep=0pt, minimum size=1.45em,
     font=\sffamily\bfseries\footnotesize]{#1};\hspace{0.5em}#2\par}

% Glosa "¿por qué?" bajo un paso — narración en lenguaje cercano
\newcommand{\porque}[1]{\par\vspace{0.1em}\noindent\hspace{2.0em}%
  {\sffamily\footnotesize\itshape\color{subjectcolor}¿por qué?\,}%
  {\sffamily\footnotesize\itshape\color{black!55}#1}\par}

% Operación glosada: cuenta a la izquierda + porqué al margen derecho
% (estilo "comentario al lado del código"). Úsala dentro de ejemplo.
% \opg{ecuación}{porqué}   — primer arg en modo display (mat, cap0)
% \opgt{etiqueta}{porqué}  — primer arg en texto sans bold (física, ciencias)
\newcommand{\opg}[2]{\par\vspace{0.25em}\noindent
  \makebox[0.46\linewidth][l]{$\displaystyle #1$}%
  {\sffamily\footnotesize\itshape\color{black!52}#2}\par}
\newcommand{\opgt}[2]{\par\vspace{0.25em}\noindent
  \makebox[0.46\linewidth][l]{\sffamily\bfseries\color{subjectcolor!80}#1}%
  {\sffamily\footnotesize\itshape\color{black!52}#2}\par}

% --- NOTA / aparte único (sabías · dato · video · margen) --------------
\newtcolorbox{nota}[1][]{
  enhanced, breakable, sharp corners, boxrule=0pt, frame hidden,
  colback=black!4!white,
  borderline west={2.2pt}{0pt}{black!35},
  left=5mm, right=4mm, top=3mm, bottom=3mm,
  before skip=\bloqueskip, after skip=\bloqueskip,
  before upper={\ifblank{#1}{}{\blocklabel{black!55}{Nota}{#1}}}
}

% --- ALERTA / advertencia (única en ámbar, uso escaso) -----------------
\newtcolorbox{alerta}[1]{
  enhanced, breakable, sharp corners, boxrule=0pt, frame hidden,
  colback=alertcolor!7!white,
  borderline west={2.2pt}{0pt}{alertcolor},
  left=5mm, right=4mm, top=3mm, bottom=3mm,
  before skip=\bloqueskip, after skip=\bloqueskip,
  before upper={\blocklabel{alertcolor}{Atención}{#1}}
}
% Sub-ítem de alerta (antes \caso)
\newcommand{\caso}[2]{\par\vspace{0.35em}\noindent
  {\sffamily\bfseries\color{alertcolor!85!black}#1.}\enspace #2\par}

% --- CONTRASTE (la virtud y su exceso → pregunta que cuestiona) --------
%   \contraste{Concepto}{La virtud / lo valioso}{Su exceso / su riesgo}{Pregunta}
\newtcolorbox{contrastebox}{
  enhanced, breakable, sharp corners, boxrule=0pt, frame hidden,
  colback=subjectcolor!5!white,
  borderline west={2.2pt}{0pt}{subjectcolor},
  left=5mm, right=4mm, top=3mm, bottom=3mm,
  before skip=\bloqueskip, after skip=\bloqueskip,
}
\newcommand{\contraste}[4]{%
\begin{contrastebox}
\blocklabel{subjectcolor}{Cuestiona esto}{\faBalanceScale\ #1}
\noindent\begin{minipage}[t]{0.45\linewidth}
  {\sffamily\footnotesize\bfseries\color{subjectcolor!85!black}La virtud}\par
  \vspace{0.2em}#2
\end{minipage}\hfill
{\color{subjectcolor!35}\vrule width 0.4pt}\hfill
\begin{minipage}[t]{0.45\linewidth}
  {\sffamily\footnotesize\bfseries\color{alertcolor!85!black}Su exceso}\par
  \vspace{0.2em}#3
\end{minipage}
\par\vspace{0.6em}
{\color{subjectcolor!25}\titlerule[0.6pt]}\par\vspace{0.4em}
\noindent{\sffamily\itshape\color{subjectcolor}\faArrowRight\ #4}
\end{contrastebox}}

% --- EJEMPLO (sin caja: regla fina + etiqueta numerada) ----------------
\newcounter{ejemplo}[section]
\newenvironment{ejemplo}[1]
  {\par\addvspace{\bloqueskip}\refstepcounter{ejemplo}%
   \noindent{\color{subjectcolor!28}\rule{\linewidth}{0.5pt}}\par\vspace{0.4em}%
   \noindent{\sffamily\bfseries\color{subjectcolor}Ejemplo \theejemplo.}%
   \enspace{\sffamily\bfseries\color{black!78}#1}\par\vspace{0.3em}}
  {\par\addvspace{\bloqueskip}}

% --- PRÁCTICA (ejercicios / hoja / ICFES) ------------------------------
\newtcolorbox{practica}[1]{
  enhanced, breakable, sharp corners, boxrule=0pt, frame hidden,
  colback=white,
  borderline north={1.4pt}{0pt}{subjectcolor},
  left=0mm, right=0mm, top=3mm, bottom=2mm,
  before skip=\bloqueskip, after skip=\bloqueskip,
  before upper={\blocklabel{subjectcolor}{Práctica}{#1}\raggedright}
}

% Ejercicio numerado con nivel opcional [básico|intermedio|reto]
\newcounter{ejer}[section]
\newcommand{\ejer}[2][]{%
  \stepcounter{ejer}\par\vspace{0.5em}\noindent
  {\sffamily\bfseries\color{subjectcolor}\theejer.}\enspace #2%
  \ifblank{#1}{}{\ \textcolor{black!45}{\footnotesize[\textit{#1}]}}\par}

% Niveles = un solo color a tres intensidades (NO nuevos colores)
\newcommand{\nivel}[2]{\par\smallskip\noindent
  \textcolor{subjectcolor!#1}{\small\faCircle}\enspace
  {\sffamily\bfseries\footnotesize\color{black!70}#2}\par}
\newcommand{\basico}{\nivel{40}{Básico}}
\newcommand{\intermedio}{\nivel{70}{Intermedio}}
\newcommand{\reto}{\nivel{100}{Reto}}

% Selección múltiple tipo ICFES — opciones con burbuja (estilo Saber)
\newcounter{pregicfes}[section]
\newcommand{\bubbleletter}[1]{\tikz[baseline=(bc.base)]\node[circle,
  draw=subjectcolor!55, line width=0.6pt, inner sep=1.5pt,
  font=\sffamily\bfseries\scriptsize, text=subjectcolor!90!black](bc){#1};}
\newcommand{\icfes}[5]{%
  \stepcounter{pregicfes}\par\vspace{0.6em}\noindent
  {\sffamily\bfseries\color{subjectcolor}\thepregicfes.}\enspace #1\par
  \begin{enumerate}[label={\protect\bubbleletter{\Alph*}}, leftmargin=1.5cm,
                    itemsep=3pt, topsep=3pt]
    \item #2 \item #3 \item #4 \item #5
  \end{enumerate}}
% Pregunta abierta de argumentación
\newcommand{\pregabierta}[1]{\par\vspace{0.6em}\noindent
  {\sffamily\bfseries\color{subjectcolor}\faComment\ Argumenta:}\enspace
  \textit{#1}\par\lineaescritura\lineaescritura\lineaescritura}

% --- LÍNEAS DE RESPUESTA -----------------------------------------------
\newcommand{\linea}[1][7cm]{\underline{\hspace{#1}}}
\newcommand{\lineaescritura}{\par\vspace{0.9em}\noindent
  {\color{black!30}\rule{\linewidth}{0.3pt}}}

% =====================================================================
%  EVALUACIÓN TIPO ICFES — cierre de unidad (banner + 2 columnas)
%  Uso:  \begin{evaluacion} \begin{multicols}{2} \icfes... \end{multicols}
%        \pregabierta{...} \end{evaluacion}
% =====================================================================
\newtcolorbox{evaluacion}[1][cierre de sección]{
  enhanced, breakable, sharp corners,
  colback=white, colframe=subjectcolor, boxrule=1pt,
  colbacktitle=subjectcolor, coltitle=white,
  fonttitle=\InterDisplaySB\normalsize,
  title={\faClipboardCheck\ \ EVALUACIÓN TIPO ICFES\ \ %
         \textnormal{\sffamily\small· #1}},
  toptitle=2mm, bottomtitle=2mm, left=4mm, right=4mm, top=3.5mm, bottom=3mm,
  before={\clearpage},
  before upper={\small\itshape\color{black!55}%
    Selección múltiple. Marca una sola opción por pregunta.\par\medskip\raggedright}
}

% =====================================================================
%  HOJA DE TRABAJO — recortable: borde punteado, escala de grises,
%  en página aparte (clearpage antes; el contenido que sigue continúa).
% =====================================================================
\newtcolorbox{hoja}[1]{
  enhanced, breakable, sharp corners,
  colback=black!2!white, boxrule=0pt, frame hidden,
  borderline={0.8pt}{0pt}{black!50, dash pattern=on 3pt off 2.5pt},
  before={\clearpage}, after={\par},
  left=6mm, right=6mm, top=6mm, bottom=6mm,
  before upper={%
    {\sffamily\bfseries\footnotesize\color{black!60}\faCut\ \ HOJA DE TRABAJO}%
    {\color{black!35}\ \textbf{·}\ }{\sffamily\bfseries\footnotesize\color{black!80}#1}%
    \hfill{\sffamily\footnotesize\color{black!55}%
      Nombre:\ \underline{\hspace{4.2cm}}\quad Fecha:\ \underline{\hspace{2cm}}}\par
    \vspace{2mm}{\color{black!30}\hrule height 0.4pt}\vspace{3mm}\raggedright}
}

% =====================================================================
%  TABLAHOJA — tabla con TODAS las líneas (rejilla) para hojas de trabajo.
%  Celdas altas para llenar a mano. Centrada. El encabezado se marca con
%  \rowcolor en la primera fila.
%  Uso:  \begin{tablahoja}{ccc} enc & enc & enc \\ a & b & c \\ \end{tablahoja}
% =====================================================================
\NewDocumentEnvironment{tablahoja}{m}{%
  \par\smallskip\centering\small\setlength{\tabcolsep}{4pt}%
  \begin{NiceTabular}{#1}[hvlines, cell-space-top-limit=5pt,
     cell-space-bottom-limit=9pt, colortbl-like]%
}{%
  \end{NiceTabular}\par\smallskip
}

% =====================================================================
%  FIGURAS — imagen real o placeholder automático (\IfFileExists)
% =====================================================================
\newcounter{figura}[section]
\newcommand{\figura}[3][0.7\linewidth]{% [ancho]{archivo}{pie}
  \par\addvspace{\bloqueskip}\refstepcounter{figura}%
  \begin{center}%
  \IfFileExists{assets/imgs/#2}{%
    \includegraphics[width=#1]{assets/imgs/#2}}{%
    \begin{tikzpicture}
      \node[draw=black!30, dashed, line width=0.6pt, fill=black!3,
            text width=0.66\linewidth, minimum height=2.6cm,
            align=center, inner sep=8pt,
            text=black!55, font=\sffamily\footnotesize]
        {{\Large\faImage}\\[4pt]\textbf{[ imagen pendiente ]}\\#3};
    \end{tikzpicture}}%
  \par\vspace{0.4em}%
  {\footnotesize\sffamily\color{black!60}%
    \textbf{Figura \thefigura.}\enspace #3}%
  \end{center}\par\addvspace{\bloqueskip}}

% =====================================================================
%  APERTURA DE SECCIÓN — marco branded de serie + título + "qué verás"
%  El "marco" es CONSTANTE en todas las secciones (= identidad de marca):
%  número grande de serie, kicker, cuadro-firma, franja de acento y sello
%  del instituto. El arte (imagen real o placeholder) va DENTRO del marco.
%  Uso (1ª línea de cada sNN.tex):
%    \aperturaseccion{ruta/img.png}{Título}{Una o dos frases.}
%  Personaliza por documento (opcional):  \renewcommand{\seccionkicker}{...}
% =====================================================================
\newcounter{secportada}
\newcounter{unidadnum}                       % nº de unidad actual (lo fija \aperturacapitulo)
% Número visible de sección = Unidad.Sección (ej. 1.2). Si no hay unidad
% definida (unidadnum=0), cae al número simple de sección.
\newcommand{\numseccion}{\ifnum\value{unidadnum}>0 \theunidadnum.\thesecportada\else\thesecportada\fi}
\newcommand{\seccionkicker}{\hdrcapitulo}   % por defecto, lo del encabezado

% El escenario branded (contenedor constante del arte)
\newtcolorbox{secstage}{%
  enhanced, sharp corners, boxrule=0pt, frame hidden,
  colback=subjectcolor!6, halign=center, valign=center,
  height=5.0cm, left=2mm, right=2mm, top=10mm, bottom=8mm,
  overlay={%
    % cuadro-firma + kicker (arriba-izquierda)
    \fill[subjectcolor] ([shift={(4mm,-4mm)}]frame.north west) rectangle ++(2.6mm,-2.6mm);
    \node[anchor=north west, font=\InterDisplaySB\footnotesize, text=subjectcolor!85]
      at ([shift={(8.5mm,-3.4mm)}]frame.north west) {\seccionkicker};
    % número de serie grande y tenue (arriba-derecha) = marca
    \node[anchor=north east, font=\InterDisplayBlk, text=subjectcolor!18, inner sep=0pt]
      at ([shift={(-4mm,-1.5mm)}]frame.north east)
      {\fontsize{34}{34}\selectfont \numseccion};
    % franja de acento (abajo-izquierda)
    \fill[subjectcolor] ([shift={(4mm,4mm)}]frame.south west) rectangle ++(2.8cm,2.2mm);
    % sello del instituto (abajo-derecha)
    \node[anchor=south east, font=\sffamily\scriptsize, text=subjectcolor!55]
      at ([shift={(-4mm,3.4mm)}]frame.south east) {ITSTZ · Zapatoca};
  }}

\newcommand{\aperturaseccion}[3]{% {archivo}{título}{intro}
  \clearpage\refstepcounter{secportada}%
  \begin{secstage}
  \IfFileExists{assets/imgs/#1}{%
    \includegraphics[width=0.9\linewidth, height=3.2cm, keepaspectratio]{assets/imgs/#1}}{%
    \begin{tikzpicture}
      \node[draw=subjectcolor!35, dashed, line width=0.6pt, fill=white,
            minimum width=0.86\linewidth, minimum height=2.9cm,
            align=center, text=subjectcolor!50, font=\sffamily\small]
        {{\fontsize{22}{22}\selectfont\faImage}\\[5pt]
         \textbf{[ imagen pendiente ]}\\[1pt]
         \scriptsize\ttfamily assets/imgs/#1};
    \end{tikzpicture}}%
  \end{secstage}
  \vspace{0.5em}
  {\color{subjectcolor}\rule{3.2cm}{2.5pt}}\par\vspace{0.3em}
  \section{{\color{subjectcolor}\numseccion}\enspace #2}%
  {\InterLight\fontsize{12.5}{17}\selectfont\color{black!58} #3\par}
  \vspace{0.4em}}

% =====================================================================
%  DESTACADO — "lo esencial": triángulo lateral + degradado muy sutil
% =====================================================================
\newtcolorbox{destacado}[1]{
  enhanced, breakable, sharp corners, boxrule=0pt, frame hidden,
  interior style={left color=subjectcolor!16, right color=subjectcolor!2},
  borderline west={2.4pt}{0pt}{subjectcolor},
  left=7mm, right=5mm, top=3.5mm, bottom=3.5mm,
  before skip=\bloqueskip, after skip=\bloqueskip,
  overlay unbroken and first={%
    \fill[subjectcolor] ([yshift=-4.5mm]frame.north west)
        -- ++(5.5pt,2.75pt) -- ++(0,-5.5pt) -- cycle;},
  before upper={\blocklabel{subjectcolor}{Lo esencial}{#1}}
}

% --- Resaltado inline suave ---
\newtcbox{\resaltar}{on line, nobeforeafter, boxsep=0pt,
  left=2.5pt, right=2.5pt, top=1pt, bottom=1pt, arc=2.5pt,
  colback=subjectcolor!15, colframe=subjectcolor!15}

% =====================================================================
%  EXTENSIONES SANTILLANA (opt-in) — cajas temáticas para dar variedad
%  visual y romper la monotonía. NO alteran los entornos base; úsalas
%  donde aporten. Mantienen el color de asignatura (coherencia de marca).
% =====================================================================

% --- ¿SABÍAS QUE? — curiosidad breve · tarjeta redondeada con bombillo --
\newtcolorbox{sabias}{
  enhanced, breakable, arc=2.5mm, boxrule=0.8pt,
  colback=subjectcolor!6!white, colframe=subjectcolor!30,
  left=4mm, right=4mm, top=3mm, bottom=3mm,
  before skip=\bloqueskip, after skip=\bloqueskip,
  before upper={{\sffamily\bfseries\footnotesize\color{subjectcolor}%
    \faLightbulb\enspace\MakeUppercase{¿Sabías que?}}%
    \par\nopagebreak\vspace{0.45em}}
}

% --- CIENCIA Y SOCIEDAD — impacto/ética · banner de título -------------
\newtcolorbox{cienciasociedad}[1]{
  enhanced, breakable, sharp corners,
  colback=subjectcolor!5!white, colframe=subjectcolor, boxrule=0pt,
  colbacktitle=subjectcolor, coltitle=white,
  fonttitle=\sffamily\bfseries\footnotesize,
  title={\faGlobeAmericas\ \ \MakeUppercase{Ciencia y sociedad}%
         \ifblank{#1}{}{\ \ \textnormal{\sffamily\small· #1}}},
  toptitle=1.6mm, bottomtitle=1.6mm,
  left=4mm, right=4mm, top=3mm, bottom=3mm,
  before skip=\bloqueskip, after skip=\bloqueskip,
}

% --- MANOS A LA OBRA — práctica experimental · marco sólido con matraz --
\newtcolorbox{laboratorio}[1]{
  enhanced, breakable, sharp corners, arc=0pt,
  colback=white, colframe=subjectcolor!55, boxrule=1pt,
  left=4mm, right=4mm, top=3mm, bottom=3mm,
  before skip=\bloqueskip, after skip=\bloqueskip,
  before upper={{\sffamily\bfseries\footnotesize\color{subjectcolor}%
    \faFlask\enspace\MakeUppercase{Manos a la obra}}%
    \ifblank{#1}{}{\ {\color{black!35}\textbf{·}}\ %
      {\sffamily\bfseries\footnotesize\color{black!78}#1}}%
    \par\nopagebreak\vspace{0.45em}}
}
% Lista de materiales dentro de \begin{laboratorio}
\newcommand{\materiales}[1]{\par\noindent
  {\sffamily\bfseries\footnotesize\color{black!70}Materiales:}\ %
  {\footnotesize #1}\par\vspace{0.3em}}

% --- IMAGEN INLINE con texto alrededor (Santillana) -------------------
%   \figinline[r|l]{ancho}{archivo}{pie}   — colócala al inicio de un párrafo
\newcommand{\figinline}[4][r]{%
  \begin{wrapfigure}{#1}{#2}
  \centering\vspace{-0.5\baselineskip}
  \IfFileExists{assets/imgs/#3}{%
    \includegraphics[width=\linewidth]{assets/imgs/#3}}{%
    \fbox{\parbox[c][3cm][c]{0.92\linewidth}{\centering\sffamily\scriptsize
      \color{black!55}{\Large\faImage}\\[3pt][imagen pendiente]\\\ttfamily\tiny #3}}}%
  \par\vspace{2pt}{\footnotesize\sffamily\color{black!60}\textbf{Fig.}\enspace #4}%
  \vspace{-0.4\baselineskip}
  \end{wrapfigure}}

% --- APERTURA CON HERO LATERAL (imagen grande al lado del título) ------
%   \aperturahero{archivo}{título}{intro}  — alternativa a \aperturaseccion
\newcommand{\aperturahero}[3]{%
  \clearpage\refstepcounter{secportada}\stepcounter{section}%
  \addcontentsline{toc}{section}{\numseccion\enspace #2}%
  \noindent\begin{minipage}[c]{0.62\linewidth}
    {\InterDisplaySB\footnotesize\color{subjectcolor!85}\MakeUppercase{\seccionkicker}}\par
    \vspace{1.8mm}
    {\LARGE\InterDisplaySB\color{subjectcolor}%
      {\color{subjectcolor!35}\numseccion}\,\,#2}\par
    \vspace{2mm}{\color{subjectcolor!30}\rule{\linewidth}{0.8pt}}\par\vspace{2mm}
    {\InterLight\fontsize{12}{16}\selectfont\color{black!58}#3}
  \end{minipage}\hfill
  \begin{minipage}[c]{0.34\linewidth}\centering
    \IfFileExists{assets/imgs/#1}{%
      \includegraphics[width=\linewidth]{assets/imgs/#1}}{%
      \fbox{\parbox[c][3.4cm][c]{0.95\linewidth}{\centering\sffamily\scriptsize
        \color{black!55}{\Large\faImage}\\[3pt][imagen pendiente]\\\ttfamily\tiny #1}}}%
  \end{minipage}\par\vspace{1.3em}}

% --- MAPAS (conceptual y mental) — estilos TikZ con texto nítido -------
\tikzset{
  croot/.style={draw=subjectcolor, fill=subjectcolor, text=white, rounded corners=3pt,
    font=\sffamily\bfseries\footnotesize, align=center, inner sep=5pt, minimum height=9mm},
  cnode/.style={draw=subjectcolor!55, fill=subjectcolor!8, rounded corners=2.5pt,
    font=\sffamily\footnotesize, align=center, inner sep=4pt, text width=2.5cm},
  cleaf/.style={draw=subjectcolor!45, fill=white, rounded corners=2.5pt,
    font=\sffamily\footnotesize, align=center, inner sep=4pt},
  clink/.style={-{Stealth[length=5pt]}, subjectcolor!65, semithick},
  cline/.style={subjectcolor!55, semithick},
}
% Caja contenedora para un mapa (título + regla superior)
\newtcolorbox{mapabox}[2][Mapa conceptual]{
  enhanced, sharp corners, boxrule=0pt, frame hidden, colback=white,
  borderline north={1.4pt}{0pt}{subjectcolor},
  left=1mm, right=1mm, top=3mm, bottom=2mm,
  before skip=\bloqueskip, after skip=\bloqueskip,
  before upper={\blocklabel{subjectcolor}{#1}{#2}}
}

% --- Papel cuadriculado para graficar a mano (hojas de trabajo) --------
%   \papelcuadricula[filas]{ancho en cm}   (cada celda = 0.6 cm)
\newcommand{\papelcuadricula}[2][8]{%
  \par\smallskip\begin{center}
  \begin{tikzpicture}
    \draw[step=0.6cm, black!16, line width=0.3pt] (0,0) grid (#2,{#1*0.6});
    \draw[black!45, line width=0.7pt] (0,0) -- (#2,0);
    \draw[black!45, line width=0.7pt] (0,0) -- (0,{#1*0.6});
  \end{tikzpicture}
  \end{center}\par\smallskip}

% =====================================================================
%  APERTURA DE CAPÍTULO — título grande, número marca de agua, aire
% =====================================================================
\newcommand{\aperturacapitulo}[3]{% {número}{título}{subtítulo}
  \setcounter{unidadnum}{#1}\setcounter{secportada}{0}% fija unidad y reinicia secciones
  \thispagestyle{empty}%
  \begin{tikzpicture}[remember picture, overlay]
    \node[anchor=north east, text=subjectcolor!9,
          font=\InterDisplayThin]
      at ([xshift=-0.9cm, yshift=-2.2cm]current page.north east)
      {\fontsize{210}{210}\selectfont #1};
  \end{tikzpicture}%
  \begingroup\raggedright\null\vspace{3.8cm}\par
  {\InterDisplaySB\fontsize{12}{14}\selectfont\color{subjectcolor}%
    CAP\'ITULO\ #1}\par
  \vspace{0.18cm}{\color{black!16}\rule{4cm}{1pt}}\par
  \vspace{0.9cm}
  {\InterDisplaySB\fontsize{46}{50}\selectfont\color{black!88}#2}\par
  \vspace{0.5cm}{\color{subjectcolor}\rule{3.2cm}{2.5pt}}\par
  \vspace{0.55cm}
  {\InterLight\fontsize{15}{21}\selectfont\color{black!55}#3}\par
  \vfill
  {\sffamily\footnotesize\color{black!45}\hdrinstituto\quad
    \textcolor{black!22}{·}\quad Prof.\ Juan Diego A.\ Prada R.}\par
  \endgroup\clearpage}

% =====================================================================
%  DESBLOQUEA — pie discreto, no caja
% =====================================================================
\newcommand{\desbloquea}[1]{%
  \par\addvspace{\bloqueskip}%
  \noindent{\color{subjectcolor!25}\rule{\linewidth}{0.5pt}}\par\vspace{0.4em}%
  {\footnotesize\sffamily\color{black!60}%
    \textcolor{subjectcolor}{\faUnlock}\ \textbf{Dominar esto te desbloquea:}\ #1}\par}

%     % =====================================================================
%  TEMA · CIENCIAS NATURALES — variante del design system (GUÍAS)
%  Identidad visual propia, distinta de Matemáticas (azul/pizarra):
%    · color base ESMERALDA (\setsubject{cnat})
%    · acento LIMA orgánico solo para ornamentos (no rompe la regla de
%      "un solo color de asignatura")
%    · iconografía de ciencias (hoja · ADN · matraz · átomo)
%    · reutiliza las cajas de ciencia que ya trae el preamble:
%      \begin{laboratorio}, \begin{sabias}, \begin{cienciasociedad}
%
%  Sirve a futuro para BIOLOGÍA, QUÍMICA y FÍSICA EXPERIMENTAL.
%
%  INVOCACIÓN (una sola línea, tras % =====================================================================
%  DESIGN SYSTEM v5 — EDITORIAL ITSTZ
%  Filosofía: texto primero · un solo color de asignatura · sin sombras
%  · sin marcos pesados · tratamientos planos (filete lateral / regla).
%
%  Vocabulario semántico (esto es TODO lo disponible):
%    \section / \subsection      jerarquía
%    definicion{título}          concepto/regla clave (filete color)
%    ejemplo{título}             ejemplo resuelto (numerado, sin caja)
%    procedimiento{título}       pasos (filete neutro) + \paso{n}{texto}
%    nota[título]                aparte único gris (sabías/dato/video/margen)
%    alerta{título}              advertencia ámbar (usar rara vez)
%    practica{título}            ejercicios/hoja/ICFES (regla superior)
%    \figura{archivo}{pie}       imagen real o placeholder automático
%    \desbloquea{lista}          pie discreto de prerrequisitos
%
%  REGLA DE ORO: las secciones (sNN.tex) NO llevan \vspace, color ni
%  minipage de maquetación. Solo contenido semántico. El diseño vive aquí.
% =====================================================================
\documentclass[11pt, letterpaper, twoside]{article}

\usepackage{fontspec}
\usepackage{polyglossia}
\setdefaultlanguage{spanish}
\usepackage[inner=2.1cm, outer=3.7cm, top=2.6cm, bottom=2.3cm,
            marginparwidth=2.55cm, marginparsep=0.45cm,
            headheight=16pt, headsep=18pt, footskip=24pt]{geometry}
\usepackage{microtype}
\usepackage{multicol}
\usepackage{xcolor}
\usepackage{graphicx}
\usepackage{tikz}
\usetikzlibrary{calc, arrows.meta, positioning}
\usepackage{wrapfig}
\usepackage[most]{tcolorbox}
\usepackage{amsmath, amssymb}
\usepackage{siunitx}
\usepackage{fontawesome5}
\usepackage{fancyhdr}
\usepackage{titlesec}
\usepackage{enumitem}
\usepackage{etoolbox}
\usepackage{titletoc}
\usepackage{array}
\usepackage{booktabs}
\usepackage{colortbl}
\usepackage{nicematrix}   % rejilla completa (hvlines) para tablas de hojas

% --- TIPOGRAFÍA -------------------------------------------------------
% Cuerpo: Charter (serif cálida y muy legible). La matemática (Latin
% Modern, serif) combina con ella. Títulos/etiquetas/UI: Inter Display.
\setmainfont{Charter}[
  UprightFont    = {Charter},
  BoldFont       = {Charter Bold},
  ItalicFont     = {Charter Italic},
  BoldItalicFont = {Charter Bold Italic},
  Ligatures = TeX, Numbers = Lining ]
\setsansfont{Inter}[
  UprightFont = Inter-Regular, BoldFont = Inter-SemiBold,
  ItalicFont  = Inter-Italic, Ligatures = TeX, Numbers = Lining ]
\newfontfamily\InterLight{Inter-Light}[ItalicFont=Inter-LightItalic, Ligatures=TeX]
\newfontfamily\InterDisplay{Inter Display}[Ligatures=TeX]
\newfontfamily\InterDisplaySB{Inter Display SemiBold}[Ligatures=TeX]
\newfontfamily\InterDisplayThin{Inter Display Thin}[Ligatures=TeX]
\newfontfamily\InterDisplayBlk{Inter Display Black}[Ligatures=TeX]

% --- RITMO --------------------------------------------------------------
\linespread{1.22}
\setlength{\parindent}{0pt}
\setlength{\parskip}{0.5em}
\newlength{\bloqueskip}\setlength{\bloqueskip}{1.0em}   % aire entre bloques
\setlength{\columnsep}{1.6em}                            % medianil multicols

% =====================================================================
%  COLOR — UNA sola variable de asignatura + neutros. Sin arcoíris.
% =====================================================================
\definecolor{mat} {RGB}{ 30,  64, 120}   % azul profundo
\definecolor{cnat}{RGB}{ 16, 110,  80}   % esmeralda
\definecolor{csoc}{RGB}{165,  75,  45}   % terracota
\definecolor{eco} {RGB}{ 70,  55, 130}   % índigo
\definecolor{alertcolor}{RGB}{180, 95, 10}  % ámbar — solo advertencias

\colorlet{subjectcolor}{mat}             % por defecto
\newcommand{\setsubject}[1]{\colorlet{subjectcolor}{#1}}

% =====================================================================
%  JERARQUÍA TIPOGRÁFICA
% =====================================================================
\setcounter{secnumdepth}{0}
\titleformat{\section}
  {\LARGE\InterDisplaySB\color{subjectcolor}}{}{0pt}{}
  [\vspace{2pt}{\color{subjectcolor!30}\titlerule[0.8pt]}]
\titlespacing*{\section}{0pt}{3.2ex plus 0.6ex}{1.8ex plus 0.3ex}

\titleformat{\subsection}
  {\large\InterDisplaySB\color{black!82}}{}{0pt}{}
\titlespacing*{\subsection}{0pt}{2.0ex plus 0.4ex}{0.5ex}

% --- TABLA DE CONTENIDO -------------------------------------------------
\titlecontents{section}[0em]
  {\addvspace{0.45pc}\sffamily\bfseries\small\color{black!85}}
  {}{}{\titlerule*[0.35pc]{.}\contentspage}
\titlecontents{subsection}[1.4em]
  {\sffamily\footnotesize\color{black!60}}
  {}{}{\titlerule*[0.35pc]{.}\contentspage}

% =====================================================================
%  ENCABEZADO / PIE — editorial limpio (regla fina, sin franjas)
% =====================================================================
\newcommand{\hdrinstituto}{Instituto Técnico Santo Tomás · Zapatoca}
\newcommand{\hdrcapitulo}{}            % se fija en el documento maestro
\pagestyle{fancy}\fancyhf{}
\renewcommand{\headrulewidth}{0pt}
\fancyhead[LE]{\footnotesize\sffamily\color{black!55}\hdrinstituto}
\fancyhead[RO]{\footnotesize\sffamily\color{subjectcolor}\hdrcapitulo}
\fancyhead[LO,RE]{}
\renewcommand{\headrule}{\vspace{2pt}{\color{black!18}\hrule height 0.4pt}}

% --- PIE: número exterior (cuadro de color) + crédito interior + filete ---
\renewcommand{\footrulewidth}{0pt}
\renewcommand{\footrule}{{\color{black!15}\hrule height 0.4pt}\vspace{3pt}}
\newcommand{\pagenumbox}{%
  \tikz[baseline=-3.2pt]{\node[fill=subjectcolor, text=white, inner xsep=7pt,
     inner ysep=3pt, font=\sffamily\bfseries\footnotesize]{\thepage};}}
\fancyfoot[RO,LE]{\pagenumbox}
\fancyfoot[LO,RE]{\footnotesize\sffamily\color{black!45}%
  \hdrcapitulo\ \textcolor{black!22}{·}\ Prof.\ Juan Diego A.\ Prada R.}

% =====================================================================
%  HELPERS
% =====================================================================
% Etiqueta de bloque:  TAG · título   (color del tag configurable)
\newcommand{\blocklabel}[3]{% #1=color  #2=TAG  #3=título
  {\sffamily\bfseries\footnotesize\color{#1}\MakeUppercase{#2}}%
  \ifblank{#3}{}{\,{\color{black!35}\textbf{·}}\,%
    {\sffamily\bfseries\footnotesize\color{black!78}#3}}%
  \par\nopagebreak\vspace{0.4em}}

% --- CAJAS AL MARGEN — personajes históricos e hitos -------------------
% Van en el margen exterior ancho (\marginpar respeta la geometría twoside
% y flota para no chocar). Son contenido semántico: colócalas en el cuerpo,
% cerca del párrafo que comentan, ENTRE párrafos. NO usarlas dentro de
% multicols, tcolorbox, tablas, ni en la misma línea que un \section.
%   \personaje{Nombre}{años · rol}{descripción breve}
%   \hito{año}{qué pasó (breve)}
\setlength{\marginparpush}{8pt}
\newcommand{\margenbox}[1]{\marginpar{%
  \begingroup\raggedright\footnotesize\sffamily
  \setlength{\parskip}{0.18em}\linespread{1.05}\selectfont
  #1\par\endgroup}}
\newcommand{\personaje}[3]{\margenbox{%
  {\color{subjectcolor}\rule{\linewidth}{1.1pt}}\par
  {\InterDisplaySB\scriptsize\color{subjectcolor}#1}\par
  {\fontsize{6.8pt}{8pt}\selectfont\itshape\color{black!50}#2}\par
  \vspace{0.15em}{\scriptsize\color{black!72}#3}}}
\newcommand{\hito}[2]{\margenbox{%
  {\sffamily\bfseries\scriptsize\colorbox{subjectcolor}{\textcolor{white}{\,#1\,}}}\par
  \vspace{0.15em}{\scriptsize\color{black!72}#2}}}

% =====================================================================
%  ENTORNOS — todos planos: filete lateral, sin marco, sin sombra
% =====================================================================

% --- DEFINICIÓN / concepto clave ---------------------------------------
\newtcolorbox{definicion}[1]{
  enhanced, breakable, sharp corners, boxrule=0pt, frame hidden,
  colback=subjectcolor!6!white,
  borderline west={2.2pt}{0pt}{subjectcolor},
  left=5mm, right=4mm, top=3mm, bottom=3mm,
  before skip=\bloqueskip, after skip=\bloqueskip,
  before upper={\blocklabel{subjectcolor}{Definición}{#1}}
}

% --- PROCEDIMIENTO / pasos ---------------------------------------------
\newtcolorbox{procedimiento}[1]{
  enhanced, breakable, sharp corners, boxrule=0pt, frame hidden,
  colback=black!3!white,
  borderline west={2.2pt}{0pt}{black!45},
  left=5mm, right=4mm, top=3mm, bottom=3mm,
  before skip=\bloqueskip, after skip=\bloqueskip,
  before upper={\blocklabel{black!60}{Procedimiento}{#1}}
}
\newcommand{\paso}[2]{%
  \par\vspace{0.45em}\noindent
  \tikz[baseline=-0.6ex]\node[circle, fill=subjectcolor, text=white,
     inner sep=0pt, minimum size=1.45em,
     font=\sffamily\bfseries\footnotesize]{#1};\hspace{0.5em}#2\par}

% Glosa "¿por qué?" bajo un paso — narración en lenguaje cercano
\newcommand{\porque}[1]{\par\vspace{0.1em}\noindent\hspace{2.0em}%
  {\sffamily\footnotesize\itshape\color{subjectcolor}¿por qué?\,}%
  {\sffamily\footnotesize\itshape\color{black!55}#1}\par}

% Operación glosada: cuenta a la izquierda + porqué al margen derecho
% (estilo "comentario al lado del código"). Úsala dentro de ejemplo.
% \opg{ecuación}{porqué}   — primer arg en modo display (mat, cap0)
% \opgt{etiqueta}{porqué}  — primer arg en texto sans bold (física, ciencias)
\newcommand{\opg}[2]{\par\vspace{0.25em}\noindent
  \makebox[0.46\linewidth][l]{$\displaystyle #1$}%
  {\sffamily\footnotesize\itshape\color{black!52}#2}\par}
\newcommand{\opgt}[2]{\par\vspace{0.25em}\noindent
  \makebox[0.46\linewidth][l]{\sffamily\bfseries\color{subjectcolor!80}#1}%
  {\sffamily\footnotesize\itshape\color{black!52}#2}\par}

% --- NOTA / aparte único (sabías · dato · video · margen) --------------
\newtcolorbox{nota}[1][]{
  enhanced, breakable, sharp corners, boxrule=0pt, frame hidden,
  colback=black!4!white,
  borderline west={2.2pt}{0pt}{black!35},
  left=5mm, right=4mm, top=3mm, bottom=3mm,
  before skip=\bloqueskip, after skip=\bloqueskip,
  before upper={\ifblank{#1}{}{\blocklabel{black!55}{Nota}{#1}}}
}

% --- ALERTA / advertencia (única en ámbar, uso escaso) -----------------
\newtcolorbox{alerta}[1]{
  enhanced, breakable, sharp corners, boxrule=0pt, frame hidden,
  colback=alertcolor!7!white,
  borderline west={2.2pt}{0pt}{alertcolor},
  left=5mm, right=4mm, top=3mm, bottom=3mm,
  before skip=\bloqueskip, after skip=\bloqueskip,
  before upper={\blocklabel{alertcolor}{Atención}{#1}}
}
% Sub-ítem de alerta (antes \caso)
\newcommand{\caso}[2]{\par\vspace{0.35em}\noindent
  {\sffamily\bfseries\color{alertcolor!85!black}#1.}\enspace #2\par}

% --- CONTRASTE (la virtud y su exceso → pregunta que cuestiona) --------
%   \contraste{Concepto}{La virtud / lo valioso}{Su exceso / su riesgo}{Pregunta}
\newtcolorbox{contrastebox}{
  enhanced, breakable, sharp corners, boxrule=0pt, frame hidden,
  colback=subjectcolor!5!white,
  borderline west={2.2pt}{0pt}{subjectcolor},
  left=5mm, right=4mm, top=3mm, bottom=3mm,
  before skip=\bloqueskip, after skip=\bloqueskip,
}
\newcommand{\contraste}[4]{%
\begin{contrastebox}
\blocklabel{subjectcolor}{Cuestiona esto}{\faBalanceScale\ #1}
\noindent\begin{minipage}[t]{0.45\linewidth}
  {\sffamily\footnotesize\bfseries\color{subjectcolor!85!black}La virtud}\par
  \vspace{0.2em}#2
\end{minipage}\hfill
{\color{subjectcolor!35}\vrule width 0.4pt}\hfill
\begin{minipage}[t]{0.45\linewidth}
  {\sffamily\footnotesize\bfseries\color{alertcolor!85!black}Su exceso}\par
  \vspace{0.2em}#3
\end{minipage}
\par\vspace{0.6em}
{\color{subjectcolor!25}\titlerule[0.6pt]}\par\vspace{0.4em}
\noindent{\sffamily\itshape\color{subjectcolor}\faArrowRight\ #4}
\end{contrastebox}}

% --- EJEMPLO (sin caja: regla fina + etiqueta numerada) ----------------
\newcounter{ejemplo}[section]
\newenvironment{ejemplo}[1]
  {\par\addvspace{\bloqueskip}\refstepcounter{ejemplo}%
   \noindent{\color{subjectcolor!28}\rule{\linewidth}{0.5pt}}\par\vspace{0.4em}%
   \noindent{\sffamily\bfseries\color{subjectcolor}Ejemplo \theejemplo.}%
   \enspace{\sffamily\bfseries\color{black!78}#1}\par\vspace{0.3em}}
  {\par\addvspace{\bloqueskip}}

% --- PRÁCTICA (ejercicios / hoja / ICFES) ------------------------------
\newtcolorbox{practica}[1]{
  enhanced, breakable, sharp corners, boxrule=0pt, frame hidden,
  colback=white,
  borderline north={1.4pt}{0pt}{subjectcolor},
  left=0mm, right=0mm, top=3mm, bottom=2mm,
  before skip=\bloqueskip, after skip=\bloqueskip,
  before upper={\blocklabel{subjectcolor}{Práctica}{#1}\raggedright}
}

% Ejercicio numerado con nivel opcional [básico|intermedio|reto]
\newcounter{ejer}[section]
\newcommand{\ejer}[2][]{%
  \stepcounter{ejer}\par\vspace{0.5em}\noindent
  {\sffamily\bfseries\color{subjectcolor}\theejer.}\enspace #2%
  \ifblank{#1}{}{\ \textcolor{black!45}{\footnotesize[\textit{#1}]}}\par}

% Niveles = un solo color a tres intensidades (NO nuevos colores)
\newcommand{\nivel}[2]{\par\smallskip\noindent
  \textcolor{subjectcolor!#1}{\small\faCircle}\enspace
  {\sffamily\bfseries\footnotesize\color{black!70}#2}\par}
\newcommand{\basico}{\nivel{40}{Básico}}
\newcommand{\intermedio}{\nivel{70}{Intermedio}}
\newcommand{\reto}{\nivel{100}{Reto}}

% Selección múltiple tipo ICFES — opciones con burbuja (estilo Saber)
\newcounter{pregicfes}[section]
\newcommand{\bubbleletter}[1]{\tikz[baseline=(bc.base)]\node[circle,
  draw=subjectcolor!55, line width=0.6pt, inner sep=1.5pt,
  font=\sffamily\bfseries\scriptsize, text=subjectcolor!90!black](bc){#1};}
\newcommand{\icfes}[5]{%
  \stepcounter{pregicfes}\par\vspace{0.6em}\noindent
  {\sffamily\bfseries\color{subjectcolor}\thepregicfes.}\enspace #1\par
  \begin{enumerate}[label={\protect\bubbleletter{\Alph*}}, leftmargin=1.5cm,
                    itemsep=3pt, topsep=3pt]
    \item #2 \item #3 \item #4 \item #5
  \end{enumerate}}
% Pregunta abierta de argumentación
\newcommand{\pregabierta}[1]{\par\vspace{0.6em}\noindent
  {\sffamily\bfseries\color{subjectcolor}\faComment\ Argumenta:}\enspace
  \textit{#1}\par\lineaescritura\lineaescritura\lineaescritura}

% --- LÍNEAS DE RESPUESTA -----------------------------------------------
\newcommand{\linea}[1][7cm]{\underline{\hspace{#1}}}
\newcommand{\lineaescritura}{\par\vspace{0.9em}\noindent
  {\color{black!30}\rule{\linewidth}{0.3pt}}}

% =====================================================================
%  EVALUACIÓN TIPO ICFES — cierre de unidad (banner + 2 columnas)
%  Uso:  \begin{evaluacion} \begin{multicols}{2} \icfes... \end{multicols}
%        \pregabierta{...} \end{evaluacion}
% =====================================================================
\newtcolorbox{evaluacion}[1][cierre de sección]{
  enhanced, breakable, sharp corners,
  colback=white, colframe=subjectcolor, boxrule=1pt,
  colbacktitle=subjectcolor, coltitle=white,
  fonttitle=\InterDisplaySB\normalsize,
  title={\faClipboardCheck\ \ EVALUACIÓN TIPO ICFES\ \ %
         \textnormal{\sffamily\small· #1}},
  toptitle=2mm, bottomtitle=2mm, left=4mm, right=4mm, top=3.5mm, bottom=3mm,
  before={\clearpage},
  before upper={\small\itshape\color{black!55}%
    Selección múltiple. Marca una sola opción por pregunta.\par\medskip\raggedright}
}

% =====================================================================
%  HOJA DE TRABAJO — recortable: borde punteado, escala de grises,
%  en página aparte (clearpage antes; el contenido que sigue continúa).
% =====================================================================
\newtcolorbox{hoja}[1]{
  enhanced, breakable, sharp corners,
  colback=black!2!white, boxrule=0pt, frame hidden,
  borderline={0.8pt}{0pt}{black!50, dash pattern=on 3pt off 2.5pt},
  before={\clearpage}, after={\par},
  left=6mm, right=6mm, top=6mm, bottom=6mm,
  before upper={%
    {\sffamily\bfseries\footnotesize\color{black!60}\faCut\ \ HOJA DE TRABAJO}%
    {\color{black!35}\ \textbf{·}\ }{\sffamily\bfseries\footnotesize\color{black!80}#1}%
    \hfill{\sffamily\footnotesize\color{black!55}%
      Nombre:\ \underline{\hspace{4.2cm}}\quad Fecha:\ \underline{\hspace{2cm}}}\par
    \vspace{2mm}{\color{black!30}\hrule height 0.4pt}\vspace{3mm}\raggedright}
}

% =====================================================================
%  TABLAHOJA — tabla con TODAS las líneas (rejilla) para hojas de trabajo.
%  Celdas altas para llenar a mano. Centrada. El encabezado se marca con
%  \rowcolor en la primera fila.
%  Uso:  \begin{tablahoja}{ccc} enc & enc & enc \\ a & b & c \\ \end{tablahoja}
% =====================================================================
\NewDocumentEnvironment{tablahoja}{m}{%
  \par\smallskip\centering\small\setlength{\tabcolsep}{4pt}%
  \begin{NiceTabular}{#1}[hvlines, cell-space-top-limit=5pt,
     cell-space-bottom-limit=9pt, colortbl-like]%
}{%
  \end{NiceTabular}\par\smallskip
}

% =====================================================================
%  FIGURAS — imagen real o placeholder automático (\IfFileExists)
% =====================================================================
\newcounter{figura}[section]
\newcommand{\figura}[3][0.7\linewidth]{% [ancho]{archivo}{pie}
  \par\addvspace{\bloqueskip}\refstepcounter{figura}%
  \begin{center}%
  \IfFileExists{assets/imgs/#2}{%
    \includegraphics[width=#1]{assets/imgs/#2}}{%
    \begin{tikzpicture}
      \node[draw=black!30, dashed, line width=0.6pt, fill=black!3,
            text width=0.66\linewidth, minimum height=2.6cm,
            align=center, inner sep=8pt,
            text=black!55, font=\sffamily\footnotesize]
        {{\Large\faImage}\\[4pt]\textbf{[ imagen pendiente ]}\\#3};
    \end{tikzpicture}}%
  \par\vspace{0.4em}%
  {\footnotesize\sffamily\color{black!60}%
    \textbf{Figura \thefigura.}\enspace #3}%
  \end{center}\par\addvspace{\bloqueskip}}

% =====================================================================
%  APERTURA DE SECCIÓN — marco branded de serie + título + "qué verás"
%  El "marco" es CONSTANTE en todas las secciones (= identidad de marca):
%  número grande de serie, kicker, cuadro-firma, franja de acento y sello
%  del instituto. El arte (imagen real o placeholder) va DENTRO del marco.
%  Uso (1ª línea de cada sNN.tex):
%    \aperturaseccion{ruta/img.png}{Título}{Una o dos frases.}
%  Personaliza por documento (opcional):  \renewcommand{\seccionkicker}{...}
% =====================================================================
\newcounter{secportada}
\newcounter{unidadnum}                       % nº de unidad actual (lo fija \aperturacapitulo)
% Número visible de sección = Unidad.Sección (ej. 1.2). Si no hay unidad
% definida (unidadnum=0), cae al número simple de sección.
\newcommand{\numseccion}{\ifnum\value{unidadnum}>0 \theunidadnum.\thesecportada\else\thesecportada\fi}
\newcommand{\seccionkicker}{\hdrcapitulo}   % por defecto, lo del encabezado

% El escenario branded (contenedor constante del arte)
\newtcolorbox{secstage}{%
  enhanced, sharp corners, boxrule=0pt, frame hidden,
  colback=subjectcolor!6, halign=center, valign=center,
  height=5.0cm, left=2mm, right=2mm, top=10mm, bottom=8mm,
  overlay={%
    % cuadro-firma + kicker (arriba-izquierda)
    \fill[subjectcolor] ([shift={(4mm,-4mm)}]frame.north west) rectangle ++(2.6mm,-2.6mm);
    \node[anchor=north west, font=\InterDisplaySB\footnotesize, text=subjectcolor!85]
      at ([shift={(8.5mm,-3.4mm)}]frame.north west) {\seccionkicker};
    % número de serie grande y tenue (arriba-derecha) = marca
    \node[anchor=north east, font=\InterDisplayBlk, text=subjectcolor!18, inner sep=0pt]
      at ([shift={(-4mm,-1.5mm)}]frame.north east)
      {\fontsize{34}{34}\selectfont \numseccion};
    % franja de acento (abajo-izquierda)
    \fill[subjectcolor] ([shift={(4mm,4mm)}]frame.south west) rectangle ++(2.8cm,2.2mm);
    % sello del instituto (abajo-derecha)
    \node[anchor=south east, font=\sffamily\scriptsize, text=subjectcolor!55]
      at ([shift={(-4mm,3.4mm)}]frame.south east) {ITSTZ · Zapatoca};
  }}

\newcommand{\aperturaseccion}[3]{% {archivo}{título}{intro}
  \clearpage\refstepcounter{secportada}%
  \begin{secstage}
  \IfFileExists{assets/imgs/#1}{%
    \includegraphics[width=0.9\linewidth, height=3.2cm, keepaspectratio]{assets/imgs/#1}}{%
    \begin{tikzpicture}
      \node[draw=subjectcolor!35, dashed, line width=0.6pt, fill=white,
            minimum width=0.86\linewidth, minimum height=2.9cm,
            align=center, text=subjectcolor!50, font=\sffamily\small]
        {{\fontsize{22}{22}\selectfont\faImage}\\[5pt]
         \textbf{[ imagen pendiente ]}\\[1pt]
         \scriptsize\ttfamily assets/imgs/#1};
    \end{tikzpicture}}%
  \end{secstage}
  \vspace{0.5em}
  {\color{subjectcolor}\rule{3.2cm}{2.5pt}}\par\vspace{0.3em}
  \section{{\color{subjectcolor}\numseccion}\enspace #2}%
  {\InterLight\fontsize{12.5}{17}\selectfont\color{black!58} #3\par}
  \vspace{0.4em}}

% =====================================================================
%  DESTACADO — "lo esencial": triángulo lateral + degradado muy sutil
% =====================================================================
\newtcolorbox{destacado}[1]{
  enhanced, breakable, sharp corners, boxrule=0pt, frame hidden,
  interior style={left color=subjectcolor!16, right color=subjectcolor!2},
  borderline west={2.4pt}{0pt}{subjectcolor},
  left=7mm, right=5mm, top=3.5mm, bottom=3.5mm,
  before skip=\bloqueskip, after skip=\bloqueskip,
  overlay unbroken and first={%
    \fill[subjectcolor] ([yshift=-4.5mm]frame.north west)
        -- ++(5.5pt,2.75pt) -- ++(0,-5.5pt) -- cycle;},
  before upper={\blocklabel{subjectcolor}{Lo esencial}{#1}}
}

% --- Resaltado inline suave ---
\newtcbox{\resaltar}{on line, nobeforeafter, boxsep=0pt,
  left=2.5pt, right=2.5pt, top=1pt, bottom=1pt, arc=2.5pt,
  colback=subjectcolor!15, colframe=subjectcolor!15}

% =====================================================================
%  EXTENSIONES SANTILLANA (opt-in) — cajas temáticas para dar variedad
%  visual y romper la monotonía. NO alteran los entornos base; úsalas
%  donde aporten. Mantienen el color de asignatura (coherencia de marca).
% =====================================================================

% --- ¿SABÍAS QUE? — curiosidad breve · tarjeta redondeada con bombillo --
\newtcolorbox{sabias}{
  enhanced, breakable, arc=2.5mm, boxrule=0.8pt,
  colback=subjectcolor!6!white, colframe=subjectcolor!30,
  left=4mm, right=4mm, top=3mm, bottom=3mm,
  before skip=\bloqueskip, after skip=\bloqueskip,
  before upper={{\sffamily\bfseries\footnotesize\color{subjectcolor}%
    \faLightbulb\enspace\MakeUppercase{¿Sabías que?}}%
    \par\nopagebreak\vspace{0.45em}}
}

% --- CIENCIA Y SOCIEDAD — impacto/ética · banner de título -------------
\newtcolorbox{cienciasociedad}[1]{
  enhanced, breakable, sharp corners,
  colback=subjectcolor!5!white, colframe=subjectcolor, boxrule=0pt,
  colbacktitle=subjectcolor, coltitle=white,
  fonttitle=\sffamily\bfseries\footnotesize,
  title={\faGlobeAmericas\ \ \MakeUppercase{Ciencia y sociedad}%
         \ifblank{#1}{}{\ \ \textnormal{\sffamily\small· #1}}},
  toptitle=1.6mm, bottomtitle=1.6mm,
  left=4mm, right=4mm, top=3mm, bottom=3mm,
  before skip=\bloqueskip, after skip=\bloqueskip,
}

% --- MANOS A LA OBRA — práctica experimental · marco sólido con matraz --
\newtcolorbox{laboratorio}[1]{
  enhanced, breakable, sharp corners, arc=0pt,
  colback=white, colframe=subjectcolor!55, boxrule=1pt,
  left=4mm, right=4mm, top=3mm, bottom=3mm,
  before skip=\bloqueskip, after skip=\bloqueskip,
  before upper={{\sffamily\bfseries\footnotesize\color{subjectcolor}%
    \faFlask\enspace\MakeUppercase{Manos a la obra}}%
    \ifblank{#1}{}{\ {\color{black!35}\textbf{·}}\ %
      {\sffamily\bfseries\footnotesize\color{black!78}#1}}%
    \par\nopagebreak\vspace{0.45em}}
}
% Lista de materiales dentro de \begin{laboratorio}
\newcommand{\materiales}[1]{\par\noindent
  {\sffamily\bfseries\footnotesize\color{black!70}Materiales:}\ %
  {\footnotesize #1}\par\vspace{0.3em}}

% --- IMAGEN INLINE con texto alrededor (Santillana) -------------------
%   \figinline[r|l]{ancho}{archivo}{pie}   — colócala al inicio de un párrafo
\newcommand{\figinline}[4][r]{%
  \begin{wrapfigure}{#1}{#2}
  \centering\vspace{-0.5\baselineskip}
  \IfFileExists{assets/imgs/#3}{%
    \includegraphics[width=\linewidth]{assets/imgs/#3}}{%
    \fbox{\parbox[c][3cm][c]{0.92\linewidth}{\centering\sffamily\scriptsize
      \color{black!55}{\Large\faImage}\\[3pt][imagen pendiente]\\\ttfamily\tiny #3}}}%
  \par\vspace{2pt}{\footnotesize\sffamily\color{black!60}\textbf{Fig.}\enspace #4}%
  \vspace{-0.4\baselineskip}
  \end{wrapfigure}}

% --- APERTURA CON HERO LATERAL (imagen grande al lado del título) ------
%   \aperturahero{archivo}{título}{intro}  — alternativa a \aperturaseccion
\newcommand{\aperturahero}[3]{%
  \clearpage\refstepcounter{secportada}\stepcounter{section}%
  \addcontentsline{toc}{section}{\numseccion\enspace #2}%
  \noindent\begin{minipage}[c]{0.62\linewidth}
    {\InterDisplaySB\footnotesize\color{subjectcolor!85}\MakeUppercase{\seccionkicker}}\par
    \vspace{1.8mm}
    {\LARGE\InterDisplaySB\color{subjectcolor}%
      {\color{subjectcolor!35}\numseccion}\,\,#2}\par
    \vspace{2mm}{\color{subjectcolor!30}\rule{\linewidth}{0.8pt}}\par\vspace{2mm}
    {\InterLight\fontsize{12}{16}\selectfont\color{black!58}#3}
  \end{minipage}\hfill
  \begin{minipage}[c]{0.34\linewidth}\centering
    \IfFileExists{assets/imgs/#1}{%
      \includegraphics[width=\linewidth]{assets/imgs/#1}}{%
      \fbox{\parbox[c][3.4cm][c]{0.95\linewidth}{\centering\sffamily\scriptsize
        \color{black!55}{\Large\faImage}\\[3pt][imagen pendiente]\\\ttfamily\tiny #1}}}%
  \end{minipage}\par\vspace{1.3em}}

% --- MAPAS (conceptual y mental) — estilos TikZ con texto nítido -------
\tikzset{
  croot/.style={draw=subjectcolor, fill=subjectcolor, text=white, rounded corners=3pt,
    font=\sffamily\bfseries\footnotesize, align=center, inner sep=5pt, minimum height=9mm},
  cnode/.style={draw=subjectcolor!55, fill=subjectcolor!8, rounded corners=2.5pt,
    font=\sffamily\footnotesize, align=center, inner sep=4pt, text width=2.5cm},
  cleaf/.style={draw=subjectcolor!45, fill=white, rounded corners=2.5pt,
    font=\sffamily\footnotesize, align=center, inner sep=4pt},
  clink/.style={-{Stealth[length=5pt]}, subjectcolor!65, semithick},
  cline/.style={subjectcolor!55, semithick},
}
% Caja contenedora para un mapa (título + regla superior)
\newtcolorbox{mapabox}[2][Mapa conceptual]{
  enhanced, sharp corners, boxrule=0pt, frame hidden, colback=white,
  borderline north={1.4pt}{0pt}{subjectcolor},
  left=1mm, right=1mm, top=3mm, bottom=2mm,
  before skip=\bloqueskip, after skip=\bloqueskip,
  before upper={\blocklabel{subjectcolor}{#1}{#2}}
}

% --- Papel cuadriculado para graficar a mano (hojas de trabajo) --------
%   \papelcuadricula[filas]{ancho en cm}   (cada celda = 0.6 cm)
\newcommand{\papelcuadricula}[2][8]{%
  \par\smallskip\begin{center}
  \begin{tikzpicture}
    \draw[step=0.6cm, black!16, line width=0.3pt] (0,0) grid (#2,{#1*0.6});
    \draw[black!45, line width=0.7pt] (0,0) -- (#2,0);
    \draw[black!45, line width=0.7pt] (0,0) -- (0,{#1*0.6});
  \end{tikzpicture}
  \end{center}\par\smallskip}

% =====================================================================
%  APERTURA DE CAPÍTULO — título grande, número marca de agua, aire
% =====================================================================
\newcommand{\aperturacapitulo}[3]{% {número}{título}{subtítulo}
  \setcounter{unidadnum}{#1}\setcounter{secportada}{0}% fija unidad y reinicia secciones
  \thispagestyle{empty}%
  \begin{tikzpicture}[remember picture, overlay]
    \node[anchor=north east, text=subjectcolor!9,
          font=\InterDisplayThin]
      at ([xshift=-0.9cm, yshift=-2.2cm]current page.north east)
      {\fontsize{210}{210}\selectfont #1};
  \end{tikzpicture}%
  \begingroup\raggedright\null\vspace{3.8cm}\par
  {\InterDisplaySB\fontsize{12}{14}\selectfont\color{subjectcolor}%
    CAP\'ITULO\ #1}\par
  \vspace{0.18cm}{\color{black!16}\rule{4cm}{1pt}}\par
  \vspace{0.9cm}
  {\InterDisplaySB\fontsize{46}{50}\selectfont\color{black!88}#2}\par
  \vspace{0.5cm}{\color{subjectcolor}\rule{3.2cm}{2.5pt}}\par
  \vspace{0.55cm}
  {\InterLight\fontsize{15}{21}\selectfont\color{black!55}#3}\par
  \vfill
  {\sffamily\footnotesize\color{black!45}\hdrinstituto\quad
    \textcolor{black!22}{·}\quad Prof.\ Juan Diego A.\ Prada R.}\par
  \endgroup\clearpage}

% =====================================================================
%  DESBLOQUEA — pie discreto, no caja
% =====================================================================
\newcommand{\desbloquea}[1]{%
  \par\addvspace{\bloqueskip}%
  \noindent{\color{subjectcolor!25}\rule{\linewidth}{0.5pt}}\par\vspace{0.4em}%
  {\footnotesize\sffamily\color{black!60}%
    \textcolor{subjectcolor}{\faUnlock}\ \textbf{Dominar esto te desbloquea:}\ #1}\par}
):
%     % =====================================================================
%  TEMA · CIENCIAS NATURALES — variante del design system (GUÍAS)
%  Identidad visual propia, distinta de Matemáticas (azul/pizarra):
%    · color base ESMERALDA (\setsubject{cnat})
%    · acento LIMA orgánico solo para ornamentos (no rompe la regla de
%      "un solo color de asignatura")
%    · iconografía de ciencias (hoja · ADN · matraz · átomo)
%    · reutiliza las cajas de ciencia que ya trae el preamble:
%      \begin{laboratorio}, \begin{sabias}, \begin{cienciasociedad}
%
%  Sirve a futuro para BIOLOGÍA, QUÍMICA y FÍSICA EXPERIMENTAL.
%
%  INVOCACIÓN (una sola línea, tras % =====================================================================
%  DESIGN SYSTEM v5 — EDITORIAL ITSTZ
%  Filosofía: texto primero · un solo color de asignatura · sin sombras
%  · sin marcos pesados · tratamientos planos (filete lateral / regla).
%
%  Vocabulario semántico (esto es TODO lo disponible):
%    \section / \subsection      jerarquía
%    definicion{título}          concepto/regla clave (filete color)
%    ejemplo{título}             ejemplo resuelto (numerado, sin caja)
%    procedimiento{título}       pasos (filete neutro) + \paso{n}{texto}
%    nota[título]                aparte único gris (sabías/dato/video/margen)
%    alerta{título}              advertencia ámbar (usar rara vez)
%    practica{título}            ejercicios/hoja/ICFES (regla superior)
%    \figura{archivo}{pie}       imagen real o placeholder automático
%    \desbloquea{lista}          pie discreto de prerrequisitos
%
%  REGLA DE ORO: las secciones (sNN.tex) NO llevan \vspace, color ni
%  minipage de maquetación. Solo contenido semántico. El diseño vive aquí.
% =====================================================================
\documentclass[11pt, letterpaper, twoside]{article}

\usepackage{fontspec}
\usepackage{polyglossia}
\setdefaultlanguage{spanish}
\usepackage[inner=2.1cm, outer=3.7cm, top=2.6cm, bottom=2.3cm,
            marginparwidth=2.55cm, marginparsep=0.45cm,
            headheight=16pt, headsep=18pt, footskip=24pt]{geometry}
\usepackage{microtype}
\usepackage{multicol}
\usepackage{xcolor}
\usepackage{graphicx}
\usepackage{tikz}
\usetikzlibrary{calc, arrows.meta, positioning}
\usepackage{wrapfig}
\usepackage[most]{tcolorbox}
\usepackage{amsmath, amssymb}
\usepackage{siunitx}
\usepackage{fontawesome5}
\usepackage{fancyhdr}
\usepackage{titlesec}
\usepackage{enumitem}
\usepackage{etoolbox}
\usepackage{titletoc}
\usepackage{array}
\usepackage{booktabs}
\usepackage{colortbl}
\usepackage{nicematrix}   % rejilla completa (hvlines) para tablas de hojas

% --- TIPOGRAFÍA -------------------------------------------------------
% Cuerpo: Charter (serif cálida y muy legible). La matemática (Latin
% Modern, serif) combina con ella. Títulos/etiquetas/UI: Inter Display.
\setmainfont{Charter}[
  UprightFont    = {Charter},
  BoldFont       = {Charter Bold},
  ItalicFont     = {Charter Italic},
  BoldItalicFont = {Charter Bold Italic},
  Ligatures = TeX, Numbers = Lining ]
\setsansfont{Inter}[
  UprightFont = Inter-Regular, BoldFont = Inter-SemiBold,
  ItalicFont  = Inter-Italic, Ligatures = TeX, Numbers = Lining ]
\newfontfamily\InterLight{Inter-Light}[ItalicFont=Inter-LightItalic, Ligatures=TeX]
\newfontfamily\InterDisplay{Inter Display}[Ligatures=TeX]
\newfontfamily\InterDisplaySB{Inter Display SemiBold}[Ligatures=TeX]
\newfontfamily\InterDisplayThin{Inter Display Thin}[Ligatures=TeX]
\newfontfamily\InterDisplayBlk{Inter Display Black}[Ligatures=TeX]

% --- RITMO --------------------------------------------------------------
\linespread{1.22}
\setlength{\parindent}{0pt}
\setlength{\parskip}{0.5em}
\newlength{\bloqueskip}\setlength{\bloqueskip}{1.0em}   % aire entre bloques
\setlength{\columnsep}{1.6em}                            % medianil multicols

% =====================================================================
%  COLOR — UNA sola variable de asignatura + neutros. Sin arcoíris.
% =====================================================================
\definecolor{mat} {RGB}{ 30,  64, 120}   % azul profundo
\definecolor{cnat}{RGB}{ 16, 110,  80}   % esmeralda
\definecolor{csoc}{RGB}{165,  75,  45}   % terracota
\definecolor{eco} {RGB}{ 70,  55, 130}   % índigo
\definecolor{alertcolor}{RGB}{180, 95, 10}  % ámbar — solo advertencias

\colorlet{subjectcolor}{mat}             % por defecto
\newcommand{\setsubject}[1]{\colorlet{subjectcolor}{#1}}

% =====================================================================
%  JERARQUÍA TIPOGRÁFICA
% =====================================================================
\setcounter{secnumdepth}{0}
\titleformat{\section}
  {\LARGE\InterDisplaySB\color{subjectcolor}}{}{0pt}{}
  [\vspace{2pt}{\color{subjectcolor!30}\titlerule[0.8pt]}]
\titlespacing*{\section}{0pt}{3.2ex plus 0.6ex}{1.8ex plus 0.3ex}

\titleformat{\subsection}
  {\large\InterDisplaySB\color{black!82}}{}{0pt}{}
\titlespacing*{\subsection}{0pt}{2.0ex plus 0.4ex}{0.5ex}

% --- TABLA DE CONTENIDO -------------------------------------------------
\titlecontents{section}[0em]
  {\addvspace{0.45pc}\sffamily\bfseries\small\color{black!85}}
  {}{}{\titlerule*[0.35pc]{.}\contentspage}
\titlecontents{subsection}[1.4em]
  {\sffamily\footnotesize\color{black!60}}
  {}{}{\titlerule*[0.35pc]{.}\contentspage}

% =====================================================================
%  ENCABEZADO / PIE — editorial limpio (regla fina, sin franjas)
% =====================================================================
\newcommand{\hdrinstituto}{Instituto Técnico Santo Tomás · Zapatoca}
\newcommand{\hdrcapitulo}{}            % se fija en el documento maestro
\pagestyle{fancy}\fancyhf{}
\renewcommand{\headrulewidth}{0pt}
\fancyhead[LE]{\footnotesize\sffamily\color{black!55}\hdrinstituto}
\fancyhead[RO]{\footnotesize\sffamily\color{subjectcolor}\hdrcapitulo}
\fancyhead[LO,RE]{}
\renewcommand{\headrule}{\vspace{2pt}{\color{black!18}\hrule height 0.4pt}}

% --- PIE: número exterior (cuadro de color) + crédito interior + filete ---
\renewcommand{\footrulewidth}{0pt}
\renewcommand{\footrule}{{\color{black!15}\hrule height 0.4pt}\vspace{3pt}}
\newcommand{\pagenumbox}{%
  \tikz[baseline=-3.2pt]{\node[fill=subjectcolor, text=white, inner xsep=7pt,
     inner ysep=3pt, font=\sffamily\bfseries\footnotesize]{\thepage};}}
\fancyfoot[RO,LE]{\pagenumbox}
\fancyfoot[LO,RE]{\footnotesize\sffamily\color{black!45}%
  \hdrcapitulo\ \textcolor{black!22}{·}\ Prof.\ Juan Diego A.\ Prada R.}

% =====================================================================
%  HELPERS
% =====================================================================
% Etiqueta de bloque:  TAG · título   (color del tag configurable)
\newcommand{\blocklabel}[3]{% #1=color  #2=TAG  #3=título
  {\sffamily\bfseries\footnotesize\color{#1}\MakeUppercase{#2}}%
  \ifblank{#3}{}{\,{\color{black!35}\textbf{·}}\,%
    {\sffamily\bfseries\footnotesize\color{black!78}#3}}%
  \par\nopagebreak\vspace{0.4em}}

% --- CAJAS AL MARGEN — personajes históricos e hitos -------------------
% Van en el margen exterior ancho (\marginpar respeta la geometría twoside
% y flota para no chocar). Son contenido semántico: colócalas en el cuerpo,
% cerca del párrafo que comentan, ENTRE párrafos. NO usarlas dentro de
% multicols, tcolorbox, tablas, ni en la misma línea que un \section.
%   \personaje{Nombre}{años · rol}{descripción breve}
%   \hito{año}{qué pasó (breve)}
\setlength{\marginparpush}{8pt}
\newcommand{\margenbox}[1]{\marginpar{%
  \begingroup\raggedright\footnotesize\sffamily
  \setlength{\parskip}{0.18em}\linespread{1.05}\selectfont
  #1\par\endgroup}}
\newcommand{\personaje}[3]{\margenbox{%
  {\color{subjectcolor}\rule{\linewidth}{1.1pt}}\par
  {\InterDisplaySB\scriptsize\color{subjectcolor}#1}\par
  {\fontsize{6.8pt}{8pt}\selectfont\itshape\color{black!50}#2}\par
  \vspace{0.15em}{\scriptsize\color{black!72}#3}}}
\newcommand{\hito}[2]{\margenbox{%
  {\sffamily\bfseries\scriptsize\colorbox{subjectcolor}{\textcolor{white}{\,#1\,}}}\par
  \vspace{0.15em}{\scriptsize\color{black!72}#2}}}

% =====================================================================
%  ENTORNOS — todos planos: filete lateral, sin marco, sin sombra
% =====================================================================

% --- DEFINICIÓN / concepto clave ---------------------------------------
\newtcolorbox{definicion}[1]{
  enhanced, breakable, sharp corners, boxrule=0pt, frame hidden,
  colback=subjectcolor!6!white,
  borderline west={2.2pt}{0pt}{subjectcolor},
  left=5mm, right=4mm, top=3mm, bottom=3mm,
  before skip=\bloqueskip, after skip=\bloqueskip,
  before upper={\blocklabel{subjectcolor}{Definición}{#1}}
}

% --- PROCEDIMIENTO / pasos ---------------------------------------------
\newtcolorbox{procedimiento}[1]{
  enhanced, breakable, sharp corners, boxrule=0pt, frame hidden,
  colback=black!3!white,
  borderline west={2.2pt}{0pt}{black!45},
  left=5mm, right=4mm, top=3mm, bottom=3mm,
  before skip=\bloqueskip, after skip=\bloqueskip,
  before upper={\blocklabel{black!60}{Procedimiento}{#1}}
}
\newcommand{\paso}[2]{%
  \par\vspace{0.45em}\noindent
  \tikz[baseline=-0.6ex]\node[circle, fill=subjectcolor, text=white,
     inner sep=0pt, minimum size=1.45em,
     font=\sffamily\bfseries\footnotesize]{#1};\hspace{0.5em}#2\par}

% Glosa "¿por qué?" bajo un paso — narración en lenguaje cercano
\newcommand{\porque}[1]{\par\vspace{0.1em}\noindent\hspace{2.0em}%
  {\sffamily\footnotesize\itshape\color{subjectcolor}¿por qué?\,}%
  {\sffamily\footnotesize\itshape\color{black!55}#1}\par}

% Operación glosada: cuenta a la izquierda + porqué al margen derecho
% (estilo "comentario al lado del código"). Úsala dentro de ejemplo.
% \opg{ecuación}{porqué}   — primer arg en modo display (mat, cap0)
% \opgt{etiqueta}{porqué}  — primer arg en texto sans bold (física, ciencias)
\newcommand{\opg}[2]{\par\vspace{0.25em}\noindent
  \makebox[0.46\linewidth][l]{$\displaystyle #1$}%
  {\sffamily\footnotesize\itshape\color{black!52}#2}\par}
\newcommand{\opgt}[2]{\par\vspace{0.25em}\noindent
  \makebox[0.46\linewidth][l]{\sffamily\bfseries\color{subjectcolor!80}#1}%
  {\sffamily\footnotesize\itshape\color{black!52}#2}\par}

% --- NOTA / aparte único (sabías · dato · video · margen) --------------
\newtcolorbox{nota}[1][]{
  enhanced, breakable, sharp corners, boxrule=0pt, frame hidden,
  colback=black!4!white,
  borderline west={2.2pt}{0pt}{black!35},
  left=5mm, right=4mm, top=3mm, bottom=3mm,
  before skip=\bloqueskip, after skip=\bloqueskip,
  before upper={\ifblank{#1}{}{\blocklabel{black!55}{Nota}{#1}}}
}

% --- ALERTA / advertencia (única en ámbar, uso escaso) -----------------
\newtcolorbox{alerta}[1]{
  enhanced, breakable, sharp corners, boxrule=0pt, frame hidden,
  colback=alertcolor!7!white,
  borderline west={2.2pt}{0pt}{alertcolor},
  left=5mm, right=4mm, top=3mm, bottom=3mm,
  before skip=\bloqueskip, after skip=\bloqueskip,
  before upper={\blocklabel{alertcolor}{Atención}{#1}}
}
% Sub-ítem de alerta (antes \caso)
\newcommand{\caso}[2]{\par\vspace{0.35em}\noindent
  {\sffamily\bfseries\color{alertcolor!85!black}#1.}\enspace #2\par}

% --- CONTRASTE (la virtud y su exceso → pregunta que cuestiona) --------
%   \contraste{Concepto}{La virtud / lo valioso}{Su exceso / su riesgo}{Pregunta}
\newtcolorbox{contrastebox}{
  enhanced, breakable, sharp corners, boxrule=0pt, frame hidden,
  colback=subjectcolor!5!white,
  borderline west={2.2pt}{0pt}{subjectcolor},
  left=5mm, right=4mm, top=3mm, bottom=3mm,
  before skip=\bloqueskip, after skip=\bloqueskip,
}
\newcommand{\contraste}[4]{%
\begin{contrastebox}
\blocklabel{subjectcolor}{Cuestiona esto}{\faBalanceScale\ #1}
\noindent\begin{minipage}[t]{0.45\linewidth}
  {\sffamily\footnotesize\bfseries\color{subjectcolor!85!black}La virtud}\par
  \vspace{0.2em}#2
\end{minipage}\hfill
{\color{subjectcolor!35}\vrule width 0.4pt}\hfill
\begin{minipage}[t]{0.45\linewidth}
  {\sffamily\footnotesize\bfseries\color{alertcolor!85!black}Su exceso}\par
  \vspace{0.2em}#3
\end{minipage}
\par\vspace{0.6em}
{\color{subjectcolor!25}\titlerule[0.6pt]}\par\vspace{0.4em}
\noindent{\sffamily\itshape\color{subjectcolor}\faArrowRight\ #4}
\end{contrastebox}}

% --- EJEMPLO (sin caja: regla fina + etiqueta numerada) ----------------
\newcounter{ejemplo}[section]
\newenvironment{ejemplo}[1]
  {\par\addvspace{\bloqueskip}\refstepcounter{ejemplo}%
   \noindent{\color{subjectcolor!28}\rule{\linewidth}{0.5pt}}\par\vspace{0.4em}%
   \noindent{\sffamily\bfseries\color{subjectcolor}Ejemplo \theejemplo.}%
   \enspace{\sffamily\bfseries\color{black!78}#1}\par\vspace{0.3em}}
  {\par\addvspace{\bloqueskip}}

% --- PRÁCTICA (ejercicios / hoja / ICFES) ------------------------------
\newtcolorbox{practica}[1]{
  enhanced, breakable, sharp corners, boxrule=0pt, frame hidden,
  colback=white,
  borderline north={1.4pt}{0pt}{subjectcolor},
  left=0mm, right=0mm, top=3mm, bottom=2mm,
  before skip=\bloqueskip, after skip=\bloqueskip,
  before upper={\blocklabel{subjectcolor}{Práctica}{#1}\raggedright}
}

% Ejercicio numerado con nivel opcional [básico|intermedio|reto]
\newcounter{ejer}[section]
\newcommand{\ejer}[2][]{%
  \stepcounter{ejer}\par\vspace{0.5em}\noindent
  {\sffamily\bfseries\color{subjectcolor}\theejer.}\enspace #2%
  \ifblank{#1}{}{\ \textcolor{black!45}{\footnotesize[\textit{#1}]}}\par}

% Niveles = un solo color a tres intensidades (NO nuevos colores)
\newcommand{\nivel}[2]{\par\smallskip\noindent
  \textcolor{subjectcolor!#1}{\small\faCircle}\enspace
  {\sffamily\bfseries\footnotesize\color{black!70}#2}\par}
\newcommand{\basico}{\nivel{40}{Básico}}
\newcommand{\intermedio}{\nivel{70}{Intermedio}}
\newcommand{\reto}{\nivel{100}{Reto}}

% Selección múltiple tipo ICFES — opciones con burbuja (estilo Saber)
\newcounter{pregicfes}[section]
\newcommand{\bubbleletter}[1]{\tikz[baseline=(bc.base)]\node[circle,
  draw=subjectcolor!55, line width=0.6pt, inner sep=1.5pt,
  font=\sffamily\bfseries\scriptsize, text=subjectcolor!90!black](bc){#1};}
\newcommand{\icfes}[5]{%
  \stepcounter{pregicfes}\par\vspace{0.6em}\noindent
  {\sffamily\bfseries\color{subjectcolor}\thepregicfes.}\enspace #1\par
  \begin{enumerate}[label={\protect\bubbleletter{\Alph*}}, leftmargin=1.5cm,
                    itemsep=3pt, topsep=3pt]
    \item #2 \item #3 \item #4 \item #5
  \end{enumerate}}
% Pregunta abierta de argumentación
\newcommand{\pregabierta}[1]{\par\vspace{0.6em}\noindent
  {\sffamily\bfseries\color{subjectcolor}\faComment\ Argumenta:}\enspace
  \textit{#1}\par\lineaescritura\lineaescritura\lineaescritura}

% --- LÍNEAS DE RESPUESTA -----------------------------------------------
\newcommand{\linea}[1][7cm]{\underline{\hspace{#1}}}
\newcommand{\lineaescritura}{\par\vspace{0.9em}\noindent
  {\color{black!30}\rule{\linewidth}{0.3pt}}}

% =====================================================================
%  EVALUACIÓN TIPO ICFES — cierre de unidad (banner + 2 columnas)
%  Uso:  \begin{evaluacion} \begin{multicols}{2} \icfes... \end{multicols}
%        \pregabierta{...} \end{evaluacion}
% =====================================================================
\newtcolorbox{evaluacion}[1][cierre de sección]{
  enhanced, breakable, sharp corners,
  colback=white, colframe=subjectcolor, boxrule=1pt,
  colbacktitle=subjectcolor, coltitle=white,
  fonttitle=\InterDisplaySB\normalsize,
  title={\faClipboardCheck\ \ EVALUACIÓN TIPO ICFES\ \ %
         \textnormal{\sffamily\small· #1}},
  toptitle=2mm, bottomtitle=2mm, left=4mm, right=4mm, top=3.5mm, bottom=3mm,
  before={\clearpage},
  before upper={\small\itshape\color{black!55}%
    Selección múltiple. Marca una sola opción por pregunta.\par\medskip\raggedright}
}

% =====================================================================
%  HOJA DE TRABAJO — recortable: borde punteado, escala de grises,
%  en página aparte (clearpage antes; el contenido que sigue continúa).
% =====================================================================
\newtcolorbox{hoja}[1]{
  enhanced, breakable, sharp corners,
  colback=black!2!white, boxrule=0pt, frame hidden,
  borderline={0.8pt}{0pt}{black!50, dash pattern=on 3pt off 2.5pt},
  before={\clearpage}, after={\par},
  left=6mm, right=6mm, top=6mm, bottom=6mm,
  before upper={%
    {\sffamily\bfseries\footnotesize\color{black!60}\faCut\ \ HOJA DE TRABAJO}%
    {\color{black!35}\ \textbf{·}\ }{\sffamily\bfseries\footnotesize\color{black!80}#1}%
    \hfill{\sffamily\footnotesize\color{black!55}%
      Nombre:\ \underline{\hspace{4.2cm}}\quad Fecha:\ \underline{\hspace{2cm}}}\par
    \vspace{2mm}{\color{black!30}\hrule height 0.4pt}\vspace{3mm}\raggedright}
}

% =====================================================================
%  TABLAHOJA — tabla con TODAS las líneas (rejilla) para hojas de trabajo.
%  Celdas altas para llenar a mano. Centrada. El encabezado se marca con
%  \rowcolor en la primera fila.
%  Uso:  \begin{tablahoja}{ccc} enc & enc & enc \\ a & b & c \\ \end{tablahoja}
% =====================================================================
\NewDocumentEnvironment{tablahoja}{m}{%
  \par\smallskip\centering\small\setlength{\tabcolsep}{4pt}%
  \begin{NiceTabular}{#1}[hvlines, cell-space-top-limit=5pt,
     cell-space-bottom-limit=9pt, colortbl-like]%
}{%
  \end{NiceTabular}\par\smallskip
}

% =====================================================================
%  FIGURAS — imagen real o placeholder automático (\IfFileExists)
% =====================================================================
\newcounter{figura}[section]
\newcommand{\figura}[3][0.7\linewidth]{% [ancho]{archivo}{pie}
  \par\addvspace{\bloqueskip}\refstepcounter{figura}%
  \begin{center}%
  \IfFileExists{assets/imgs/#2}{%
    \includegraphics[width=#1]{assets/imgs/#2}}{%
    \begin{tikzpicture}
      \node[draw=black!30, dashed, line width=0.6pt, fill=black!3,
            text width=0.66\linewidth, minimum height=2.6cm,
            align=center, inner sep=8pt,
            text=black!55, font=\sffamily\footnotesize]
        {{\Large\faImage}\\[4pt]\textbf{[ imagen pendiente ]}\\#3};
    \end{tikzpicture}}%
  \par\vspace{0.4em}%
  {\footnotesize\sffamily\color{black!60}%
    \textbf{Figura \thefigura.}\enspace #3}%
  \end{center}\par\addvspace{\bloqueskip}}

% =====================================================================
%  APERTURA DE SECCIÓN — marco branded de serie + título + "qué verás"
%  El "marco" es CONSTANTE en todas las secciones (= identidad de marca):
%  número grande de serie, kicker, cuadro-firma, franja de acento y sello
%  del instituto. El arte (imagen real o placeholder) va DENTRO del marco.
%  Uso (1ª línea de cada sNN.tex):
%    \aperturaseccion{ruta/img.png}{Título}{Una o dos frases.}
%  Personaliza por documento (opcional):  \renewcommand{\seccionkicker}{...}
% =====================================================================
\newcounter{secportada}
\newcounter{unidadnum}                       % nº de unidad actual (lo fija \aperturacapitulo)
% Número visible de sección = Unidad.Sección (ej. 1.2). Si no hay unidad
% definida (unidadnum=0), cae al número simple de sección.
\newcommand{\numseccion}{\ifnum\value{unidadnum}>0 \theunidadnum.\thesecportada\else\thesecportada\fi}
\newcommand{\seccionkicker}{\hdrcapitulo}   % por defecto, lo del encabezado

% El escenario branded (contenedor constante del arte)
\newtcolorbox{secstage}{%
  enhanced, sharp corners, boxrule=0pt, frame hidden,
  colback=subjectcolor!6, halign=center, valign=center,
  height=5.0cm, left=2mm, right=2mm, top=10mm, bottom=8mm,
  overlay={%
    % cuadro-firma + kicker (arriba-izquierda)
    \fill[subjectcolor] ([shift={(4mm,-4mm)}]frame.north west) rectangle ++(2.6mm,-2.6mm);
    \node[anchor=north west, font=\InterDisplaySB\footnotesize, text=subjectcolor!85]
      at ([shift={(8.5mm,-3.4mm)}]frame.north west) {\seccionkicker};
    % número de serie grande y tenue (arriba-derecha) = marca
    \node[anchor=north east, font=\InterDisplayBlk, text=subjectcolor!18, inner sep=0pt]
      at ([shift={(-4mm,-1.5mm)}]frame.north east)
      {\fontsize{34}{34}\selectfont \numseccion};
    % franja de acento (abajo-izquierda)
    \fill[subjectcolor] ([shift={(4mm,4mm)}]frame.south west) rectangle ++(2.8cm,2.2mm);
    % sello del instituto (abajo-derecha)
    \node[anchor=south east, font=\sffamily\scriptsize, text=subjectcolor!55]
      at ([shift={(-4mm,3.4mm)}]frame.south east) {ITSTZ · Zapatoca};
  }}

\newcommand{\aperturaseccion}[3]{% {archivo}{título}{intro}
  \clearpage\refstepcounter{secportada}%
  \begin{secstage}
  \IfFileExists{assets/imgs/#1}{%
    \includegraphics[width=0.9\linewidth, height=3.2cm, keepaspectratio]{assets/imgs/#1}}{%
    \begin{tikzpicture}
      \node[draw=subjectcolor!35, dashed, line width=0.6pt, fill=white,
            minimum width=0.86\linewidth, minimum height=2.9cm,
            align=center, text=subjectcolor!50, font=\sffamily\small]
        {{\fontsize{22}{22}\selectfont\faImage}\\[5pt]
         \textbf{[ imagen pendiente ]}\\[1pt]
         \scriptsize\ttfamily assets/imgs/#1};
    \end{tikzpicture}}%
  \end{secstage}
  \vspace{0.5em}
  {\color{subjectcolor}\rule{3.2cm}{2.5pt}}\par\vspace{0.3em}
  \section{{\color{subjectcolor}\numseccion}\enspace #2}%
  {\InterLight\fontsize{12.5}{17}\selectfont\color{black!58} #3\par}
  \vspace{0.4em}}

% =====================================================================
%  DESTACADO — "lo esencial": triángulo lateral + degradado muy sutil
% =====================================================================
\newtcolorbox{destacado}[1]{
  enhanced, breakable, sharp corners, boxrule=0pt, frame hidden,
  interior style={left color=subjectcolor!16, right color=subjectcolor!2},
  borderline west={2.4pt}{0pt}{subjectcolor},
  left=7mm, right=5mm, top=3.5mm, bottom=3.5mm,
  before skip=\bloqueskip, after skip=\bloqueskip,
  overlay unbroken and first={%
    \fill[subjectcolor] ([yshift=-4.5mm]frame.north west)
        -- ++(5.5pt,2.75pt) -- ++(0,-5.5pt) -- cycle;},
  before upper={\blocklabel{subjectcolor}{Lo esencial}{#1}}
}

% --- Resaltado inline suave ---
\newtcbox{\resaltar}{on line, nobeforeafter, boxsep=0pt,
  left=2.5pt, right=2.5pt, top=1pt, bottom=1pt, arc=2.5pt,
  colback=subjectcolor!15, colframe=subjectcolor!15}

% =====================================================================
%  EXTENSIONES SANTILLANA (opt-in) — cajas temáticas para dar variedad
%  visual y romper la monotonía. NO alteran los entornos base; úsalas
%  donde aporten. Mantienen el color de asignatura (coherencia de marca).
% =====================================================================

% --- ¿SABÍAS QUE? — curiosidad breve · tarjeta redondeada con bombillo --
\newtcolorbox{sabias}{
  enhanced, breakable, arc=2.5mm, boxrule=0.8pt,
  colback=subjectcolor!6!white, colframe=subjectcolor!30,
  left=4mm, right=4mm, top=3mm, bottom=3mm,
  before skip=\bloqueskip, after skip=\bloqueskip,
  before upper={{\sffamily\bfseries\footnotesize\color{subjectcolor}%
    \faLightbulb\enspace\MakeUppercase{¿Sabías que?}}%
    \par\nopagebreak\vspace{0.45em}}
}

% --- CIENCIA Y SOCIEDAD — impacto/ética · banner de título -------------
\newtcolorbox{cienciasociedad}[1]{
  enhanced, breakable, sharp corners,
  colback=subjectcolor!5!white, colframe=subjectcolor, boxrule=0pt,
  colbacktitle=subjectcolor, coltitle=white,
  fonttitle=\sffamily\bfseries\footnotesize,
  title={\faGlobeAmericas\ \ \MakeUppercase{Ciencia y sociedad}%
         \ifblank{#1}{}{\ \ \textnormal{\sffamily\small· #1}}},
  toptitle=1.6mm, bottomtitle=1.6mm,
  left=4mm, right=4mm, top=3mm, bottom=3mm,
  before skip=\bloqueskip, after skip=\bloqueskip,
}

% --- MANOS A LA OBRA — práctica experimental · marco sólido con matraz --
\newtcolorbox{laboratorio}[1]{
  enhanced, breakable, sharp corners, arc=0pt,
  colback=white, colframe=subjectcolor!55, boxrule=1pt,
  left=4mm, right=4mm, top=3mm, bottom=3mm,
  before skip=\bloqueskip, after skip=\bloqueskip,
  before upper={{\sffamily\bfseries\footnotesize\color{subjectcolor}%
    \faFlask\enspace\MakeUppercase{Manos a la obra}}%
    \ifblank{#1}{}{\ {\color{black!35}\textbf{·}}\ %
      {\sffamily\bfseries\footnotesize\color{black!78}#1}}%
    \par\nopagebreak\vspace{0.45em}}
}
% Lista de materiales dentro de \begin{laboratorio}
\newcommand{\materiales}[1]{\par\noindent
  {\sffamily\bfseries\footnotesize\color{black!70}Materiales:}\ %
  {\footnotesize #1}\par\vspace{0.3em}}

% --- IMAGEN INLINE con texto alrededor (Santillana) -------------------
%   \figinline[r|l]{ancho}{archivo}{pie}   — colócala al inicio de un párrafo
\newcommand{\figinline}[4][r]{%
  \begin{wrapfigure}{#1}{#2}
  \centering\vspace{-0.5\baselineskip}
  \IfFileExists{assets/imgs/#3}{%
    \includegraphics[width=\linewidth]{assets/imgs/#3}}{%
    \fbox{\parbox[c][3cm][c]{0.92\linewidth}{\centering\sffamily\scriptsize
      \color{black!55}{\Large\faImage}\\[3pt][imagen pendiente]\\\ttfamily\tiny #3}}}%
  \par\vspace{2pt}{\footnotesize\sffamily\color{black!60}\textbf{Fig.}\enspace #4}%
  \vspace{-0.4\baselineskip}
  \end{wrapfigure}}

% --- APERTURA CON HERO LATERAL (imagen grande al lado del título) ------
%   \aperturahero{archivo}{título}{intro}  — alternativa a \aperturaseccion
\newcommand{\aperturahero}[3]{%
  \clearpage\refstepcounter{secportada}\stepcounter{section}%
  \addcontentsline{toc}{section}{\numseccion\enspace #2}%
  \noindent\begin{minipage}[c]{0.62\linewidth}
    {\InterDisplaySB\footnotesize\color{subjectcolor!85}\MakeUppercase{\seccionkicker}}\par
    \vspace{1.8mm}
    {\LARGE\InterDisplaySB\color{subjectcolor}%
      {\color{subjectcolor!35}\numseccion}\,\,#2}\par
    \vspace{2mm}{\color{subjectcolor!30}\rule{\linewidth}{0.8pt}}\par\vspace{2mm}
    {\InterLight\fontsize{12}{16}\selectfont\color{black!58}#3}
  \end{minipage}\hfill
  \begin{minipage}[c]{0.34\linewidth}\centering
    \IfFileExists{assets/imgs/#1}{%
      \includegraphics[width=\linewidth]{assets/imgs/#1}}{%
      \fbox{\parbox[c][3.4cm][c]{0.95\linewidth}{\centering\sffamily\scriptsize
        \color{black!55}{\Large\faImage}\\[3pt][imagen pendiente]\\\ttfamily\tiny #1}}}%
  \end{minipage}\par\vspace{1.3em}}

% --- MAPAS (conceptual y mental) — estilos TikZ con texto nítido -------
\tikzset{
  croot/.style={draw=subjectcolor, fill=subjectcolor, text=white, rounded corners=3pt,
    font=\sffamily\bfseries\footnotesize, align=center, inner sep=5pt, minimum height=9mm},
  cnode/.style={draw=subjectcolor!55, fill=subjectcolor!8, rounded corners=2.5pt,
    font=\sffamily\footnotesize, align=center, inner sep=4pt, text width=2.5cm},
  cleaf/.style={draw=subjectcolor!45, fill=white, rounded corners=2.5pt,
    font=\sffamily\footnotesize, align=center, inner sep=4pt},
  clink/.style={-{Stealth[length=5pt]}, subjectcolor!65, semithick},
  cline/.style={subjectcolor!55, semithick},
}
% Caja contenedora para un mapa (título + regla superior)
\newtcolorbox{mapabox}[2][Mapa conceptual]{
  enhanced, sharp corners, boxrule=0pt, frame hidden, colback=white,
  borderline north={1.4pt}{0pt}{subjectcolor},
  left=1mm, right=1mm, top=3mm, bottom=2mm,
  before skip=\bloqueskip, after skip=\bloqueskip,
  before upper={\blocklabel{subjectcolor}{#1}{#2}}
}

% --- Papel cuadriculado para graficar a mano (hojas de trabajo) --------
%   \papelcuadricula[filas]{ancho en cm}   (cada celda = 0.6 cm)
\newcommand{\papelcuadricula}[2][8]{%
  \par\smallskip\begin{center}
  \begin{tikzpicture}
    \draw[step=0.6cm, black!16, line width=0.3pt] (0,0) grid (#2,{#1*0.6});
    \draw[black!45, line width=0.7pt] (0,0) -- (#2,0);
    \draw[black!45, line width=0.7pt] (0,0) -- (0,{#1*0.6});
  \end{tikzpicture}
  \end{center}\par\smallskip}

% =====================================================================
%  APERTURA DE CAPÍTULO — título grande, número marca de agua, aire
% =====================================================================
\newcommand{\aperturacapitulo}[3]{% {número}{título}{subtítulo}
  \setcounter{unidadnum}{#1}\setcounter{secportada}{0}% fija unidad y reinicia secciones
  \thispagestyle{empty}%
  \begin{tikzpicture}[remember picture, overlay]
    \node[anchor=north east, text=subjectcolor!9,
          font=\InterDisplayThin]
      at ([xshift=-0.9cm, yshift=-2.2cm]current page.north east)
      {\fontsize{210}{210}\selectfont #1};
  \end{tikzpicture}%
  \begingroup\raggedright\null\vspace{3.8cm}\par
  {\InterDisplaySB\fontsize{12}{14}\selectfont\color{subjectcolor}%
    CAP\'ITULO\ #1}\par
  \vspace{0.18cm}{\color{black!16}\rule{4cm}{1pt}}\par
  \vspace{0.9cm}
  {\InterDisplaySB\fontsize{46}{50}\selectfont\color{black!88}#2}\par
  \vspace{0.5cm}{\color{subjectcolor}\rule{3.2cm}{2.5pt}}\par
  \vspace{0.55cm}
  {\InterLight\fontsize{15}{21}\selectfont\color{black!55}#3}\par
  \vfill
  {\sffamily\footnotesize\color{black!45}\hdrinstituto\quad
    \textcolor{black!22}{·}\quad Prof.\ Juan Diego A.\ Prada R.}\par
  \endgroup\clearpage}

% =====================================================================
%  DESBLOQUEA — pie discreto, no caja
% =====================================================================
\newcommand{\desbloquea}[1]{%
  \par\addvspace{\bloqueskip}%
  \noindent{\color{subjectcolor!25}\rule{\linewidth}{0.5pt}}\par\vspace{0.4em}%
  {\footnotesize\sffamily\color{black!60}%
    \textcolor{subjectcolor}{\faUnlock}\ \textbf{Dominar esto te desbloquea:}\ #1}\par}
):
%     % =====================================================================
%  TEMA · CIENCIAS NATURALES — variante del design system (GUÍAS)
%  Identidad visual propia, distinta de Matemáticas (azul/pizarra):
%    · color base ESMERALDA (\setsubject{cnat})
%    · acento LIMA orgánico solo para ornamentos (no rompe la regla de
%      "un solo color de asignatura")
%    · iconografía de ciencias (hoja · ADN · matraz · átomo)
%    · reutiliza las cajas de ciencia que ya trae el preamble:
%      \begin{laboratorio}, \begin{sabias}, \begin{cienciasociedad}
%
%  Sirve a futuro para BIOLOGÍA, QUÍMICA y FÍSICA EXPERIMENTAL.
%
%  INVOCACIÓN (una sola línea, tras \input{design/preamble.tex}):
%     \input{design/tema-cnat.tex}
%  El doc puede seguir personalizando \hdrcapitulo, \seccionkicker, etc.
% =====================================================================

% --- Color base de la asignatura --------------------------------------
\setsubject{cnat}                 % esmeralda RGB(16,110,80)

% Acento LIMA orgánico — SOLO ornamentos (hoja, filete). Nunca contenido:
% la regla "un solo color de asignatura" se mantiene intacta.
\definecolor{cnatlima}{RGB}{124,169,60}

% --- Iconografía semántica de ciencias (fontawesome5, ya cargado) -----
\newcommand{\icobio}{{\color{subjectcolor}\faLeaf}}     % biología
\newcommand{\icoquim}{{\color{subjectcolor}\faFlask}}   % química
\newcommand{\icofis}{{\color{subjectcolor}\faAtom}}     % física
\newcommand{\icogen}{{\color{subjectcolor}\faDna}}      % genética

% --- Kicker temático por defecto (el documento puede renovarlo) -------
\renewcommand{\seccionkicker}{\faLeaf\enspace Ciencias Naturales}

% --- Nombre científico: género + especie en itálica + nombre común ----
%   \especie{Pisum}{sativum}{arveja}  ->  *Pisum sativum* (arveja)
\newcommand{\especie}[3]{\textit{#1 #2}\ (#3)}

% --- Genotipo en línea: preserva mayúsculas/minúsculas de los alelos ---
%   \genotipo{Aa}  (NUNCA aplicar \MakeUppercase: borraría dominante/recesivo)
\newcommand{\genotipo}[1]{{\sffamily\bfseries\color{subjectcolor}#1}}

% --- Filete orgánico con hoja: cierre decorativo opcional de bloque ----
\newcommand{\filetehoja}{%
  \par\vspace{0.4em}\noindent
  {\color{cnatlima}\faLeaf}\hspace{0.3em}%
  {\color{subjectcolor!35}\rule[0.35ex]{0.86\linewidth}{0.7pt}}\par}

%  El doc puede seguir personalizando \hdrcapitulo, \seccionkicker, etc.
% =====================================================================

% --- Color base de la asignatura --------------------------------------
\setsubject{cnat}                 % esmeralda RGB(16,110,80)

% Acento LIMA orgánico — SOLO ornamentos (hoja, filete). Nunca contenido:
% la regla "un solo color de asignatura" se mantiene intacta.
\definecolor{cnatlima}{RGB}{124,169,60}

% --- Iconografía semántica de ciencias (fontawesome5, ya cargado) -----
\newcommand{\icobio}{{\color{subjectcolor}\faLeaf}}     % biología
\newcommand{\icoquim}{{\color{subjectcolor}\faFlask}}   % química
\newcommand{\icofis}{{\color{subjectcolor}\faAtom}}     % física
\newcommand{\icogen}{{\color{subjectcolor}\faDna}}      % genética

% --- Kicker temático por defecto (el documento puede renovarlo) -------
\renewcommand{\seccionkicker}{\faLeaf\enspace Ciencias Naturales}

% --- Nombre científico: género + especie en itálica + nombre común ----
%   \especie{Pisum}{sativum}{arveja}  ->  *Pisum sativum* (arveja)
\newcommand{\especie}[3]{\textit{#1 #2}\ (#3)}

% --- Genotipo en línea: preserva mayúsculas/minúsculas de los alelos ---
%   \genotipo{Aa}  (NUNCA aplicar \MakeUppercase: borraría dominante/recesivo)
\newcommand{\genotipo}[1]{{\sffamily\bfseries\color{subjectcolor}#1}}

% --- Filete orgánico con hoja: cierre decorativo opcional de bloque ----
\newcommand{\filetehoja}{%
  \par\vspace{0.4em}\noindent
  {\color{cnatlima}\faLeaf}\hspace{0.3em}%
  {\color{subjectcolor!35}\rule[0.35ex]{0.86\linewidth}{0.7pt}}\par}

%  El doc puede seguir personalizando \hdrcapitulo, \seccionkicker, etc.
% =====================================================================

% --- Color base de la asignatura --------------------------------------
\setsubject{cnat}                 % esmeralda RGB(16,110,80)

% Acento LIMA orgánico — SOLO ornamentos (hoja, filete). Nunca contenido:
% la regla "un solo color de asignatura" se mantiene intacta.
\definecolor{cnatlima}{RGB}{124,169,60}

% --- Iconografía semántica de ciencias (fontawesome5, ya cargado) -----
\newcommand{\icobio}{{\color{subjectcolor}\faLeaf}}     % biología
\newcommand{\icoquim}{{\color{subjectcolor}\faFlask}}   % química
\newcommand{\icofis}{{\color{subjectcolor}\faAtom}}     % física
\newcommand{\icogen}{{\color{subjectcolor}\faDna}}      % genética

% --- Kicker temático por defecto (el documento puede renovarlo) -------
\renewcommand{\seccionkicker}{\faLeaf\enspace Ciencias Naturales}

% --- Nombre científico: género + especie en itálica + nombre común ----
%   \especie{Pisum}{sativum}{arveja}  ->  *Pisum sativum* (arveja)
\newcommand{\especie}[3]{\textit{#1 #2}\ (#3)}

% --- Genotipo en línea: preserva mayúsculas/minúsculas de los alelos ---
%   \genotipo{Aa}  (NUNCA aplicar \MakeUppercase: borraría dominante/recesivo)
\newcommand{\genotipo}[1]{{\sffamily\bfseries\color{subjectcolor}#1}}

% --- Filete orgánico con hoja: cierre decorativo opcional de bloque ----
\newcommand{\filetehoja}{%
  \par\vspace{0.4em}\noindent
  {\color{cnatlima}\faLeaf}\hspace{0.3em}%
  {\color{subjectcolor!35}\rule[0.35ex]{0.86\linewidth}{0.7pt}}\par}
     % si es ciencias (verde); mate va en azul
%     % =====================================================================
%  TEMA · MODO TABLERO — "clave del profe" a TODO EL ANCHO
%  Rompe a propósito el layout de libro (twoside angosto): página ancha
%  de una sola columna, banner moderno, imágenes libres al lado con el
%  texto rodeándolas (\figinline), gráficos flat TikZ/SVG y TODO RESUELTO.
%
%  Pensado para material que el PROFE resuelve en el tablero (taller/quiz
%  con solución). NO es el estilo Santillana del kit: es intencionalmente
%  más "revista/póster", más aire, más ancho.
%
%  INVOCACIÓN (tras % =====================================================================
%  DESIGN SYSTEM v5 — EDITORIAL ITSTZ
%  Filosofía: texto primero · un solo color de asignatura · sin sombras
%  · sin marcos pesados · tratamientos planos (filete lateral / regla).
%
%  Vocabulario semántico (esto es TODO lo disponible):
%    \section / \subsection      jerarquía
%    definicion{título}          concepto/regla clave (filete color)
%    ejemplo{título}             ejemplo resuelto (numerado, sin caja)
%    procedimiento{título}       pasos (filete neutro) + \paso{n}{texto}
%    nota[título]                aparte único gris (sabías/dato/video/margen)
%    alerta{título}              advertencia ámbar (usar rara vez)
%    practica{título}            ejercicios/hoja/ICFES (regla superior)
%    \figura{archivo}{pie}       imagen real o placeholder automático
%    \desbloquea{lista}          pie discreto de prerrequisitos
%
%  REGLA DE ORO: las secciones (sNN.tex) NO llevan \vspace, color ni
%  minipage de maquetación. Solo contenido semántico. El diseño vive aquí.
% =====================================================================
\documentclass[11pt, letterpaper, twoside]{article}

\usepackage{fontspec}
\usepackage{polyglossia}
\setdefaultlanguage{spanish}
\usepackage[inner=2.1cm, outer=3.7cm, top=2.6cm, bottom=2.3cm,
            marginparwidth=2.55cm, marginparsep=0.45cm,
            headheight=16pt, headsep=18pt, footskip=24pt]{geometry}
\usepackage{microtype}
\usepackage{multicol}
\usepackage{xcolor}
\usepackage{graphicx}
\usepackage{tikz}
\usetikzlibrary{calc, arrows.meta, positioning}
\usepackage{wrapfig}
\usepackage[most]{tcolorbox}
\usepackage{amsmath, amssymb}
\usepackage{siunitx}
\usepackage{fontawesome5}
\usepackage{fancyhdr}
\usepackage{titlesec}
\usepackage{enumitem}
\usepackage{etoolbox}
\usepackage{titletoc}
\usepackage{array}
\usepackage{booktabs}
\usepackage{colortbl}
\usepackage{nicematrix}   % rejilla completa (hvlines) para tablas de hojas

% --- TIPOGRAFÍA -------------------------------------------------------
% Cuerpo: Charter (serif cálida y muy legible). La matemática (Latin
% Modern, serif) combina con ella. Títulos/etiquetas/UI: Inter Display.
\setmainfont{Charter}[
  UprightFont    = {Charter},
  BoldFont       = {Charter Bold},
  ItalicFont     = {Charter Italic},
  BoldItalicFont = {Charter Bold Italic},
  Ligatures = TeX, Numbers = Lining ]
\setsansfont{Inter}[
  UprightFont = Inter-Regular, BoldFont = Inter-SemiBold,
  ItalicFont  = Inter-Italic, Ligatures = TeX, Numbers = Lining ]
\newfontfamily\InterLight{Inter-Light}[ItalicFont=Inter-LightItalic, Ligatures=TeX]
\newfontfamily\InterDisplay{Inter Display}[Ligatures=TeX]
\newfontfamily\InterDisplaySB{Inter Display SemiBold}[Ligatures=TeX]
\newfontfamily\InterDisplayThin{Inter Display Thin}[Ligatures=TeX]
\newfontfamily\InterDisplayBlk{Inter Display Black}[Ligatures=TeX]

% --- RITMO --------------------------------------------------------------
\linespread{1.22}
\setlength{\parindent}{0pt}
\setlength{\parskip}{0.5em}
\newlength{\bloqueskip}\setlength{\bloqueskip}{1.0em}   % aire entre bloques
\setlength{\columnsep}{1.6em}                            % medianil multicols

% =====================================================================
%  COLOR — UNA sola variable de asignatura + neutros. Sin arcoíris.
% =====================================================================
\definecolor{mat} {RGB}{ 30,  64, 120}   % azul profundo
\definecolor{cnat}{RGB}{ 16, 110,  80}   % esmeralda
\definecolor{csoc}{RGB}{165,  75,  45}   % terracota
\definecolor{eco} {RGB}{ 70,  55, 130}   % índigo
\definecolor{alertcolor}{RGB}{180, 95, 10}  % ámbar — solo advertencias

\colorlet{subjectcolor}{mat}             % por defecto
\newcommand{\setsubject}[1]{\colorlet{subjectcolor}{#1}}

% =====================================================================
%  JERARQUÍA TIPOGRÁFICA
% =====================================================================
\setcounter{secnumdepth}{0}
\titleformat{\section}
  {\LARGE\InterDisplaySB\color{subjectcolor}}{}{0pt}{}
  [\vspace{2pt}{\color{subjectcolor!30}\titlerule[0.8pt]}]
\titlespacing*{\section}{0pt}{3.2ex plus 0.6ex}{1.8ex plus 0.3ex}

\titleformat{\subsection}
  {\large\InterDisplaySB\color{black!82}}{}{0pt}{}
\titlespacing*{\subsection}{0pt}{2.0ex plus 0.4ex}{0.5ex}

% --- TABLA DE CONTENIDO -------------------------------------------------
\titlecontents{section}[0em]
  {\addvspace{0.45pc}\sffamily\bfseries\small\color{black!85}}
  {}{}{\titlerule*[0.35pc]{.}\contentspage}
\titlecontents{subsection}[1.4em]
  {\sffamily\footnotesize\color{black!60}}
  {}{}{\titlerule*[0.35pc]{.}\contentspage}

% =====================================================================
%  ENCABEZADO / PIE — editorial limpio (regla fina, sin franjas)
% =====================================================================
\newcommand{\hdrinstituto}{Instituto Técnico Santo Tomás · Zapatoca}
\newcommand{\hdrcapitulo}{}            % se fija en el documento maestro
\pagestyle{fancy}\fancyhf{}
\renewcommand{\headrulewidth}{0pt}
\fancyhead[LE]{\footnotesize\sffamily\color{black!55}\hdrinstituto}
\fancyhead[RO]{\footnotesize\sffamily\color{subjectcolor}\hdrcapitulo}
\fancyhead[LO,RE]{}
\renewcommand{\headrule}{\vspace{2pt}{\color{black!18}\hrule height 0.4pt}}

% --- PIE: número exterior (cuadro de color) + crédito interior + filete ---
\renewcommand{\footrulewidth}{0pt}
\renewcommand{\footrule}{{\color{black!15}\hrule height 0.4pt}\vspace{3pt}}
\newcommand{\pagenumbox}{%
  \tikz[baseline=-3.2pt]{\node[fill=subjectcolor, text=white, inner xsep=7pt,
     inner ysep=3pt, font=\sffamily\bfseries\footnotesize]{\thepage};}}
\fancyfoot[RO,LE]{\pagenumbox}
\fancyfoot[LO,RE]{\footnotesize\sffamily\color{black!45}%
  \hdrcapitulo\ \textcolor{black!22}{·}\ Prof.\ Juan Diego A.\ Prada R.}

% =====================================================================
%  HELPERS
% =====================================================================
% Etiqueta de bloque:  TAG · título   (color del tag configurable)
\newcommand{\blocklabel}[3]{% #1=color  #2=TAG  #3=título
  {\sffamily\bfseries\footnotesize\color{#1}\MakeUppercase{#2}}%
  \ifblank{#3}{}{\,{\color{black!35}\textbf{·}}\,%
    {\sffamily\bfseries\footnotesize\color{black!78}#3}}%
  \par\nopagebreak\vspace{0.4em}}

% --- CAJAS AL MARGEN — personajes históricos e hitos -------------------
% Van en el margen exterior ancho (\marginpar respeta la geometría twoside
% y flota para no chocar). Son contenido semántico: colócalas en el cuerpo,
% cerca del párrafo que comentan, ENTRE párrafos. NO usarlas dentro de
% multicols, tcolorbox, tablas, ni en la misma línea que un \section.
%   \personaje{Nombre}{años · rol}{descripción breve}
%   \hito{año}{qué pasó (breve)}
\setlength{\marginparpush}{8pt}
\newcommand{\margenbox}[1]{\marginpar{%
  \begingroup\raggedright\footnotesize\sffamily
  \setlength{\parskip}{0.18em}\linespread{1.05}\selectfont
  #1\par\endgroup}}
\newcommand{\personaje}[3]{\margenbox{%
  {\color{subjectcolor}\rule{\linewidth}{1.1pt}}\par
  {\InterDisplaySB\scriptsize\color{subjectcolor}#1}\par
  {\fontsize{6.8pt}{8pt}\selectfont\itshape\color{black!50}#2}\par
  \vspace{0.15em}{\scriptsize\color{black!72}#3}}}
\newcommand{\hito}[2]{\margenbox{%
  {\sffamily\bfseries\scriptsize\colorbox{subjectcolor}{\textcolor{white}{\,#1\,}}}\par
  \vspace{0.15em}{\scriptsize\color{black!72}#2}}}

% =====================================================================
%  ENTORNOS — todos planos: filete lateral, sin marco, sin sombra
% =====================================================================

% --- DEFINICIÓN / concepto clave ---------------------------------------
\newtcolorbox{definicion}[1]{
  enhanced, breakable, sharp corners, boxrule=0pt, frame hidden,
  colback=subjectcolor!6!white,
  borderline west={2.2pt}{0pt}{subjectcolor},
  left=5mm, right=4mm, top=3mm, bottom=3mm,
  before skip=\bloqueskip, after skip=\bloqueskip,
  before upper={\blocklabel{subjectcolor}{Definición}{#1}}
}

% --- PROCEDIMIENTO / pasos ---------------------------------------------
\newtcolorbox{procedimiento}[1]{
  enhanced, breakable, sharp corners, boxrule=0pt, frame hidden,
  colback=black!3!white,
  borderline west={2.2pt}{0pt}{black!45},
  left=5mm, right=4mm, top=3mm, bottom=3mm,
  before skip=\bloqueskip, after skip=\bloqueskip,
  before upper={\blocklabel{black!60}{Procedimiento}{#1}}
}
\newcommand{\paso}[2]{%
  \par\vspace{0.45em}\noindent
  \tikz[baseline=-0.6ex]\node[circle, fill=subjectcolor, text=white,
     inner sep=0pt, minimum size=1.45em,
     font=\sffamily\bfseries\footnotesize]{#1};\hspace{0.5em}#2\par}

% Glosa "¿por qué?" bajo un paso — narración en lenguaje cercano
\newcommand{\porque}[1]{\par\vspace{0.1em}\noindent\hspace{2.0em}%
  {\sffamily\footnotesize\itshape\color{subjectcolor}¿por qué?\,}%
  {\sffamily\footnotesize\itshape\color{black!55}#1}\par}

% Operación glosada: cuenta a la izquierda + porqué al margen derecho
% (estilo "comentario al lado del código"). Úsala dentro de ejemplo.
% \opg{ecuación}{porqué}   — primer arg en modo display (mat, cap0)
% \opgt{etiqueta}{porqué}  — primer arg en texto sans bold (física, ciencias)
\newcommand{\opg}[2]{\par\vspace{0.25em}\noindent
  \makebox[0.46\linewidth][l]{$\displaystyle #1$}%
  {\sffamily\footnotesize\itshape\color{black!52}#2}\par}
\newcommand{\opgt}[2]{\par\vspace{0.25em}\noindent
  \makebox[0.46\linewidth][l]{\sffamily\bfseries\color{subjectcolor!80}#1}%
  {\sffamily\footnotesize\itshape\color{black!52}#2}\par}

% --- NOTA / aparte único (sabías · dato · video · margen) --------------
\newtcolorbox{nota}[1][]{
  enhanced, breakable, sharp corners, boxrule=0pt, frame hidden,
  colback=black!4!white,
  borderline west={2.2pt}{0pt}{black!35},
  left=5mm, right=4mm, top=3mm, bottom=3mm,
  before skip=\bloqueskip, after skip=\bloqueskip,
  before upper={\ifblank{#1}{}{\blocklabel{black!55}{Nota}{#1}}}
}

% --- ALERTA / advertencia (única en ámbar, uso escaso) -----------------
\newtcolorbox{alerta}[1]{
  enhanced, breakable, sharp corners, boxrule=0pt, frame hidden,
  colback=alertcolor!7!white,
  borderline west={2.2pt}{0pt}{alertcolor},
  left=5mm, right=4mm, top=3mm, bottom=3mm,
  before skip=\bloqueskip, after skip=\bloqueskip,
  before upper={\blocklabel{alertcolor}{Atención}{#1}}
}
% Sub-ítem de alerta (antes \caso)
\newcommand{\caso}[2]{\par\vspace{0.35em}\noindent
  {\sffamily\bfseries\color{alertcolor!85!black}#1.}\enspace #2\par}

% --- CONTRASTE (la virtud y su exceso → pregunta que cuestiona) --------
%   \contraste{Concepto}{La virtud / lo valioso}{Su exceso / su riesgo}{Pregunta}
\newtcolorbox{contrastebox}{
  enhanced, breakable, sharp corners, boxrule=0pt, frame hidden,
  colback=subjectcolor!5!white,
  borderline west={2.2pt}{0pt}{subjectcolor},
  left=5mm, right=4mm, top=3mm, bottom=3mm,
  before skip=\bloqueskip, after skip=\bloqueskip,
}
\newcommand{\contraste}[4]{%
\begin{contrastebox}
\blocklabel{subjectcolor}{Cuestiona esto}{\faBalanceScale\ #1}
\noindent\begin{minipage}[t]{0.45\linewidth}
  {\sffamily\footnotesize\bfseries\color{subjectcolor!85!black}La virtud}\par
  \vspace{0.2em}#2
\end{minipage}\hfill
{\color{subjectcolor!35}\vrule width 0.4pt}\hfill
\begin{minipage}[t]{0.45\linewidth}
  {\sffamily\footnotesize\bfseries\color{alertcolor!85!black}Su exceso}\par
  \vspace{0.2em}#3
\end{minipage}
\par\vspace{0.6em}
{\color{subjectcolor!25}\titlerule[0.6pt]}\par\vspace{0.4em}
\noindent{\sffamily\itshape\color{subjectcolor}\faArrowRight\ #4}
\end{contrastebox}}

% --- EJEMPLO (sin caja: regla fina + etiqueta numerada) ----------------
\newcounter{ejemplo}[section]
\newenvironment{ejemplo}[1]
  {\par\addvspace{\bloqueskip}\refstepcounter{ejemplo}%
   \noindent{\color{subjectcolor!28}\rule{\linewidth}{0.5pt}}\par\vspace{0.4em}%
   \noindent{\sffamily\bfseries\color{subjectcolor}Ejemplo \theejemplo.}%
   \enspace{\sffamily\bfseries\color{black!78}#1}\par\vspace{0.3em}}
  {\par\addvspace{\bloqueskip}}

% --- PRÁCTICA (ejercicios / hoja / ICFES) ------------------------------
\newtcolorbox{practica}[1]{
  enhanced, breakable, sharp corners, boxrule=0pt, frame hidden,
  colback=white,
  borderline north={1.4pt}{0pt}{subjectcolor},
  left=0mm, right=0mm, top=3mm, bottom=2mm,
  before skip=\bloqueskip, after skip=\bloqueskip,
  before upper={\blocklabel{subjectcolor}{Práctica}{#1}\raggedright}
}

% Ejercicio numerado con nivel opcional [básico|intermedio|reto]
\newcounter{ejer}[section]
\newcommand{\ejer}[2][]{%
  \stepcounter{ejer}\par\vspace{0.5em}\noindent
  {\sffamily\bfseries\color{subjectcolor}\theejer.}\enspace #2%
  \ifblank{#1}{}{\ \textcolor{black!45}{\footnotesize[\textit{#1}]}}\par}

% Niveles = un solo color a tres intensidades (NO nuevos colores)
\newcommand{\nivel}[2]{\par\smallskip\noindent
  \textcolor{subjectcolor!#1}{\small\faCircle}\enspace
  {\sffamily\bfseries\footnotesize\color{black!70}#2}\par}
\newcommand{\basico}{\nivel{40}{Básico}}
\newcommand{\intermedio}{\nivel{70}{Intermedio}}
\newcommand{\reto}{\nivel{100}{Reto}}

% Selección múltiple tipo ICFES — opciones con burbuja (estilo Saber)
\newcounter{pregicfes}[section]
\newcommand{\bubbleletter}[1]{\tikz[baseline=(bc.base)]\node[circle,
  draw=subjectcolor!55, line width=0.6pt, inner sep=1.5pt,
  font=\sffamily\bfseries\scriptsize, text=subjectcolor!90!black](bc){#1};}
\newcommand{\icfes}[5]{%
  \stepcounter{pregicfes}\par\vspace{0.6em}\noindent
  {\sffamily\bfseries\color{subjectcolor}\thepregicfes.}\enspace #1\par
  \begin{enumerate}[label={\protect\bubbleletter{\Alph*}}, leftmargin=1.5cm,
                    itemsep=3pt, topsep=3pt]
    \item #2 \item #3 \item #4 \item #5
  \end{enumerate}}
% Pregunta abierta de argumentación
\newcommand{\pregabierta}[1]{\par\vspace{0.6em}\noindent
  {\sffamily\bfseries\color{subjectcolor}\faComment\ Argumenta:}\enspace
  \textit{#1}\par\lineaescritura\lineaescritura\lineaescritura}

% --- LÍNEAS DE RESPUESTA -----------------------------------------------
\newcommand{\linea}[1][7cm]{\underline{\hspace{#1}}}
\newcommand{\lineaescritura}{\par\vspace{0.9em}\noindent
  {\color{black!30}\rule{\linewidth}{0.3pt}}}

% =====================================================================
%  EVALUACIÓN TIPO ICFES — cierre de unidad (banner + 2 columnas)
%  Uso:  \begin{evaluacion} \begin{multicols}{2} \icfes... \end{multicols}
%        \pregabierta{...} \end{evaluacion}
% =====================================================================
\newtcolorbox{evaluacion}[1][cierre de sección]{
  enhanced, breakable, sharp corners,
  colback=white, colframe=subjectcolor, boxrule=1pt,
  colbacktitle=subjectcolor, coltitle=white,
  fonttitle=\InterDisplaySB\normalsize,
  title={\faClipboardCheck\ \ EVALUACIÓN TIPO ICFES\ \ %
         \textnormal{\sffamily\small· #1}},
  toptitle=2mm, bottomtitle=2mm, left=4mm, right=4mm, top=3.5mm, bottom=3mm,
  before={\clearpage},
  before upper={\small\itshape\color{black!55}%
    Selección múltiple. Marca una sola opción por pregunta.\par\medskip\raggedright}
}

% =====================================================================
%  HOJA DE TRABAJO — recortable: borde punteado, escala de grises,
%  en página aparte (clearpage antes; el contenido que sigue continúa).
% =====================================================================
\newtcolorbox{hoja}[1]{
  enhanced, breakable, sharp corners,
  colback=black!2!white, boxrule=0pt, frame hidden,
  borderline={0.8pt}{0pt}{black!50, dash pattern=on 3pt off 2.5pt},
  before={\clearpage}, after={\par},
  left=6mm, right=6mm, top=6mm, bottom=6mm,
  before upper={%
    {\sffamily\bfseries\footnotesize\color{black!60}\faCut\ \ HOJA DE TRABAJO}%
    {\color{black!35}\ \textbf{·}\ }{\sffamily\bfseries\footnotesize\color{black!80}#1}%
    \hfill{\sffamily\footnotesize\color{black!55}%
      Nombre:\ \underline{\hspace{4.2cm}}\quad Fecha:\ \underline{\hspace{2cm}}}\par
    \vspace{2mm}{\color{black!30}\hrule height 0.4pt}\vspace{3mm}\raggedright}
}

% =====================================================================
%  TABLAHOJA — tabla con TODAS las líneas (rejilla) para hojas de trabajo.
%  Celdas altas para llenar a mano. Centrada. El encabezado se marca con
%  \rowcolor en la primera fila.
%  Uso:  \begin{tablahoja}{ccc} enc & enc & enc \\ a & b & c \\ \end{tablahoja}
% =====================================================================
\NewDocumentEnvironment{tablahoja}{m}{%
  \par\smallskip\centering\small\setlength{\tabcolsep}{4pt}%
  \begin{NiceTabular}{#1}[hvlines, cell-space-top-limit=5pt,
     cell-space-bottom-limit=9pt, colortbl-like]%
}{%
  \end{NiceTabular}\par\smallskip
}

% =====================================================================
%  FIGURAS — imagen real o placeholder automático (\IfFileExists)
% =====================================================================
\newcounter{figura}[section]
\newcommand{\figura}[3][0.7\linewidth]{% [ancho]{archivo}{pie}
  \par\addvspace{\bloqueskip}\refstepcounter{figura}%
  \begin{center}%
  \IfFileExists{assets/imgs/#2}{%
    \includegraphics[width=#1]{assets/imgs/#2}}{%
    \begin{tikzpicture}
      \node[draw=black!30, dashed, line width=0.6pt, fill=black!3,
            text width=0.66\linewidth, minimum height=2.6cm,
            align=center, inner sep=8pt,
            text=black!55, font=\sffamily\footnotesize]
        {{\Large\faImage}\\[4pt]\textbf{[ imagen pendiente ]}\\#3};
    \end{tikzpicture}}%
  \par\vspace{0.4em}%
  {\footnotesize\sffamily\color{black!60}%
    \textbf{Figura \thefigura.}\enspace #3}%
  \end{center}\par\addvspace{\bloqueskip}}

% =====================================================================
%  APERTURA DE SECCIÓN — marco branded de serie + título + "qué verás"
%  El "marco" es CONSTANTE en todas las secciones (= identidad de marca):
%  número grande de serie, kicker, cuadro-firma, franja de acento y sello
%  del instituto. El arte (imagen real o placeholder) va DENTRO del marco.
%  Uso (1ª línea de cada sNN.tex):
%    \aperturaseccion{ruta/img.png}{Título}{Una o dos frases.}
%  Personaliza por documento (opcional):  \renewcommand{\seccionkicker}{...}
% =====================================================================
\newcounter{secportada}
\newcounter{unidadnum}                       % nº de unidad actual (lo fija \aperturacapitulo)
% Número visible de sección = Unidad.Sección (ej. 1.2). Si no hay unidad
% definida (unidadnum=0), cae al número simple de sección.
\newcommand{\numseccion}{\ifnum\value{unidadnum}>0 \theunidadnum.\thesecportada\else\thesecportada\fi}
\newcommand{\seccionkicker}{\hdrcapitulo}   % por defecto, lo del encabezado

% El escenario branded (contenedor constante del arte)
\newtcolorbox{secstage}{%
  enhanced, sharp corners, boxrule=0pt, frame hidden,
  colback=subjectcolor!6, halign=center, valign=center,
  height=5.0cm, left=2mm, right=2mm, top=10mm, bottom=8mm,
  overlay={%
    % cuadro-firma + kicker (arriba-izquierda)
    \fill[subjectcolor] ([shift={(4mm,-4mm)}]frame.north west) rectangle ++(2.6mm,-2.6mm);
    \node[anchor=north west, font=\InterDisplaySB\footnotesize, text=subjectcolor!85]
      at ([shift={(8.5mm,-3.4mm)}]frame.north west) {\seccionkicker};
    % número de serie grande y tenue (arriba-derecha) = marca
    \node[anchor=north east, font=\InterDisplayBlk, text=subjectcolor!18, inner sep=0pt]
      at ([shift={(-4mm,-1.5mm)}]frame.north east)
      {\fontsize{34}{34}\selectfont \numseccion};
    % franja de acento (abajo-izquierda)
    \fill[subjectcolor] ([shift={(4mm,4mm)}]frame.south west) rectangle ++(2.8cm,2.2mm);
    % sello del instituto (abajo-derecha)
    \node[anchor=south east, font=\sffamily\scriptsize, text=subjectcolor!55]
      at ([shift={(-4mm,3.4mm)}]frame.south east) {ITSTZ · Zapatoca};
  }}

\newcommand{\aperturaseccion}[3]{% {archivo}{título}{intro}
  \clearpage\refstepcounter{secportada}%
  \begin{secstage}
  \IfFileExists{assets/imgs/#1}{%
    \includegraphics[width=0.9\linewidth, height=3.2cm, keepaspectratio]{assets/imgs/#1}}{%
    \begin{tikzpicture}
      \node[draw=subjectcolor!35, dashed, line width=0.6pt, fill=white,
            minimum width=0.86\linewidth, minimum height=2.9cm,
            align=center, text=subjectcolor!50, font=\sffamily\small]
        {{\fontsize{22}{22}\selectfont\faImage}\\[5pt]
         \textbf{[ imagen pendiente ]}\\[1pt]
         \scriptsize\ttfamily assets/imgs/#1};
    \end{tikzpicture}}%
  \end{secstage}
  \vspace{0.5em}
  {\color{subjectcolor}\rule{3.2cm}{2.5pt}}\par\vspace{0.3em}
  \section{{\color{subjectcolor}\numseccion}\enspace #2}%
  {\InterLight\fontsize{12.5}{17}\selectfont\color{black!58} #3\par}
  \vspace{0.4em}}

% =====================================================================
%  DESTACADO — "lo esencial": triángulo lateral + degradado muy sutil
% =====================================================================
\newtcolorbox{destacado}[1]{
  enhanced, breakable, sharp corners, boxrule=0pt, frame hidden,
  interior style={left color=subjectcolor!16, right color=subjectcolor!2},
  borderline west={2.4pt}{0pt}{subjectcolor},
  left=7mm, right=5mm, top=3.5mm, bottom=3.5mm,
  before skip=\bloqueskip, after skip=\bloqueskip,
  overlay unbroken and first={%
    \fill[subjectcolor] ([yshift=-4.5mm]frame.north west)
        -- ++(5.5pt,2.75pt) -- ++(0,-5.5pt) -- cycle;},
  before upper={\blocklabel{subjectcolor}{Lo esencial}{#1}}
}

% --- Resaltado inline suave ---
\newtcbox{\resaltar}{on line, nobeforeafter, boxsep=0pt,
  left=2.5pt, right=2.5pt, top=1pt, bottom=1pt, arc=2.5pt,
  colback=subjectcolor!15, colframe=subjectcolor!15}

% =====================================================================
%  EXTENSIONES SANTILLANA (opt-in) — cajas temáticas para dar variedad
%  visual y romper la monotonía. NO alteran los entornos base; úsalas
%  donde aporten. Mantienen el color de asignatura (coherencia de marca).
% =====================================================================

% --- ¿SABÍAS QUE? — curiosidad breve · tarjeta redondeada con bombillo --
\newtcolorbox{sabias}{
  enhanced, breakable, arc=2.5mm, boxrule=0.8pt,
  colback=subjectcolor!6!white, colframe=subjectcolor!30,
  left=4mm, right=4mm, top=3mm, bottom=3mm,
  before skip=\bloqueskip, after skip=\bloqueskip,
  before upper={{\sffamily\bfseries\footnotesize\color{subjectcolor}%
    \faLightbulb\enspace\MakeUppercase{¿Sabías que?}}%
    \par\nopagebreak\vspace{0.45em}}
}

% --- CIENCIA Y SOCIEDAD — impacto/ética · banner de título -------------
\newtcolorbox{cienciasociedad}[1]{
  enhanced, breakable, sharp corners,
  colback=subjectcolor!5!white, colframe=subjectcolor, boxrule=0pt,
  colbacktitle=subjectcolor, coltitle=white,
  fonttitle=\sffamily\bfseries\footnotesize,
  title={\faGlobeAmericas\ \ \MakeUppercase{Ciencia y sociedad}%
         \ifblank{#1}{}{\ \ \textnormal{\sffamily\small· #1}}},
  toptitle=1.6mm, bottomtitle=1.6mm,
  left=4mm, right=4mm, top=3mm, bottom=3mm,
  before skip=\bloqueskip, after skip=\bloqueskip,
}

% --- MANOS A LA OBRA — práctica experimental · marco sólido con matraz --
\newtcolorbox{laboratorio}[1]{
  enhanced, breakable, sharp corners, arc=0pt,
  colback=white, colframe=subjectcolor!55, boxrule=1pt,
  left=4mm, right=4mm, top=3mm, bottom=3mm,
  before skip=\bloqueskip, after skip=\bloqueskip,
  before upper={{\sffamily\bfseries\footnotesize\color{subjectcolor}%
    \faFlask\enspace\MakeUppercase{Manos a la obra}}%
    \ifblank{#1}{}{\ {\color{black!35}\textbf{·}}\ %
      {\sffamily\bfseries\footnotesize\color{black!78}#1}}%
    \par\nopagebreak\vspace{0.45em}}
}
% Lista de materiales dentro de \begin{laboratorio}
\newcommand{\materiales}[1]{\par\noindent
  {\sffamily\bfseries\footnotesize\color{black!70}Materiales:}\ %
  {\footnotesize #1}\par\vspace{0.3em}}

% --- IMAGEN INLINE con texto alrededor (Santillana) -------------------
%   \figinline[r|l]{ancho}{archivo}{pie}   — colócala al inicio de un párrafo
\newcommand{\figinline}[4][r]{%
  \begin{wrapfigure}{#1}{#2}
  \centering\vspace{-0.5\baselineskip}
  \IfFileExists{assets/imgs/#3}{%
    \includegraphics[width=\linewidth]{assets/imgs/#3}}{%
    \fbox{\parbox[c][3cm][c]{0.92\linewidth}{\centering\sffamily\scriptsize
      \color{black!55}{\Large\faImage}\\[3pt][imagen pendiente]\\\ttfamily\tiny #3}}}%
  \par\vspace{2pt}{\footnotesize\sffamily\color{black!60}\textbf{Fig.}\enspace #4}%
  \vspace{-0.4\baselineskip}
  \end{wrapfigure}}

% --- APERTURA CON HERO LATERAL (imagen grande al lado del título) ------
%   \aperturahero{archivo}{título}{intro}  — alternativa a \aperturaseccion
\newcommand{\aperturahero}[3]{%
  \clearpage\refstepcounter{secportada}\stepcounter{section}%
  \addcontentsline{toc}{section}{\numseccion\enspace #2}%
  \noindent\begin{minipage}[c]{0.62\linewidth}
    {\InterDisplaySB\footnotesize\color{subjectcolor!85}\MakeUppercase{\seccionkicker}}\par
    \vspace{1.8mm}
    {\LARGE\InterDisplaySB\color{subjectcolor}%
      {\color{subjectcolor!35}\numseccion}\,\,#2}\par
    \vspace{2mm}{\color{subjectcolor!30}\rule{\linewidth}{0.8pt}}\par\vspace{2mm}
    {\InterLight\fontsize{12}{16}\selectfont\color{black!58}#3}
  \end{minipage}\hfill
  \begin{minipage}[c]{0.34\linewidth}\centering
    \IfFileExists{assets/imgs/#1}{%
      \includegraphics[width=\linewidth]{assets/imgs/#1}}{%
      \fbox{\parbox[c][3.4cm][c]{0.95\linewidth}{\centering\sffamily\scriptsize
        \color{black!55}{\Large\faImage}\\[3pt][imagen pendiente]\\\ttfamily\tiny #1}}}%
  \end{minipage}\par\vspace{1.3em}}

% --- MAPAS (conceptual y mental) — estilos TikZ con texto nítido -------
\tikzset{
  croot/.style={draw=subjectcolor, fill=subjectcolor, text=white, rounded corners=3pt,
    font=\sffamily\bfseries\footnotesize, align=center, inner sep=5pt, minimum height=9mm},
  cnode/.style={draw=subjectcolor!55, fill=subjectcolor!8, rounded corners=2.5pt,
    font=\sffamily\footnotesize, align=center, inner sep=4pt, text width=2.5cm},
  cleaf/.style={draw=subjectcolor!45, fill=white, rounded corners=2.5pt,
    font=\sffamily\footnotesize, align=center, inner sep=4pt},
  clink/.style={-{Stealth[length=5pt]}, subjectcolor!65, semithick},
  cline/.style={subjectcolor!55, semithick},
}
% Caja contenedora para un mapa (título + regla superior)
\newtcolorbox{mapabox}[2][Mapa conceptual]{
  enhanced, sharp corners, boxrule=0pt, frame hidden, colback=white,
  borderline north={1.4pt}{0pt}{subjectcolor},
  left=1mm, right=1mm, top=3mm, bottom=2mm,
  before skip=\bloqueskip, after skip=\bloqueskip,
  before upper={\blocklabel{subjectcolor}{#1}{#2}}
}

% --- Papel cuadriculado para graficar a mano (hojas de trabajo) --------
%   \papelcuadricula[filas]{ancho en cm}   (cada celda = 0.6 cm)
\newcommand{\papelcuadricula}[2][8]{%
  \par\smallskip\begin{center}
  \begin{tikzpicture}
    \draw[step=0.6cm, black!16, line width=0.3pt] (0,0) grid (#2,{#1*0.6});
    \draw[black!45, line width=0.7pt] (0,0) -- (#2,0);
    \draw[black!45, line width=0.7pt] (0,0) -- (0,{#1*0.6});
  \end{tikzpicture}
  \end{center}\par\smallskip}

% =====================================================================
%  APERTURA DE CAPÍTULO — título grande, número marca de agua, aire
% =====================================================================
\newcommand{\aperturacapitulo}[3]{% {número}{título}{subtítulo}
  \setcounter{unidadnum}{#1}\setcounter{secportada}{0}% fija unidad y reinicia secciones
  \thispagestyle{empty}%
  \begin{tikzpicture}[remember picture, overlay]
    \node[anchor=north east, text=subjectcolor!9,
          font=\InterDisplayThin]
      at ([xshift=-0.9cm, yshift=-2.2cm]current page.north east)
      {\fontsize{210}{210}\selectfont #1};
  \end{tikzpicture}%
  \begingroup\raggedright\null\vspace{3.8cm}\par
  {\InterDisplaySB\fontsize{12}{14}\selectfont\color{subjectcolor}%
    CAP\'ITULO\ #1}\par
  \vspace{0.18cm}{\color{black!16}\rule{4cm}{1pt}}\par
  \vspace{0.9cm}
  {\InterDisplaySB\fontsize{46}{50}\selectfont\color{black!88}#2}\par
  \vspace{0.5cm}{\color{subjectcolor}\rule{3.2cm}{2.5pt}}\par
  \vspace{0.55cm}
  {\InterLight\fontsize{15}{21}\selectfont\color{black!55}#3}\par
  \vfill
  {\sffamily\footnotesize\color{black!45}\hdrinstituto\quad
    \textcolor{black!22}{·}\quad Prof.\ Juan Diego A.\ Prada R.}\par
  \endgroup\clearpage}

% =====================================================================
%  DESBLOQUEA — pie discreto, no caja
% =====================================================================
\newcommand{\desbloquea}[1]{%
  \par\addvspace{\bloqueskip}%
  \noindent{\color{subjectcolor!25}\rule{\linewidth}{0.5pt}}\par\vspace{0.4em}%
  {\footnotesize\sffamily\color{black!60}%
    \textcolor{subjectcolor}{\faUnlock}\ \textbf{Dominar esto te desbloquea:}\ #1}\par}
; opcional tema de asignatura):
%     % =====================================================================
%  DESIGN SYSTEM v5 — EDITORIAL ITSTZ
%  Filosofía: texto primero · un solo color de asignatura · sin sombras
%  · sin marcos pesados · tratamientos planos (filete lateral / regla).
%
%  Vocabulario semántico (esto es TODO lo disponible):
%    \section / \subsection      jerarquía
%    definicion{título}          concepto/regla clave (filete color)
%    ejemplo{título}             ejemplo resuelto (numerado, sin caja)
%    procedimiento{título}       pasos (filete neutro) + \paso{n}{texto}
%    nota[título]                aparte único gris (sabías/dato/video/margen)
%    alerta{título}              advertencia ámbar (usar rara vez)
%    practica{título}            ejercicios/hoja/ICFES (regla superior)
%    \figura{archivo}{pie}       imagen real o placeholder automático
%    \desbloquea{lista}          pie discreto de prerrequisitos
%
%  REGLA DE ORO: las secciones (sNN.tex) NO llevan \vspace, color ni
%  minipage de maquetación. Solo contenido semántico. El diseño vive aquí.
% =====================================================================
\documentclass[11pt, letterpaper, twoside]{article}

\usepackage{fontspec}
\usepackage{polyglossia}
\setdefaultlanguage{spanish}
\usepackage[inner=2.1cm, outer=3.7cm, top=2.6cm, bottom=2.3cm,
            marginparwidth=2.55cm, marginparsep=0.45cm,
            headheight=16pt, headsep=18pt, footskip=24pt]{geometry}
\usepackage{microtype}
\usepackage{multicol}
\usepackage{xcolor}
\usepackage{graphicx}
\usepackage{tikz}
\usetikzlibrary{calc, arrows.meta, positioning}
\usepackage{wrapfig}
\usepackage[most]{tcolorbox}
\usepackage{amsmath, amssymb}
\usepackage{siunitx}
\usepackage{fontawesome5}
\usepackage{fancyhdr}
\usepackage{titlesec}
\usepackage{enumitem}
\usepackage{etoolbox}
\usepackage{titletoc}
\usepackage{array}
\usepackage{booktabs}
\usepackage{colortbl}
\usepackage{nicematrix}   % rejilla completa (hvlines) para tablas de hojas

% --- TIPOGRAFÍA -------------------------------------------------------
% Cuerpo: Charter (serif cálida y muy legible). La matemática (Latin
% Modern, serif) combina con ella. Títulos/etiquetas/UI: Inter Display.
\setmainfont{Charter}[
  UprightFont    = {Charter},
  BoldFont       = {Charter Bold},
  ItalicFont     = {Charter Italic},
  BoldItalicFont = {Charter Bold Italic},
  Ligatures = TeX, Numbers = Lining ]
\setsansfont{Inter}[
  UprightFont = Inter-Regular, BoldFont = Inter-SemiBold,
  ItalicFont  = Inter-Italic, Ligatures = TeX, Numbers = Lining ]
\newfontfamily\InterLight{Inter-Light}[ItalicFont=Inter-LightItalic, Ligatures=TeX]
\newfontfamily\InterDisplay{Inter Display}[Ligatures=TeX]
\newfontfamily\InterDisplaySB{Inter Display SemiBold}[Ligatures=TeX]
\newfontfamily\InterDisplayThin{Inter Display Thin}[Ligatures=TeX]
\newfontfamily\InterDisplayBlk{Inter Display Black}[Ligatures=TeX]

% --- RITMO --------------------------------------------------------------
\linespread{1.22}
\setlength{\parindent}{0pt}
\setlength{\parskip}{0.5em}
\newlength{\bloqueskip}\setlength{\bloqueskip}{1.0em}   % aire entre bloques
\setlength{\columnsep}{1.6em}                            % medianil multicols

% =====================================================================
%  COLOR — UNA sola variable de asignatura + neutros. Sin arcoíris.
% =====================================================================
\definecolor{mat} {RGB}{ 30,  64, 120}   % azul profundo
\definecolor{cnat}{RGB}{ 16, 110,  80}   % esmeralda
\definecolor{csoc}{RGB}{165,  75,  45}   % terracota
\definecolor{eco} {RGB}{ 70,  55, 130}   % índigo
\definecolor{alertcolor}{RGB}{180, 95, 10}  % ámbar — solo advertencias

\colorlet{subjectcolor}{mat}             % por defecto
\newcommand{\setsubject}[1]{\colorlet{subjectcolor}{#1}}

% =====================================================================
%  JERARQUÍA TIPOGRÁFICA
% =====================================================================
\setcounter{secnumdepth}{0}
\titleformat{\section}
  {\LARGE\InterDisplaySB\color{subjectcolor}}{}{0pt}{}
  [\vspace{2pt}{\color{subjectcolor!30}\titlerule[0.8pt]}]
\titlespacing*{\section}{0pt}{3.2ex plus 0.6ex}{1.8ex plus 0.3ex}

\titleformat{\subsection}
  {\large\InterDisplaySB\color{black!82}}{}{0pt}{}
\titlespacing*{\subsection}{0pt}{2.0ex plus 0.4ex}{0.5ex}

% --- TABLA DE CONTENIDO -------------------------------------------------
\titlecontents{section}[0em]
  {\addvspace{0.45pc}\sffamily\bfseries\small\color{black!85}}
  {}{}{\titlerule*[0.35pc]{.}\contentspage}
\titlecontents{subsection}[1.4em]
  {\sffamily\footnotesize\color{black!60}}
  {}{}{\titlerule*[0.35pc]{.}\contentspage}

% =====================================================================
%  ENCABEZADO / PIE — editorial limpio (regla fina, sin franjas)
% =====================================================================
\newcommand{\hdrinstituto}{Instituto Técnico Santo Tomás · Zapatoca}
\newcommand{\hdrcapitulo}{}            % se fija en el documento maestro
\pagestyle{fancy}\fancyhf{}
\renewcommand{\headrulewidth}{0pt}
\fancyhead[LE]{\footnotesize\sffamily\color{black!55}\hdrinstituto}
\fancyhead[RO]{\footnotesize\sffamily\color{subjectcolor}\hdrcapitulo}
\fancyhead[LO,RE]{}
\renewcommand{\headrule}{\vspace{2pt}{\color{black!18}\hrule height 0.4pt}}

% --- PIE: número exterior (cuadro de color) + crédito interior + filete ---
\renewcommand{\footrulewidth}{0pt}
\renewcommand{\footrule}{{\color{black!15}\hrule height 0.4pt}\vspace{3pt}}
\newcommand{\pagenumbox}{%
  \tikz[baseline=-3.2pt]{\node[fill=subjectcolor, text=white, inner xsep=7pt,
     inner ysep=3pt, font=\sffamily\bfseries\footnotesize]{\thepage};}}
\fancyfoot[RO,LE]{\pagenumbox}
\fancyfoot[LO,RE]{\footnotesize\sffamily\color{black!45}%
  \hdrcapitulo\ \textcolor{black!22}{·}\ Prof.\ Juan Diego A.\ Prada R.}

% =====================================================================
%  HELPERS
% =====================================================================
% Etiqueta de bloque:  TAG · título   (color del tag configurable)
\newcommand{\blocklabel}[3]{% #1=color  #2=TAG  #3=título
  {\sffamily\bfseries\footnotesize\color{#1}\MakeUppercase{#2}}%
  \ifblank{#3}{}{\,{\color{black!35}\textbf{·}}\,%
    {\sffamily\bfseries\footnotesize\color{black!78}#3}}%
  \par\nopagebreak\vspace{0.4em}}

% --- CAJAS AL MARGEN — personajes históricos e hitos -------------------
% Van en el margen exterior ancho (\marginpar respeta la geometría twoside
% y flota para no chocar). Son contenido semántico: colócalas en el cuerpo,
% cerca del párrafo que comentan, ENTRE párrafos. NO usarlas dentro de
% multicols, tcolorbox, tablas, ni en la misma línea que un \section.
%   \personaje{Nombre}{años · rol}{descripción breve}
%   \hito{año}{qué pasó (breve)}
\setlength{\marginparpush}{8pt}
\newcommand{\margenbox}[1]{\marginpar{%
  \begingroup\raggedright\footnotesize\sffamily
  \setlength{\parskip}{0.18em}\linespread{1.05}\selectfont
  #1\par\endgroup}}
\newcommand{\personaje}[3]{\margenbox{%
  {\color{subjectcolor}\rule{\linewidth}{1.1pt}}\par
  {\InterDisplaySB\scriptsize\color{subjectcolor}#1}\par
  {\fontsize{6.8pt}{8pt}\selectfont\itshape\color{black!50}#2}\par
  \vspace{0.15em}{\scriptsize\color{black!72}#3}}}
\newcommand{\hito}[2]{\margenbox{%
  {\sffamily\bfseries\scriptsize\colorbox{subjectcolor}{\textcolor{white}{\,#1\,}}}\par
  \vspace{0.15em}{\scriptsize\color{black!72}#2}}}

% =====================================================================
%  ENTORNOS — todos planos: filete lateral, sin marco, sin sombra
% =====================================================================

% --- DEFINICIÓN / concepto clave ---------------------------------------
\newtcolorbox{definicion}[1]{
  enhanced, breakable, sharp corners, boxrule=0pt, frame hidden,
  colback=subjectcolor!6!white,
  borderline west={2.2pt}{0pt}{subjectcolor},
  left=5mm, right=4mm, top=3mm, bottom=3mm,
  before skip=\bloqueskip, after skip=\bloqueskip,
  before upper={\blocklabel{subjectcolor}{Definición}{#1}}
}

% --- PROCEDIMIENTO / pasos ---------------------------------------------
\newtcolorbox{procedimiento}[1]{
  enhanced, breakable, sharp corners, boxrule=0pt, frame hidden,
  colback=black!3!white,
  borderline west={2.2pt}{0pt}{black!45},
  left=5mm, right=4mm, top=3mm, bottom=3mm,
  before skip=\bloqueskip, after skip=\bloqueskip,
  before upper={\blocklabel{black!60}{Procedimiento}{#1}}
}
\newcommand{\paso}[2]{%
  \par\vspace{0.45em}\noindent
  \tikz[baseline=-0.6ex]\node[circle, fill=subjectcolor, text=white,
     inner sep=0pt, minimum size=1.45em,
     font=\sffamily\bfseries\footnotesize]{#1};\hspace{0.5em}#2\par}

% Glosa "¿por qué?" bajo un paso — narración en lenguaje cercano
\newcommand{\porque}[1]{\par\vspace{0.1em}\noindent\hspace{2.0em}%
  {\sffamily\footnotesize\itshape\color{subjectcolor}¿por qué?\,}%
  {\sffamily\footnotesize\itshape\color{black!55}#1}\par}

% Operación glosada: cuenta a la izquierda + porqué al margen derecho
% (estilo "comentario al lado del código"). Úsala dentro de ejemplo.
% \opg{ecuación}{porqué}   — primer arg en modo display (mat, cap0)
% \opgt{etiqueta}{porqué}  — primer arg en texto sans bold (física, ciencias)
\newcommand{\opg}[2]{\par\vspace{0.25em}\noindent
  \makebox[0.46\linewidth][l]{$\displaystyle #1$}%
  {\sffamily\footnotesize\itshape\color{black!52}#2}\par}
\newcommand{\opgt}[2]{\par\vspace{0.25em}\noindent
  \makebox[0.46\linewidth][l]{\sffamily\bfseries\color{subjectcolor!80}#1}%
  {\sffamily\footnotesize\itshape\color{black!52}#2}\par}

% --- NOTA / aparte único (sabías · dato · video · margen) --------------
\newtcolorbox{nota}[1][]{
  enhanced, breakable, sharp corners, boxrule=0pt, frame hidden,
  colback=black!4!white,
  borderline west={2.2pt}{0pt}{black!35},
  left=5mm, right=4mm, top=3mm, bottom=3mm,
  before skip=\bloqueskip, after skip=\bloqueskip,
  before upper={\ifblank{#1}{}{\blocklabel{black!55}{Nota}{#1}}}
}

% --- ALERTA / advertencia (única en ámbar, uso escaso) -----------------
\newtcolorbox{alerta}[1]{
  enhanced, breakable, sharp corners, boxrule=0pt, frame hidden,
  colback=alertcolor!7!white,
  borderline west={2.2pt}{0pt}{alertcolor},
  left=5mm, right=4mm, top=3mm, bottom=3mm,
  before skip=\bloqueskip, after skip=\bloqueskip,
  before upper={\blocklabel{alertcolor}{Atención}{#1}}
}
% Sub-ítem de alerta (antes \caso)
\newcommand{\caso}[2]{\par\vspace{0.35em}\noindent
  {\sffamily\bfseries\color{alertcolor!85!black}#1.}\enspace #2\par}

% --- CONTRASTE (la virtud y su exceso → pregunta que cuestiona) --------
%   \contraste{Concepto}{La virtud / lo valioso}{Su exceso / su riesgo}{Pregunta}
\newtcolorbox{contrastebox}{
  enhanced, breakable, sharp corners, boxrule=0pt, frame hidden,
  colback=subjectcolor!5!white,
  borderline west={2.2pt}{0pt}{subjectcolor},
  left=5mm, right=4mm, top=3mm, bottom=3mm,
  before skip=\bloqueskip, after skip=\bloqueskip,
}
\newcommand{\contraste}[4]{%
\begin{contrastebox}
\blocklabel{subjectcolor}{Cuestiona esto}{\faBalanceScale\ #1}
\noindent\begin{minipage}[t]{0.45\linewidth}
  {\sffamily\footnotesize\bfseries\color{subjectcolor!85!black}La virtud}\par
  \vspace{0.2em}#2
\end{minipage}\hfill
{\color{subjectcolor!35}\vrule width 0.4pt}\hfill
\begin{minipage}[t]{0.45\linewidth}
  {\sffamily\footnotesize\bfseries\color{alertcolor!85!black}Su exceso}\par
  \vspace{0.2em}#3
\end{minipage}
\par\vspace{0.6em}
{\color{subjectcolor!25}\titlerule[0.6pt]}\par\vspace{0.4em}
\noindent{\sffamily\itshape\color{subjectcolor}\faArrowRight\ #4}
\end{contrastebox}}

% --- EJEMPLO (sin caja: regla fina + etiqueta numerada) ----------------
\newcounter{ejemplo}[section]
\newenvironment{ejemplo}[1]
  {\par\addvspace{\bloqueskip}\refstepcounter{ejemplo}%
   \noindent{\color{subjectcolor!28}\rule{\linewidth}{0.5pt}}\par\vspace{0.4em}%
   \noindent{\sffamily\bfseries\color{subjectcolor}Ejemplo \theejemplo.}%
   \enspace{\sffamily\bfseries\color{black!78}#1}\par\vspace{0.3em}}
  {\par\addvspace{\bloqueskip}}

% --- PRÁCTICA (ejercicios / hoja / ICFES) ------------------------------
\newtcolorbox{practica}[1]{
  enhanced, breakable, sharp corners, boxrule=0pt, frame hidden,
  colback=white,
  borderline north={1.4pt}{0pt}{subjectcolor},
  left=0mm, right=0mm, top=3mm, bottom=2mm,
  before skip=\bloqueskip, after skip=\bloqueskip,
  before upper={\blocklabel{subjectcolor}{Práctica}{#1}\raggedright}
}

% Ejercicio numerado con nivel opcional [básico|intermedio|reto]
\newcounter{ejer}[section]
\newcommand{\ejer}[2][]{%
  \stepcounter{ejer}\par\vspace{0.5em}\noindent
  {\sffamily\bfseries\color{subjectcolor}\theejer.}\enspace #2%
  \ifblank{#1}{}{\ \textcolor{black!45}{\footnotesize[\textit{#1}]}}\par}

% Niveles = un solo color a tres intensidades (NO nuevos colores)
\newcommand{\nivel}[2]{\par\smallskip\noindent
  \textcolor{subjectcolor!#1}{\small\faCircle}\enspace
  {\sffamily\bfseries\footnotesize\color{black!70}#2}\par}
\newcommand{\basico}{\nivel{40}{Básico}}
\newcommand{\intermedio}{\nivel{70}{Intermedio}}
\newcommand{\reto}{\nivel{100}{Reto}}

% Selección múltiple tipo ICFES — opciones con burbuja (estilo Saber)
\newcounter{pregicfes}[section]
\newcommand{\bubbleletter}[1]{\tikz[baseline=(bc.base)]\node[circle,
  draw=subjectcolor!55, line width=0.6pt, inner sep=1.5pt,
  font=\sffamily\bfseries\scriptsize, text=subjectcolor!90!black](bc){#1};}
\newcommand{\icfes}[5]{%
  \stepcounter{pregicfes}\par\vspace{0.6em}\noindent
  {\sffamily\bfseries\color{subjectcolor}\thepregicfes.}\enspace #1\par
  \begin{enumerate}[label={\protect\bubbleletter{\Alph*}}, leftmargin=1.5cm,
                    itemsep=3pt, topsep=3pt]
    \item #2 \item #3 \item #4 \item #5
  \end{enumerate}}
% Pregunta abierta de argumentación
\newcommand{\pregabierta}[1]{\par\vspace{0.6em}\noindent
  {\sffamily\bfseries\color{subjectcolor}\faComment\ Argumenta:}\enspace
  \textit{#1}\par\lineaescritura\lineaescritura\lineaescritura}

% --- LÍNEAS DE RESPUESTA -----------------------------------------------
\newcommand{\linea}[1][7cm]{\underline{\hspace{#1}}}
\newcommand{\lineaescritura}{\par\vspace{0.9em}\noindent
  {\color{black!30}\rule{\linewidth}{0.3pt}}}

% =====================================================================
%  EVALUACIÓN TIPO ICFES — cierre de unidad (banner + 2 columnas)
%  Uso:  \begin{evaluacion} \begin{multicols}{2} \icfes... \end{multicols}
%        \pregabierta{...} \end{evaluacion}
% =====================================================================
\newtcolorbox{evaluacion}[1][cierre de sección]{
  enhanced, breakable, sharp corners,
  colback=white, colframe=subjectcolor, boxrule=1pt,
  colbacktitle=subjectcolor, coltitle=white,
  fonttitle=\InterDisplaySB\normalsize,
  title={\faClipboardCheck\ \ EVALUACIÓN TIPO ICFES\ \ %
         \textnormal{\sffamily\small· #1}},
  toptitle=2mm, bottomtitle=2mm, left=4mm, right=4mm, top=3.5mm, bottom=3mm,
  before={\clearpage},
  before upper={\small\itshape\color{black!55}%
    Selección múltiple. Marca una sola opción por pregunta.\par\medskip\raggedright}
}

% =====================================================================
%  HOJA DE TRABAJO — recortable: borde punteado, escala de grises,
%  en página aparte (clearpage antes; el contenido que sigue continúa).
% =====================================================================
\newtcolorbox{hoja}[1]{
  enhanced, breakable, sharp corners,
  colback=black!2!white, boxrule=0pt, frame hidden,
  borderline={0.8pt}{0pt}{black!50, dash pattern=on 3pt off 2.5pt},
  before={\clearpage}, after={\par},
  left=6mm, right=6mm, top=6mm, bottom=6mm,
  before upper={%
    {\sffamily\bfseries\footnotesize\color{black!60}\faCut\ \ HOJA DE TRABAJO}%
    {\color{black!35}\ \textbf{·}\ }{\sffamily\bfseries\footnotesize\color{black!80}#1}%
    \hfill{\sffamily\footnotesize\color{black!55}%
      Nombre:\ \underline{\hspace{4.2cm}}\quad Fecha:\ \underline{\hspace{2cm}}}\par
    \vspace{2mm}{\color{black!30}\hrule height 0.4pt}\vspace{3mm}\raggedright}
}

% =====================================================================
%  TABLAHOJA — tabla con TODAS las líneas (rejilla) para hojas de trabajo.
%  Celdas altas para llenar a mano. Centrada. El encabezado se marca con
%  \rowcolor en la primera fila.
%  Uso:  \begin{tablahoja}{ccc} enc & enc & enc \\ a & b & c \\ \end{tablahoja}
% =====================================================================
\NewDocumentEnvironment{tablahoja}{m}{%
  \par\smallskip\centering\small\setlength{\tabcolsep}{4pt}%
  \begin{NiceTabular}{#1}[hvlines, cell-space-top-limit=5pt,
     cell-space-bottom-limit=9pt, colortbl-like]%
}{%
  \end{NiceTabular}\par\smallskip
}

% =====================================================================
%  FIGURAS — imagen real o placeholder automático (\IfFileExists)
% =====================================================================
\newcounter{figura}[section]
\newcommand{\figura}[3][0.7\linewidth]{% [ancho]{archivo}{pie}
  \par\addvspace{\bloqueskip}\refstepcounter{figura}%
  \begin{center}%
  \IfFileExists{assets/imgs/#2}{%
    \includegraphics[width=#1]{assets/imgs/#2}}{%
    \begin{tikzpicture}
      \node[draw=black!30, dashed, line width=0.6pt, fill=black!3,
            text width=0.66\linewidth, minimum height=2.6cm,
            align=center, inner sep=8pt,
            text=black!55, font=\sffamily\footnotesize]
        {{\Large\faImage}\\[4pt]\textbf{[ imagen pendiente ]}\\#3};
    \end{tikzpicture}}%
  \par\vspace{0.4em}%
  {\footnotesize\sffamily\color{black!60}%
    \textbf{Figura \thefigura.}\enspace #3}%
  \end{center}\par\addvspace{\bloqueskip}}

% =====================================================================
%  APERTURA DE SECCIÓN — marco branded de serie + título + "qué verás"
%  El "marco" es CONSTANTE en todas las secciones (= identidad de marca):
%  número grande de serie, kicker, cuadro-firma, franja de acento y sello
%  del instituto. El arte (imagen real o placeholder) va DENTRO del marco.
%  Uso (1ª línea de cada sNN.tex):
%    \aperturaseccion{ruta/img.png}{Título}{Una o dos frases.}
%  Personaliza por documento (opcional):  \renewcommand{\seccionkicker}{...}
% =====================================================================
\newcounter{secportada}
\newcounter{unidadnum}                       % nº de unidad actual (lo fija \aperturacapitulo)
% Número visible de sección = Unidad.Sección (ej. 1.2). Si no hay unidad
% definida (unidadnum=0), cae al número simple de sección.
\newcommand{\numseccion}{\ifnum\value{unidadnum}>0 \theunidadnum.\thesecportada\else\thesecportada\fi}
\newcommand{\seccionkicker}{\hdrcapitulo}   % por defecto, lo del encabezado

% El escenario branded (contenedor constante del arte)
\newtcolorbox{secstage}{%
  enhanced, sharp corners, boxrule=0pt, frame hidden,
  colback=subjectcolor!6, halign=center, valign=center,
  height=5.0cm, left=2mm, right=2mm, top=10mm, bottom=8mm,
  overlay={%
    % cuadro-firma + kicker (arriba-izquierda)
    \fill[subjectcolor] ([shift={(4mm,-4mm)}]frame.north west) rectangle ++(2.6mm,-2.6mm);
    \node[anchor=north west, font=\InterDisplaySB\footnotesize, text=subjectcolor!85]
      at ([shift={(8.5mm,-3.4mm)}]frame.north west) {\seccionkicker};
    % número de serie grande y tenue (arriba-derecha) = marca
    \node[anchor=north east, font=\InterDisplayBlk, text=subjectcolor!18, inner sep=0pt]
      at ([shift={(-4mm,-1.5mm)}]frame.north east)
      {\fontsize{34}{34}\selectfont \numseccion};
    % franja de acento (abajo-izquierda)
    \fill[subjectcolor] ([shift={(4mm,4mm)}]frame.south west) rectangle ++(2.8cm,2.2mm);
    % sello del instituto (abajo-derecha)
    \node[anchor=south east, font=\sffamily\scriptsize, text=subjectcolor!55]
      at ([shift={(-4mm,3.4mm)}]frame.south east) {ITSTZ · Zapatoca};
  }}

\newcommand{\aperturaseccion}[3]{% {archivo}{título}{intro}
  \clearpage\refstepcounter{secportada}%
  \begin{secstage}
  \IfFileExists{assets/imgs/#1}{%
    \includegraphics[width=0.9\linewidth, height=3.2cm, keepaspectratio]{assets/imgs/#1}}{%
    \begin{tikzpicture}
      \node[draw=subjectcolor!35, dashed, line width=0.6pt, fill=white,
            minimum width=0.86\linewidth, minimum height=2.9cm,
            align=center, text=subjectcolor!50, font=\sffamily\small]
        {{\fontsize{22}{22}\selectfont\faImage}\\[5pt]
         \textbf{[ imagen pendiente ]}\\[1pt]
         \scriptsize\ttfamily assets/imgs/#1};
    \end{tikzpicture}}%
  \end{secstage}
  \vspace{0.5em}
  {\color{subjectcolor}\rule{3.2cm}{2.5pt}}\par\vspace{0.3em}
  \section{{\color{subjectcolor}\numseccion}\enspace #2}%
  {\InterLight\fontsize{12.5}{17}\selectfont\color{black!58} #3\par}
  \vspace{0.4em}}

% =====================================================================
%  DESTACADO — "lo esencial": triángulo lateral + degradado muy sutil
% =====================================================================
\newtcolorbox{destacado}[1]{
  enhanced, breakable, sharp corners, boxrule=0pt, frame hidden,
  interior style={left color=subjectcolor!16, right color=subjectcolor!2},
  borderline west={2.4pt}{0pt}{subjectcolor},
  left=7mm, right=5mm, top=3.5mm, bottom=3.5mm,
  before skip=\bloqueskip, after skip=\bloqueskip,
  overlay unbroken and first={%
    \fill[subjectcolor] ([yshift=-4.5mm]frame.north west)
        -- ++(5.5pt,2.75pt) -- ++(0,-5.5pt) -- cycle;},
  before upper={\blocklabel{subjectcolor}{Lo esencial}{#1}}
}

% --- Resaltado inline suave ---
\newtcbox{\resaltar}{on line, nobeforeafter, boxsep=0pt,
  left=2.5pt, right=2.5pt, top=1pt, bottom=1pt, arc=2.5pt,
  colback=subjectcolor!15, colframe=subjectcolor!15}

% =====================================================================
%  EXTENSIONES SANTILLANA (opt-in) — cajas temáticas para dar variedad
%  visual y romper la monotonía. NO alteran los entornos base; úsalas
%  donde aporten. Mantienen el color de asignatura (coherencia de marca).
% =====================================================================

% --- ¿SABÍAS QUE? — curiosidad breve · tarjeta redondeada con bombillo --
\newtcolorbox{sabias}{
  enhanced, breakable, arc=2.5mm, boxrule=0.8pt,
  colback=subjectcolor!6!white, colframe=subjectcolor!30,
  left=4mm, right=4mm, top=3mm, bottom=3mm,
  before skip=\bloqueskip, after skip=\bloqueskip,
  before upper={{\sffamily\bfseries\footnotesize\color{subjectcolor}%
    \faLightbulb\enspace\MakeUppercase{¿Sabías que?}}%
    \par\nopagebreak\vspace{0.45em}}
}

% --- CIENCIA Y SOCIEDAD — impacto/ética · banner de título -------------
\newtcolorbox{cienciasociedad}[1]{
  enhanced, breakable, sharp corners,
  colback=subjectcolor!5!white, colframe=subjectcolor, boxrule=0pt,
  colbacktitle=subjectcolor, coltitle=white,
  fonttitle=\sffamily\bfseries\footnotesize,
  title={\faGlobeAmericas\ \ \MakeUppercase{Ciencia y sociedad}%
         \ifblank{#1}{}{\ \ \textnormal{\sffamily\small· #1}}},
  toptitle=1.6mm, bottomtitle=1.6mm,
  left=4mm, right=4mm, top=3mm, bottom=3mm,
  before skip=\bloqueskip, after skip=\bloqueskip,
}

% --- MANOS A LA OBRA — práctica experimental · marco sólido con matraz --
\newtcolorbox{laboratorio}[1]{
  enhanced, breakable, sharp corners, arc=0pt,
  colback=white, colframe=subjectcolor!55, boxrule=1pt,
  left=4mm, right=4mm, top=3mm, bottom=3mm,
  before skip=\bloqueskip, after skip=\bloqueskip,
  before upper={{\sffamily\bfseries\footnotesize\color{subjectcolor}%
    \faFlask\enspace\MakeUppercase{Manos a la obra}}%
    \ifblank{#1}{}{\ {\color{black!35}\textbf{·}}\ %
      {\sffamily\bfseries\footnotesize\color{black!78}#1}}%
    \par\nopagebreak\vspace{0.45em}}
}
% Lista de materiales dentro de \begin{laboratorio}
\newcommand{\materiales}[1]{\par\noindent
  {\sffamily\bfseries\footnotesize\color{black!70}Materiales:}\ %
  {\footnotesize #1}\par\vspace{0.3em}}

% --- IMAGEN INLINE con texto alrededor (Santillana) -------------------
%   \figinline[r|l]{ancho}{archivo}{pie}   — colócala al inicio de un párrafo
\newcommand{\figinline}[4][r]{%
  \begin{wrapfigure}{#1}{#2}
  \centering\vspace{-0.5\baselineskip}
  \IfFileExists{assets/imgs/#3}{%
    \includegraphics[width=\linewidth]{assets/imgs/#3}}{%
    \fbox{\parbox[c][3cm][c]{0.92\linewidth}{\centering\sffamily\scriptsize
      \color{black!55}{\Large\faImage}\\[3pt][imagen pendiente]\\\ttfamily\tiny #3}}}%
  \par\vspace{2pt}{\footnotesize\sffamily\color{black!60}\textbf{Fig.}\enspace #4}%
  \vspace{-0.4\baselineskip}
  \end{wrapfigure}}

% --- APERTURA CON HERO LATERAL (imagen grande al lado del título) ------
%   \aperturahero{archivo}{título}{intro}  — alternativa a \aperturaseccion
\newcommand{\aperturahero}[3]{%
  \clearpage\refstepcounter{secportada}\stepcounter{section}%
  \addcontentsline{toc}{section}{\numseccion\enspace #2}%
  \noindent\begin{minipage}[c]{0.62\linewidth}
    {\InterDisplaySB\footnotesize\color{subjectcolor!85}\MakeUppercase{\seccionkicker}}\par
    \vspace{1.8mm}
    {\LARGE\InterDisplaySB\color{subjectcolor}%
      {\color{subjectcolor!35}\numseccion}\,\,#2}\par
    \vspace{2mm}{\color{subjectcolor!30}\rule{\linewidth}{0.8pt}}\par\vspace{2mm}
    {\InterLight\fontsize{12}{16}\selectfont\color{black!58}#3}
  \end{minipage}\hfill
  \begin{minipage}[c]{0.34\linewidth}\centering
    \IfFileExists{assets/imgs/#1}{%
      \includegraphics[width=\linewidth]{assets/imgs/#1}}{%
      \fbox{\parbox[c][3.4cm][c]{0.95\linewidth}{\centering\sffamily\scriptsize
        \color{black!55}{\Large\faImage}\\[3pt][imagen pendiente]\\\ttfamily\tiny #1}}}%
  \end{minipage}\par\vspace{1.3em}}

% --- MAPAS (conceptual y mental) — estilos TikZ con texto nítido -------
\tikzset{
  croot/.style={draw=subjectcolor, fill=subjectcolor, text=white, rounded corners=3pt,
    font=\sffamily\bfseries\footnotesize, align=center, inner sep=5pt, minimum height=9mm},
  cnode/.style={draw=subjectcolor!55, fill=subjectcolor!8, rounded corners=2.5pt,
    font=\sffamily\footnotesize, align=center, inner sep=4pt, text width=2.5cm},
  cleaf/.style={draw=subjectcolor!45, fill=white, rounded corners=2.5pt,
    font=\sffamily\footnotesize, align=center, inner sep=4pt},
  clink/.style={-{Stealth[length=5pt]}, subjectcolor!65, semithick},
  cline/.style={subjectcolor!55, semithick},
}
% Caja contenedora para un mapa (título + regla superior)
\newtcolorbox{mapabox}[2][Mapa conceptual]{
  enhanced, sharp corners, boxrule=0pt, frame hidden, colback=white,
  borderline north={1.4pt}{0pt}{subjectcolor},
  left=1mm, right=1mm, top=3mm, bottom=2mm,
  before skip=\bloqueskip, after skip=\bloqueskip,
  before upper={\blocklabel{subjectcolor}{#1}{#2}}
}

% --- Papel cuadriculado para graficar a mano (hojas de trabajo) --------
%   \papelcuadricula[filas]{ancho en cm}   (cada celda = 0.6 cm)
\newcommand{\papelcuadricula}[2][8]{%
  \par\smallskip\begin{center}
  \begin{tikzpicture}
    \draw[step=0.6cm, black!16, line width=0.3pt] (0,0) grid (#2,{#1*0.6});
    \draw[black!45, line width=0.7pt] (0,0) -- (#2,0);
    \draw[black!45, line width=0.7pt] (0,0) -- (0,{#1*0.6});
  \end{tikzpicture}
  \end{center}\par\smallskip}

% =====================================================================
%  APERTURA DE CAPÍTULO — título grande, número marca de agua, aire
% =====================================================================
\newcommand{\aperturacapitulo}[3]{% {número}{título}{subtítulo}
  \setcounter{unidadnum}{#1}\setcounter{secportada}{0}% fija unidad y reinicia secciones
  \thispagestyle{empty}%
  \begin{tikzpicture}[remember picture, overlay]
    \node[anchor=north east, text=subjectcolor!9,
          font=\InterDisplayThin]
      at ([xshift=-0.9cm, yshift=-2.2cm]current page.north east)
      {\fontsize{210}{210}\selectfont #1};
  \end{tikzpicture}%
  \begingroup\raggedright\null\vspace{3.8cm}\par
  {\InterDisplaySB\fontsize{12}{14}\selectfont\color{subjectcolor}%
    CAP\'ITULO\ #1}\par
  \vspace{0.18cm}{\color{black!16}\rule{4cm}{1pt}}\par
  \vspace{0.9cm}
  {\InterDisplaySB\fontsize{46}{50}\selectfont\color{black!88}#2}\par
  \vspace{0.5cm}{\color{subjectcolor}\rule{3.2cm}{2.5pt}}\par
  \vspace{0.55cm}
  {\InterLight\fontsize{15}{21}\selectfont\color{black!55}#3}\par
  \vfill
  {\sffamily\footnotesize\color{black!45}\hdrinstituto\quad
    \textcolor{black!22}{·}\quad Prof.\ Juan Diego A.\ Prada R.}\par
  \endgroup\clearpage}

% =====================================================================
%  DESBLOQUEA — pie discreto, no caja
% =====================================================================
\newcommand{\desbloquea}[1]{%
  \par\addvspace{\bloqueskip}%
  \noindent{\color{subjectcolor!25}\rule{\linewidth}{0.5pt}}\par\vspace{0.4em}%
  {\footnotesize\sffamily\color{black!60}%
    \textcolor{subjectcolor}{\faUnlock}\ \textbf{Dominar esto te desbloquea:}\ #1}\par}

%     % =====================================================================
%  TEMA · CIENCIAS NATURALES — variante del design system (GUÍAS)
%  Identidad visual propia, distinta de Matemáticas (azul/pizarra):
%    · color base ESMERALDA (\setsubject{cnat})
%    · acento LIMA orgánico solo para ornamentos (no rompe la regla de
%      "un solo color de asignatura")
%    · iconografía de ciencias (hoja · ADN · matraz · átomo)
%    · reutiliza las cajas de ciencia que ya trae el preamble:
%      \begin{laboratorio}, \begin{sabias}, \begin{cienciasociedad}
%
%  Sirve a futuro para BIOLOGÍA, QUÍMICA y FÍSICA EXPERIMENTAL.
%
%  INVOCACIÓN (una sola línea, tras % =====================================================================
%  DESIGN SYSTEM v5 — EDITORIAL ITSTZ
%  Filosofía: texto primero · un solo color de asignatura · sin sombras
%  · sin marcos pesados · tratamientos planos (filete lateral / regla).
%
%  Vocabulario semántico (esto es TODO lo disponible):
%    \section / \subsection      jerarquía
%    definicion{título}          concepto/regla clave (filete color)
%    ejemplo{título}             ejemplo resuelto (numerado, sin caja)
%    procedimiento{título}       pasos (filete neutro) + \paso{n}{texto}
%    nota[título]                aparte único gris (sabías/dato/video/margen)
%    alerta{título}              advertencia ámbar (usar rara vez)
%    practica{título}            ejercicios/hoja/ICFES (regla superior)
%    \figura{archivo}{pie}       imagen real o placeholder automático
%    \desbloquea{lista}          pie discreto de prerrequisitos
%
%  REGLA DE ORO: las secciones (sNN.tex) NO llevan \vspace, color ni
%  minipage de maquetación. Solo contenido semántico. El diseño vive aquí.
% =====================================================================
\documentclass[11pt, letterpaper, twoside]{article}

\usepackage{fontspec}
\usepackage{polyglossia}
\setdefaultlanguage{spanish}
\usepackage[inner=2.1cm, outer=3.7cm, top=2.6cm, bottom=2.3cm,
            marginparwidth=2.55cm, marginparsep=0.45cm,
            headheight=16pt, headsep=18pt, footskip=24pt]{geometry}
\usepackage{microtype}
\usepackage{multicol}
\usepackage{xcolor}
\usepackage{graphicx}
\usepackage{tikz}
\usetikzlibrary{calc, arrows.meta, positioning}
\usepackage{wrapfig}
\usepackage[most]{tcolorbox}
\usepackage{amsmath, amssymb}
\usepackage{siunitx}
\usepackage{fontawesome5}
\usepackage{fancyhdr}
\usepackage{titlesec}
\usepackage{enumitem}
\usepackage{etoolbox}
\usepackage{titletoc}
\usepackage{array}
\usepackage{booktabs}
\usepackage{colortbl}
\usepackage{nicematrix}   % rejilla completa (hvlines) para tablas de hojas

% --- TIPOGRAFÍA -------------------------------------------------------
% Cuerpo: Charter (serif cálida y muy legible). La matemática (Latin
% Modern, serif) combina con ella. Títulos/etiquetas/UI: Inter Display.
\setmainfont{Charter}[
  UprightFont    = {Charter},
  BoldFont       = {Charter Bold},
  ItalicFont     = {Charter Italic},
  BoldItalicFont = {Charter Bold Italic},
  Ligatures = TeX, Numbers = Lining ]
\setsansfont{Inter}[
  UprightFont = Inter-Regular, BoldFont = Inter-SemiBold,
  ItalicFont  = Inter-Italic, Ligatures = TeX, Numbers = Lining ]
\newfontfamily\InterLight{Inter-Light}[ItalicFont=Inter-LightItalic, Ligatures=TeX]
\newfontfamily\InterDisplay{Inter Display}[Ligatures=TeX]
\newfontfamily\InterDisplaySB{Inter Display SemiBold}[Ligatures=TeX]
\newfontfamily\InterDisplayThin{Inter Display Thin}[Ligatures=TeX]
\newfontfamily\InterDisplayBlk{Inter Display Black}[Ligatures=TeX]

% --- RITMO --------------------------------------------------------------
\linespread{1.22}
\setlength{\parindent}{0pt}
\setlength{\parskip}{0.5em}
\newlength{\bloqueskip}\setlength{\bloqueskip}{1.0em}   % aire entre bloques
\setlength{\columnsep}{1.6em}                            % medianil multicols

% =====================================================================
%  COLOR — UNA sola variable de asignatura + neutros. Sin arcoíris.
% =====================================================================
\definecolor{mat} {RGB}{ 30,  64, 120}   % azul profundo
\definecolor{cnat}{RGB}{ 16, 110,  80}   % esmeralda
\definecolor{csoc}{RGB}{165,  75,  45}   % terracota
\definecolor{eco} {RGB}{ 70,  55, 130}   % índigo
\definecolor{alertcolor}{RGB}{180, 95, 10}  % ámbar — solo advertencias

\colorlet{subjectcolor}{mat}             % por defecto
\newcommand{\setsubject}[1]{\colorlet{subjectcolor}{#1}}

% =====================================================================
%  JERARQUÍA TIPOGRÁFICA
% =====================================================================
\setcounter{secnumdepth}{0}
\titleformat{\section}
  {\LARGE\InterDisplaySB\color{subjectcolor}}{}{0pt}{}
  [\vspace{2pt}{\color{subjectcolor!30}\titlerule[0.8pt]}]
\titlespacing*{\section}{0pt}{3.2ex plus 0.6ex}{1.8ex plus 0.3ex}

\titleformat{\subsection}
  {\large\InterDisplaySB\color{black!82}}{}{0pt}{}
\titlespacing*{\subsection}{0pt}{2.0ex plus 0.4ex}{0.5ex}

% --- TABLA DE CONTENIDO -------------------------------------------------
\titlecontents{section}[0em]
  {\addvspace{0.45pc}\sffamily\bfseries\small\color{black!85}}
  {}{}{\titlerule*[0.35pc]{.}\contentspage}
\titlecontents{subsection}[1.4em]
  {\sffamily\footnotesize\color{black!60}}
  {}{}{\titlerule*[0.35pc]{.}\contentspage}

% =====================================================================
%  ENCABEZADO / PIE — editorial limpio (regla fina, sin franjas)
% =====================================================================
\newcommand{\hdrinstituto}{Instituto Técnico Santo Tomás · Zapatoca}
\newcommand{\hdrcapitulo}{}            % se fija en el documento maestro
\pagestyle{fancy}\fancyhf{}
\renewcommand{\headrulewidth}{0pt}
\fancyhead[LE]{\footnotesize\sffamily\color{black!55}\hdrinstituto}
\fancyhead[RO]{\footnotesize\sffamily\color{subjectcolor}\hdrcapitulo}
\fancyhead[LO,RE]{}
\renewcommand{\headrule}{\vspace{2pt}{\color{black!18}\hrule height 0.4pt}}

% --- PIE: número exterior (cuadro de color) + crédito interior + filete ---
\renewcommand{\footrulewidth}{0pt}
\renewcommand{\footrule}{{\color{black!15}\hrule height 0.4pt}\vspace{3pt}}
\newcommand{\pagenumbox}{%
  \tikz[baseline=-3.2pt]{\node[fill=subjectcolor, text=white, inner xsep=7pt,
     inner ysep=3pt, font=\sffamily\bfseries\footnotesize]{\thepage};}}
\fancyfoot[RO,LE]{\pagenumbox}
\fancyfoot[LO,RE]{\footnotesize\sffamily\color{black!45}%
  \hdrcapitulo\ \textcolor{black!22}{·}\ Prof.\ Juan Diego A.\ Prada R.}

% =====================================================================
%  HELPERS
% =====================================================================
% Etiqueta de bloque:  TAG · título   (color del tag configurable)
\newcommand{\blocklabel}[3]{% #1=color  #2=TAG  #3=título
  {\sffamily\bfseries\footnotesize\color{#1}\MakeUppercase{#2}}%
  \ifblank{#3}{}{\,{\color{black!35}\textbf{·}}\,%
    {\sffamily\bfseries\footnotesize\color{black!78}#3}}%
  \par\nopagebreak\vspace{0.4em}}

% --- CAJAS AL MARGEN — personajes históricos e hitos -------------------
% Van en el margen exterior ancho (\marginpar respeta la geometría twoside
% y flota para no chocar). Son contenido semántico: colócalas en el cuerpo,
% cerca del párrafo que comentan, ENTRE párrafos. NO usarlas dentro de
% multicols, tcolorbox, tablas, ni en la misma línea que un \section.
%   \personaje{Nombre}{años · rol}{descripción breve}
%   \hito{año}{qué pasó (breve)}
\setlength{\marginparpush}{8pt}
\newcommand{\margenbox}[1]{\marginpar{%
  \begingroup\raggedright\footnotesize\sffamily
  \setlength{\parskip}{0.18em}\linespread{1.05}\selectfont
  #1\par\endgroup}}
\newcommand{\personaje}[3]{\margenbox{%
  {\color{subjectcolor}\rule{\linewidth}{1.1pt}}\par
  {\InterDisplaySB\scriptsize\color{subjectcolor}#1}\par
  {\fontsize{6.8pt}{8pt}\selectfont\itshape\color{black!50}#2}\par
  \vspace{0.15em}{\scriptsize\color{black!72}#3}}}
\newcommand{\hito}[2]{\margenbox{%
  {\sffamily\bfseries\scriptsize\colorbox{subjectcolor}{\textcolor{white}{\,#1\,}}}\par
  \vspace{0.15em}{\scriptsize\color{black!72}#2}}}

% =====================================================================
%  ENTORNOS — todos planos: filete lateral, sin marco, sin sombra
% =====================================================================

% --- DEFINICIÓN / concepto clave ---------------------------------------
\newtcolorbox{definicion}[1]{
  enhanced, breakable, sharp corners, boxrule=0pt, frame hidden,
  colback=subjectcolor!6!white,
  borderline west={2.2pt}{0pt}{subjectcolor},
  left=5mm, right=4mm, top=3mm, bottom=3mm,
  before skip=\bloqueskip, after skip=\bloqueskip,
  before upper={\blocklabel{subjectcolor}{Definición}{#1}}
}

% --- PROCEDIMIENTO / pasos ---------------------------------------------
\newtcolorbox{procedimiento}[1]{
  enhanced, breakable, sharp corners, boxrule=0pt, frame hidden,
  colback=black!3!white,
  borderline west={2.2pt}{0pt}{black!45},
  left=5mm, right=4mm, top=3mm, bottom=3mm,
  before skip=\bloqueskip, after skip=\bloqueskip,
  before upper={\blocklabel{black!60}{Procedimiento}{#1}}
}
\newcommand{\paso}[2]{%
  \par\vspace{0.45em}\noindent
  \tikz[baseline=-0.6ex]\node[circle, fill=subjectcolor, text=white,
     inner sep=0pt, minimum size=1.45em,
     font=\sffamily\bfseries\footnotesize]{#1};\hspace{0.5em}#2\par}

% Glosa "¿por qué?" bajo un paso — narración en lenguaje cercano
\newcommand{\porque}[1]{\par\vspace{0.1em}\noindent\hspace{2.0em}%
  {\sffamily\footnotesize\itshape\color{subjectcolor}¿por qué?\,}%
  {\sffamily\footnotesize\itshape\color{black!55}#1}\par}

% Operación glosada: cuenta a la izquierda + porqué al margen derecho
% (estilo "comentario al lado del código"). Úsala dentro de ejemplo.
% \opg{ecuación}{porqué}   — primer arg en modo display (mat, cap0)
% \opgt{etiqueta}{porqué}  — primer arg en texto sans bold (física, ciencias)
\newcommand{\opg}[2]{\par\vspace{0.25em}\noindent
  \makebox[0.46\linewidth][l]{$\displaystyle #1$}%
  {\sffamily\footnotesize\itshape\color{black!52}#2}\par}
\newcommand{\opgt}[2]{\par\vspace{0.25em}\noindent
  \makebox[0.46\linewidth][l]{\sffamily\bfseries\color{subjectcolor!80}#1}%
  {\sffamily\footnotesize\itshape\color{black!52}#2}\par}

% --- NOTA / aparte único (sabías · dato · video · margen) --------------
\newtcolorbox{nota}[1][]{
  enhanced, breakable, sharp corners, boxrule=0pt, frame hidden,
  colback=black!4!white,
  borderline west={2.2pt}{0pt}{black!35},
  left=5mm, right=4mm, top=3mm, bottom=3mm,
  before skip=\bloqueskip, after skip=\bloqueskip,
  before upper={\ifblank{#1}{}{\blocklabel{black!55}{Nota}{#1}}}
}

% --- ALERTA / advertencia (única en ámbar, uso escaso) -----------------
\newtcolorbox{alerta}[1]{
  enhanced, breakable, sharp corners, boxrule=0pt, frame hidden,
  colback=alertcolor!7!white,
  borderline west={2.2pt}{0pt}{alertcolor},
  left=5mm, right=4mm, top=3mm, bottom=3mm,
  before skip=\bloqueskip, after skip=\bloqueskip,
  before upper={\blocklabel{alertcolor}{Atención}{#1}}
}
% Sub-ítem de alerta (antes \caso)
\newcommand{\caso}[2]{\par\vspace{0.35em}\noindent
  {\sffamily\bfseries\color{alertcolor!85!black}#1.}\enspace #2\par}

% --- CONTRASTE (la virtud y su exceso → pregunta que cuestiona) --------
%   \contraste{Concepto}{La virtud / lo valioso}{Su exceso / su riesgo}{Pregunta}
\newtcolorbox{contrastebox}{
  enhanced, breakable, sharp corners, boxrule=0pt, frame hidden,
  colback=subjectcolor!5!white,
  borderline west={2.2pt}{0pt}{subjectcolor},
  left=5mm, right=4mm, top=3mm, bottom=3mm,
  before skip=\bloqueskip, after skip=\bloqueskip,
}
\newcommand{\contraste}[4]{%
\begin{contrastebox}
\blocklabel{subjectcolor}{Cuestiona esto}{\faBalanceScale\ #1}
\noindent\begin{minipage}[t]{0.45\linewidth}
  {\sffamily\footnotesize\bfseries\color{subjectcolor!85!black}La virtud}\par
  \vspace{0.2em}#2
\end{minipage}\hfill
{\color{subjectcolor!35}\vrule width 0.4pt}\hfill
\begin{minipage}[t]{0.45\linewidth}
  {\sffamily\footnotesize\bfseries\color{alertcolor!85!black}Su exceso}\par
  \vspace{0.2em}#3
\end{minipage}
\par\vspace{0.6em}
{\color{subjectcolor!25}\titlerule[0.6pt]}\par\vspace{0.4em}
\noindent{\sffamily\itshape\color{subjectcolor}\faArrowRight\ #4}
\end{contrastebox}}

% --- EJEMPLO (sin caja: regla fina + etiqueta numerada) ----------------
\newcounter{ejemplo}[section]
\newenvironment{ejemplo}[1]
  {\par\addvspace{\bloqueskip}\refstepcounter{ejemplo}%
   \noindent{\color{subjectcolor!28}\rule{\linewidth}{0.5pt}}\par\vspace{0.4em}%
   \noindent{\sffamily\bfseries\color{subjectcolor}Ejemplo \theejemplo.}%
   \enspace{\sffamily\bfseries\color{black!78}#1}\par\vspace{0.3em}}
  {\par\addvspace{\bloqueskip}}

% --- PRÁCTICA (ejercicios / hoja / ICFES) ------------------------------
\newtcolorbox{practica}[1]{
  enhanced, breakable, sharp corners, boxrule=0pt, frame hidden,
  colback=white,
  borderline north={1.4pt}{0pt}{subjectcolor},
  left=0mm, right=0mm, top=3mm, bottom=2mm,
  before skip=\bloqueskip, after skip=\bloqueskip,
  before upper={\blocklabel{subjectcolor}{Práctica}{#1}\raggedright}
}

% Ejercicio numerado con nivel opcional [básico|intermedio|reto]
\newcounter{ejer}[section]
\newcommand{\ejer}[2][]{%
  \stepcounter{ejer}\par\vspace{0.5em}\noindent
  {\sffamily\bfseries\color{subjectcolor}\theejer.}\enspace #2%
  \ifblank{#1}{}{\ \textcolor{black!45}{\footnotesize[\textit{#1}]}}\par}

% Niveles = un solo color a tres intensidades (NO nuevos colores)
\newcommand{\nivel}[2]{\par\smallskip\noindent
  \textcolor{subjectcolor!#1}{\small\faCircle}\enspace
  {\sffamily\bfseries\footnotesize\color{black!70}#2}\par}
\newcommand{\basico}{\nivel{40}{Básico}}
\newcommand{\intermedio}{\nivel{70}{Intermedio}}
\newcommand{\reto}{\nivel{100}{Reto}}

% Selección múltiple tipo ICFES — opciones con burbuja (estilo Saber)
\newcounter{pregicfes}[section]
\newcommand{\bubbleletter}[1]{\tikz[baseline=(bc.base)]\node[circle,
  draw=subjectcolor!55, line width=0.6pt, inner sep=1.5pt,
  font=\sffamily\bfseries\scriptsize, text=subjectcolor!90!black](bc){#1};}
\newcommand{\icfes}[5]{%
  \stepcounter{pregicfes}\par\vspace{0.6em}\noindent
  {\sffamily\bfseries\color{subjectcolor}\thepregicfes.}\enspace #1\par
  \begin{enumerate}[label={\protect\bubbleletter{\Alph*}}, leftmargin=1.5cm,
                    itemsep=3pt, topsep=3pt]
    \item #2 \item #3 \item #4 \item #5
  \end{enumerate}}
% Pregunta abierta de argumentación
\newcommand{\pregabierta}[1]{\par\vspace{0.6em}\noindent
  {\sffamily\bfseries\color{subjectcolor}\faComment\ Argumenta:}\enspace
  \textit{#1}\par\lineaescritura\lineaescritura\lineaescritura}

% --- LÍNEAS DE RESPUESTA -----------------------------------------------
\newcommand{\linea}[1][7cm]{\underline{\hspace{#1}}}
\newcommand{\lineaescritura}{\par\vspace{0.9em}\noindent
  {\color{black!30}\rule{\linewidth}{0.3pt}}}

% =====================================================================
%  EVALUACIÓN TIPO ICFES — cierre de unidad (banner + 2 columnas)
%  Uso:  \begin{evaluacion} \begin{multicols}{2} \icfes... \end{multicols}
%        \pregabierta{...} \end{evaluacion}
% =====================================================================
\newtcolorbox{evaluacion}[1][cierre de sección]{
  enhanced, breakable, sharp corners,
  colback=white, colframe=subjectcolor, boxrule=1pt,
  colbacktitle=subjectcolor, coltitle=white,
  fonttitle=\InterDisplaySB\normalsize,
  title={\faClipboardCheck\ \ EVALUACIÓN TIPO ICFES\ \ %
         \textnormal{\sffamily\small· #1}},
  toptitle=2mm, bottomtitle=2mm, left=4mm, right=4mm, top=3.5mm, bottom=3mm,
  before={\clearpage},
  before upper={\small\itshape\color{black!55}%
    Selección múltiple. Marca una sola opción por pregunta.\par\medskip\raggedright}
}

% =====================================================================
%  HOJA DE TRABAJO — recortable: borde punteado, escala de grises,
%  en página aparte (clearpage antes; el contenido que sigue continúa).
% =====================================================================
\newtcolorbox{hoja}[1]{
  enhanced, breakable, sharp corners,
  colback=black!2!white, boxrule=0pt, frame hidden,
  borderline={0.8pt}{0pt}{black!50, dash pattern=on 3pt off 2.5pt},
  before={\clearpage}, after={\par},
  left=6mm, right=6mm, top=6mm, bottom=6mm,
  before upper={%
    {\sffamily\bfseries\footnotesize\color{black!60}\faCut\ \ HOJA DE TRABAJO}%
    {\color{black!35}\ \textbf{·}\ }{\sffamily\bfseries\footnotesize\color{black!80}#1}%
    \hfill{\sffamily\footnotesize\color{black!55}%
      Nombre:\ \underline{\hspace{4.2cm}}\quad Fecha:\ \underline{\hspace{2cm}}}\par
    \vspace{2mm}{\color{black!30}\hrule height 0.4pt}\vspace{3mm}\raggedright}
}

% =====================================================================
%  TABLAHOJA — tabla con TODAS las líneas (rejilla) para hojas de trabajo.
%  Celdas altas para llenar a mano. Centrada. El encabezado se marca con
%  \rowcolor en la primera fila.
%  Uso:  \begin{tablahoja}{ccc} enc & enc & enc \\ a & b & c \\ \end{tablahoja}
% =====================================================================
\NewDocumentEnvironment{tablahoja}{m}{%
  \par\smallskip\centering\small\setlength{\tabcolsep}{4pt}%
  \begin{NiceTabular}{#1}[hvlines, cell-space-top-limit=5pt,
     cell-space-bottom-limit=9pt, colortbl-like]%
}{%
  \end{NiceTabular}\par\smallskip
}

% =====================================================================
%  FIGURAS — imagen real o placeholder automático (\IfFileExists)
% =====================================================================
\newcounter{figura}[section]
\newcommand{\figura}[3][0.7\linewidth]{% [ancho]{archivo}{pie}
  \par\addvspace{\bloqueskip}\refstepcounter{figura}%
  \begin{center}%
  \IfFileExists{assets/imgs/#2}{%
    \includegraphics[width=#1]{assets/imgs/#2}}{%
    \begin{tikzpicture}
      \node[draw=black!30, dashed, line width=0.6pt, fill=black!3,
            text width=0.66\linewidth, minimum height=2.6cm,
            align=center, inner sep=8pt,
            text=black!55, font=\sffamily\footnotesize]
        {{\Large\faImage}\\[4pt]\textbf{[ imagen pendiente ]}\\#3};
    \end{tikzpicture}}%
  \par\vspace{0.4em}%
  {\footnotesize\sffamily\color{black!60}%
    \textbf{Figura \thefigura.}\enspace #3}%
  \end{center}\par\addvspace{\bloqueskip}}

% =====================================================================
%  APERTURA DE SECCIÓN — marco branded de serie + título + "qué verás"
%  El "marco" es CONSTANTE en todas las secciones (= identidad de marca):
%  número grande de serie, kicker, cuadro-firma, franja de acento y sello
%  del instituto. El arte (imagen real o placeholder) va DENTRO del marco.
%  Uso (1ª línea de cada sNN.tex):
%    \aperturaseccion{ruta/img.png}{Título}{Una o dos frases.}
%  Personaliza por documento (opcional):  \renewcommand{\seccionkicker}{...}
% =====================================================================
\newcounter{secportada}
\newcounter{unidadnum}                       % nº de unidad actual (lo fija \aperturacapitulo)
% Número visible de sección = Unidad.Sección (ej. 1.2). Si no hay unidad
% definida (unidadnum=0), cae al número simple de sección.
\newcommand{\numseccion}{\ifnum\value{unidadnum}>0 \theunidadnum.\thesecportada\else\thesecportada\fi}
\newcommand{\seccionkicker}{\hdrcapitulo}   % por defecto, lo del encabezado

% El escenario branded (contenedor constante del arte)
\newtcolorbox{secstage}{%
  enhanced, sharp corners, boxrule=0pt, frame hidden,
  colback=subjectcolor!6, halign=center, valign=center,
  height=5.0cm, left=2mm, right=2mm, top=10mm, bottom=8mm,
  overlay={%
    % cuadro-firma + kicker (arriba-izquierda)
    \fill[subjectcolor] ([shift={(4mm,-4mm)}]frame.north west) rectangle ++(2.6mm,-2.6mm);
    \node[anchor=north west, font=\InterDisplaySB\footnotesize, text=subjectcolor!85]
      at ([shift={(8.5mm,-3.4mm)}]frame.north west) {\seccionkicker};
    % número de serie grande y tenue (arriba-derecha) = marca
    \node[anchor=north east, font=\InterDisplayBlk, text=subjectcolor!18, inner sep=0pt]
      at ([shift={(-4mm,-1.5mm)}]frame.north east)
      {\fontsize{34}{34}\selectfont \numseccion};
    % franja de acento (abajo-izquierda)
    \fill[subjectcolor] ([shift={(4mm,4mm)}]frame.south west) rectangle ++(2.8cm,2.2mm);
    % sello del instituto (abajo-derecha)
    \node[anchor=south east, font=\sffamily\scriptsize, text=subjectcolor!55]
      at ([shift={(-4mm,3.4mm)}]frame.south east) {ITSTZ · Zapatoca};
  }}

\newcommand{\aperturaseccion}[3]{% {archivo}{título}{intro}
  \clearpage\refstepcounter{secportada}%
  \begin{secstage}
  \IfFileExists{assets/imgs/#1}{%
    \includegraphics[width=0.9\linewidth, height=3.2cm, keepaspectratio]{assets/imgs/#1}}{%
    \begin{tikzpicture}
      \node[draw=subjectcolor!35, dashed, line width=0.6pt, fill=white,
            minimum width=0.86\linewidth, minimum height=2.9cm,
            align=center, text=subjectcolor!50, font=\sffamily\small]
        {{\fontsize{22}{22}\selectfont\faImage}\\[5pt]
         \textbf{[ imagen pendiente ]}\\[1pt]
         \scriptsize\ttfamily assets/imgs/#1};
    \end{tikzpicture}}%
  \end{secstage}
  \vspace{0.5em}
  {\color{subjectcolor}\rule{3.2cm}{2.5pt}}\par\vspace{0.3em}
  \section{{\color{subjectcolor}\numseccion}\enspace #2}%
  {\InterLight\fontsize{12.5}{17}\selectfont\color{black!58} #3\par}
  \vspace{0.4em}}

% =====================================================================
%  DESTACADO — "lo esencial": triángulo lateral + degradado muy sutil
% =====================================================================
\newtcolorbox{destacado}[1]{
  enhanced, breakable, sharp corners, boxrule=0pt, frame hidden,
  interior style={left color=subjectcolor!16, right color=subjectcolor!2},
  borderline west={2.4pt}{0pt}{subjectcolor},
  left=7mm, right=5mm, top=3.5mm, bottom=3.5mm,
  before skip=\bloqueskip, after skip=\bloqueskip,
  overlay unbroken and first={%
    \fill[subjectcolor] ([yshift=-4.5mm]frame.north west)
        -- ++(5.5pt,2.75pt) -- ++(0,-5.5pt) -- cycle;},
  before upper={\blocklabel{subjectcolor}{Lo esencial}{#1}}
}

% --- Resaltado inline suave ---
\newtcbox{\resaltar}{on line, nobeforeafter, boxsep=0pt,
  left=2.5pt, right=2.5pt, top=1pt, bottom=1pt, arc=2.5pt,
  colback=subjectcolor!15, colframe=subjectcolor!15}

% =====================================================================
%  EXTENSIONES SANTILLANA (opt-in) — cajas temáticas para dar variedad
%  visual y romper la monotonía. NO alteran los entornos base; úsalas
%  donde aporten. Mantienen el color de asignatura (coherencia de marca).
% =====================================================================

% --- ¿SABÍAS QUE? — curiosidad breve · tarjeta redondeada con bombillo --
\newtcolorbox{sabias}{
  enhanced, breakable, arc=2.5mm, boxrule=0.8pt,
  colback=subjectcolor!6!white, colframe=subjectcolor!30,
  left=4mm, right=4mm, top=3mm, bottom=3mm,
  before skip=\bloqueskip, after skip=\bloqueskip,
  before upper={{\sffamily\bfseries\footnotesize\color{subjectcolor}%
    \faLightbulb\enspace\MakeUppercase{¿Sabías que?}}%
    \par\nopagebreak\vspace{0.45em}}
}

% --- CIENCIA Y SOCIEDAD — impacto/ética · banner de título -------------
\newtcolorbox{cienciasociedad}[1]{
  enhanced, breakable, sharp corners,
  colback=subjectcolor!5!white, colframe=subjectcolor, boxrule=0pt,
  colbacktitle=subjectcolor, coltitle=white,
  fonttitle=\sffamily\bfseries\footnotesize,
  title={\faGlobeAmericas\ \ \MakeUppercase{Ciencia y sociedad}%
         \ifblank{#1}{}{\ \ \textnormal{\sffamily\small· #1}}},
  toptitle=1.6mm, bottomtitle=1.6mm,
  left=4mm, right=4mm, top=3mm, bottom=3mm,
  before skip=\bloqueskip, after skip=\bloqueskip,
}

% --- MANOS A LA OBRA — práctica experimental · marco sólido con matraz --
\newtcolorbox{laboratorio}[1]{
  enhanced, breakable, sharp corners, arc=0pt,
  colback=white, colframe=subjectcolor!55, boxrule=1pt,
  left=4mm, right=4mm, top=3mm, bottom=3mm,
  before skip=\bloqueskip, after skip=\bloqueskip,
  before upper={{\sffamily\bfseries\footnotesize\color{subjectcolor}%
    \faFlask\enspace\MakeUppercase{Manos a la obra}}%
    \ifblank{#1}{}{\ {\color{black!35}\textbf{·}}\ %
      {\sffamily\bfseries\footnotesize\color{black!78}#1}}%
    \par\nopagebreak\vspace{0.45em}}
}
% Lista de materiales dentro de \begin{laboratorio}
\newcommand{\materiales}[1]{\par\noindent
  {\sffamily\bfseries\footnotesize\color{black!70}Materiales:}\ %
  {\footnotesize #1}\par\vspace{0.3em}}

% --- IMAGEN INLINE con texto alrededor (Santillana) -------------------
%   \figinline[r|l]{ancho}{archivo}{pie}   — colócala al inicio de un párrafo
\newcommand{\figinline}[4][r]{%
  \begin{wrapfigure}{#1}{#2}
  \centering\vspace{-0.5\baselineskip}
  \IfFileExists{assets/imgs/#3}{%
    \includegraphics[width=\linewidth]{assets/imgs/#3}}{%
    \fbox{\parbox[c][3cm][c]{0.92\linewidth}{\centering\sffamily\scriptsize
      \color{black!55}{\Large\faImage}\\[3pt][imagen pendiente]\\\ttfamily\tiny #3}}}%
  \par\vspace{2pt}{\footnotesize\sffamily\color{black!60}\textbf{Fig.}\enspace #4}%
  \vspace{-0.4\baselineskip}
  \end{wrapfigure}}

% --- APERTURA CON HERO LATERAL (imagen grande al lado del título) ------
%   \aperturahero{archivo}{título}{intro}  — alternativa a \aperturaseccion
\newcommand{\aperturahero}[3]{%
  \clearpage\refstepcounter{secportada}\stepcounter{section}%
  \addcontentsline{toc}{section}{\numseccion\enspace #2}%
  \noindent\begin{minipage}[c]{0.62\linewidth}
    {\InterDisplaySB\footnotesize\color{subjectcolor!85}\MakeUppercase{\seccionkicker}}\par
    \vspace{1.8mm}
    {\LARGE\InterDisplaySB\color{subjectcolor}%
      {\color{subjectcolor!35}\numseccion}\,\,#2}\par
    \vspace{2mm}{\color{subjectcolor!30}\rule{\linewidth}{0.8pt}}\par\vspace{2mm}
    {\InterLight\fontsize{12}{16}\selectfont\color{black!58}#3}
  \end{minipage}\hfill
  \begin{minipage}[c]{0.34\linewidth}\centering
    \IfFileExists{assets/imgs/#1}{%
      \includegraphics[width=\linewidth]{assets/imgs/#1}}{%
      \fbox{\parbox[c][3.4cm][c]{0.95\linewidth}{\centering\sffamily\scriptsize
        \color{black!55}{\Large\faImage}\\[3pt][imagen pendiente]\\\ttfamily\tiny #1}}}%
  \end{minipage}\par\vspace{1.3em}}

% --- MAPAS (conceptual y mental) — estilos TikZ con texto nítido -------
\tikzset{
  croot/.style={draw=subjectcolor, fill=subjectcolor, text=white, rounded corners=3pt,
    font=\sffamily\bfseries\footnotesize, align=center, inner sep=5pt, minimum height=9mm},
  cnode/.style={draw=subjectcolor!55, fill=subjectcolor!8, rounded corners=2.5pt,
    font=\sffamily\footnotesize, align=center, inner sep=4pt, text width=2.5cm},
  cleaf/.style={draw=subjectcolor!45, fill=white, rounded corners=2.5pt,
    font=\sffamily\footnotesize, align=center, inner sep=4pt},
  clink/.style={-{Stealth[length=5pt]}, subjectcolor!65, semithick},
  cline/.style={subjectcolor!55, semithick},
}
% Caja contenedora para un mapa (título + regla superior)
\newtcolorbox{mapabox}[2][Mapa conceptual]{
  enhanced, sharp corners, boxrule=0pt, frame hidden, colback=white,
  borderline north={1.4pt}{0pt}{subjectcolor},
  left=1mm, right=1mm, top=3mm, bottom=2mm,
  before skip=\bloqueskip, after skip=\bloqueskip,
  before upper={\blocklabel{subjectcolor}{#1}{#2}}
}

% --- Papel cuadriculado para graficar a mano (hojas de trabajo) --------
%   \papelcuadricula[filas]{ancho en cm}   (cada celda = 0.6 cm)
\newcommand{\papelcuadricula}[2][8]{%
  \par\smallskip\begin{center}
  \begin{tikzpicture}
    \draw[step=0.6cm, black!16, line width=0.3pt] (0,0) grid (#2,{#1*0.6});
    \draw[black!45, line width=0.7pt] (0,0) -- (#2,0);
    \draw[black!45, line width=0.7pt] (0,0) -- (0,{#1*0.6});
  \end{tikzpicture}
  \end{center}\par\smallskip}

% =====================================================================
%  APERTURA DE CAPÍTULO — título grande, número marca de agua, aire
% =====================================================================
\newcommand{\aperturacapitulo}[3]{% {número}{título}{subtítulo}
  \setcounter{unidadnum}{#1}\setcounter{secportada}{0}% fija unidad y reinicia secciones
  \thispagestyle{empty}%
  \begin{tikzpicture}[remember picture, overlay]
    \node[anchor=north east, text=subjectcolor!9,
          font=\InterDisplayThin]
      at ([xshift=-0.9cm, yshift=-2.2cm]current page.north east)
      {\fontsize{210}{210}\selectfont #1};
  \end{tikzpicture}%
  \begingroup\raggedright\null\vspace{3.8cm}\par
  {\InterDisplaySB\fontsize{12}{14}\selectfont\color{subjectcolor}%
    CAP\'ITULO\ #1}\par
  \vspace{0.18cm}{\color{black!16}\rule{4cm}{1pt}}\par
  \vspace{0.9cm}
  {\InterDisplaySB\fontsize{46}{50}\selectfont\color{black!88}#2}\par
  \vspace{0.5cm}{\color{subjectcolor}\rule{3.2cm}{2.5pt}}\par
  \vspace{0.55cm}
  {\InterLight\fontsize{15}{21}\selectfont\color{black!55}#3}\par
  \vfill
  {\sffamily\footnotesize\color{black!45}\hdrinstituto\quad
    \textcolor{black!22}{·}\quad Prof.\ Juan Diego A.\ Prada R.}\par
  \endgroup\clearpage}

% =====================================================================
%  DESBLOQUEA — pie discreto, no caja
% =====================================================================
\newcommand{\desbloquea}[1]{%
  \par\addvspace{\bloqueskip}%
  \noindent{\color{subjectcolor!25}\rule{\linewidth}{0.5pt}}\par\vspace{0.4em}%
  {\footnotesize\sffamily\color{black!60}%
    \textcolor{subjectcolor}{\faUnlock}\ \textbf{Dominar esto te desbloquea:}\ #1}\par}
):
%     % =====================================================================
%  TEMA · CIENCIAS NATURALES — variante del design system (GUÍAS)
%  Identidad visual propia, distinta de Matemáticas (azul/pizarra):
%    · color base ESMERALDA (\setsubject{cnat})
%    · acento LIMA orgánico solo para ornamentos (no rompe la regla de
%      "un solo color de asignatura")
%    · iconografía de ciencias (hoja · ADN · matraz · átomo)
%    · reutiliza las cajas de ciencia que ya trae el preamble:
%      \begin{laboratorio}, \begin{sabias}, \begin{cienciasociedad}
%
%  Sirve a futuro para BIOLOGÍA, QUÍMICA y FÍSICA EXPERIMENTAL.
%
%  INVOCACIÓN (una sola línea, tras \input{design/preamble.tex}):
%     \input{design/tema-cnat.tex}
%  El doc puede seguir personalizando \hdrcapitulo, \seccionkicker, etc.
% =====================================================================

% --- Color base de la asignatura --------------------------------------
\setsubject{cnat}                 % esmeralda RGB(16,110,80)

% Acento LIMA orgánico — SOLO ornamentos (hoja, filete). Nunca contenido:
% la regla "un solo color de asignatura" se mantiene intacta.
\definecolor{cnatlima}{RGB}{124,169,60}

% --- Iconografía semántica de ciencias (fontawesome5, ya cargado) -----
\newcommand{\icobio}{{\color{subjectcolor}\faLeaf}}     % biología
\newcommand{\icoquim}{{\color{subjectcolor}\faFlask}}   % química
\newcommand{\icofis}{{\color{subjectcolor}\faAtom}}     % física
\newcommand{\icogen}{{\color{subjectcolor}\faDna}}      % genética

% --- Kicker temático por defecto (el documento puede renovarlo) -------
\renewcommand{\seccionkicker}{\faLeaf\enspace Ciencias Naturales}

% --- Nombre científico: género + especie en itálica + nombre común ----
%   \especie{Pisum}{sativum}{arveja}  ->  *Pisum sativum* (arveja)
\newcommand{\especie}[3]{\textit{#1 #2}\ (#3)}

% --- Genotipo en línea: preserva mayúsculas/minúsculas de los alelos ---
%   \genotipo{Aa}  (NUNCA aplicar \MakeUppercase: borraría dominante/recesivo)
\newcommand{\genotipo}[1]{{\sffamily\bfseries\color{subjectcolor}#1}}

% --- Filete orgánico con hoja: cierre decorativo opcional de bloque ----
\newcommand{\filetehoja}{%
  \par\vspace{0.4em}\noindent
  {\color{cnatlima}\faLeaf}\hspace{0.3em}%
  {\color{subjectcolor!35}\rule[0.35ex]{0.86\linewidth}{0.7pt}}\par}

%  El doc puede seguir personalizando \hdrcapitulo, \seccionkicker, etc.
% =====================================================================

% --- Color base de la asignatura --------------------------------------
\setsubject{cnat}                 % esmeralda RGB(16,110,80)

% Acento LIMA orgánico — SOLO ornamentos (hoja, filete). Nunca contenido:
% la regla "un solo color de asignatura" se mantiene intacta.
\definecolor{cnatlima}{RGB}{124,169,60}

% --- Iconografía semántica de ciencias (fontawesome5, ya cargado) -----
\newcommand{\icobio}{{\color{subjectcolor}\faLeaf}}     % biología
\newcommand{\icoquim}{{\color{subjectcolor}\faFlask}}   % química
\newcommand{\icofis}{{\color{subjectcolor}\faAtom}}     % física
\newcommand{\icogen}{{\color{subjectcolor}\faDna}}      % genética

% --- Kicker temático por defecto (el documento puede renovarlo) -------
\renewcommand{\seccionkicker}{\faLeaf\enspace Ciencias Naturales}

% --- Nombre científico: género + especie en itálica + nombre común ----
%   \especie{Pisum}{sativum}{arveja}  ->  *Pisum sativum* (arveja)
\newcommand{\especie}[3]{\textit{#1 #2}\ (#3)}

% --- Genotipo en línea: preserva mayúsculas/minúsculas de los alelos ---
%   \genotipo{Aa}  (NUNCA aplicar \MakeUppercase: borraría dominante/recesivo)
\newcommand{\genotipo}[1]{{\sffamily\bfseries\color{subjectcolor}#1}}

% --- Filete orgánico con hoja: cierre decorativo opcional de bloque ----
\newcommand{\filetehoja}{%
  \par\vspace{0.4em}\noindent
  {\color{cnatlima}\faLeaf}\hspace{0.3em}%
  {\color{subjectcolor!35}\rule[0.35ex]{0.86\linewidth}{0.7pt}}\par}
     % si es ciencias (verde); mate va en azul
%     % =====================================================================
%  TEMA · MODO TABLERO — "clave del profe" a TODO EL ANCHO
%  Rompe a propósito el layout de libro (twoside angosto): página ancha
%  de una sola columna, banner moderno, imágenes libres al lado con el
%  texto rodeándolas (\figinline), gráficos flat TikZ/SVG y TODO RESUELTO.
%
%  Pensado para material que el PROFE resuelve en el tablero (taller/quiz
%  con solución). NO es el estilo Santillana del kit: es intencionalmente
%  más "revista/póster", más aire, más ancho.
%
%  INVOCACIÓN (tras % =====================================================================
%  DESIGN SYSTEM v5 — EDITORIAL ITSTZ
%  Filosofía: texto primero · un solo color de asignatura · sin sombras
%  · sin marcos pesados · tratamientos planos (filete lateral / regla).
%
%  Vocabulario semántico (esto es TODO lo disponible):
%    \section / \subsection      jerarquía
%    definicion{título}          concepto/regla clave (filete color)
%    ejemplo{título}             ejemplo resuelto (numerado, sin caja)
%    procedimiento{título}       pasos (filete neutro) + \paso{n}{texto}
%    nota[título]                aparte único gris (sabías/dato/video/margen)
%    alerta{título}              advertencia ámbar (usar rara vez)
%    practica{título}            ejercicios/hoja/ICFES (regla superior)
%    \figura{archivo}{pie}       imagen real o placeholder automático
%    \desbloquea{lista}          pie discreto de prerrequisitos
%
%  REGLA DE ORO: las secciones (sNN.tex) NO llevan \vspace, color ni
%  minipage de maquetación. Solo contenido semántico. El diseño vive aquí.
% =====================================================================
\documentclass[11pt, letterpaper, twoside]{article}

\usepackage{fontspec}
\usepackage{polyglossia}
\setdefaultlanguage{spanish}
\usepackage[inner=2.1cm, outer=3.7cm, top=2.6cm, bottom=2.3cm,
            marginparwidth=2.55cm, marginparsep=0.45cm,
            headheight=16pt, headsep=18pt, footskip=24pt]{geometry}
\usepackage{microtype}
\usepackage{multicol}
\usepackage{xcolor}
\usepackage{graphicx}
\usepackage{tikz}
\usetikzlibrary{calc, arrows.meta, positioning}
\usepackage{wrapfig}
\usepackage[most]{tcolorbox}
\usepackage{amsmath, amssymb}
\usepackage{siunitx}
\usepackage{fontawesome5}
\usepackage{fancyhdr}
\usepackage{titlesec}
\usepackage{enumitem}
\usepackage{etoolbox}
\usepackage{titletoc}
\usepackage{array}
\usepackage{booktabs}
\usepackage{colortbl}
\usepackage{nicematrix}   % rejilla completa (hvlines) para tablas de hojas

% --- TIPOGRAFÍA -------------------------------------------------------
% Cuerpo: Charter (serif cálida y muy legible). La matemática (Latin
% Modern, serif) combina con ella. Títulos/etiquetas/UI: Inter Display.
\setmainfont{Charter}[
  UprightFont    = {Charter},
  BoldFont       = {Charter Bold},
  ItalicFont     = {Charter Italic},
  BoldItalicFont = {Charter Bold Italic},
  Ligatures = TeX, Numbers = Lining ]
\setsansfont{Inter}[
  UprightFont = Inter-Regular, BoldFont = Inter-SemiBold,
  ItalicFont  = Inter-Italic, Ligatures = TeX, Numbers = Lining ]
\newfontfamily\InterLight{Inter-Light}[ItalicFont=Inter-LightItalic, Ligatures=TeX]
\newfontfamily\InterDisplay{Inter Display}[Ligatures=TeX]
\newfontfamily\InterDisplaySB{Inter Display SemiBold}[Ligatures=TeX]
\newfontfamily\InterDisplayThin{Inter Display Thin}[Ligatures=TeX]
\newfontfamily\InterDisplayBlk{Inter Display Black}[Ligatures=TeX]

% --- RITMO --------------------------------------------------------------
\linespread{1.22}
\setlength{\parindent}{0pt}
\setlength{\parskip}{0.5em}
\newlength{\bloqueskip}\setlength{\bloqueskip}{1.0em}   % aire entre bloques
\setlength{\columnsep}{1.6em}                            % medianil multicols

% =====================================================================
%  COLOR — UNA sola variable de asignatura + neutros. Sin arcoíris.
% =====================================================================
\definecolor{mat} {RGB}{ 30,  64, 120}   % azul profundo
\definecolor{cnat}{RGB}{ 16, 110,  80}   % esmeralda
\definecolor{csoc}{RGB}{165,  75,  45}   % terracota
\definecolor{eco} {RGB}{ 70,  55, 130}   % índigo
\definecolor{alertcolor}{RGB}{180, 95, 10}  % ámbar — solo advertencias

\colorlet{subjectcolor}{mat}             % por defecto
\newcommand{\setsubject}[1]{\colorlet{subjectcolor}{#1}}

% =====================================================================
%  JERARQUÍA TIPOGRÁFICA
% =====================================================================
\setcounter{secnumdepth}{0}
\titleformat{\section}
  {\LARGE\InterDisplaySB\color{subjectcolor}}{}{0pt}{}
  [\vspace{2pt}{\color{subjectcolor!30}\titlerule[0.8pt]}]
\titlespacing*{\section}{0pt}{3.2ex plus 0.6ex}{1.8ex plus 0.3ex}

\titleformat{\subsection}
  {\large\InterDisplaySB\color{black!82}}{}{0pt}{}
\titlespacing*{\subsection}{0pt}{2.0ex plus 0.4ex}{0.5ex}

% --- TABLA DE CONTENIDO -------------------------------------------------
\titlecontents{section}[0em]
  {\addvspace{0.45pc}\sffamily\bfseries\small\color{black!85}}
  {}{}{\titlerule*[0.35pc]{.}\contentspage}
\titlecontents{subsection}[1.4em]
  {\sffamily\footnotesize\color{black!60}}
  {}{}{\titlerule*[0.35pc]{.}\contentspage}

% =====================================================================
%  ENCABEZADO / PIE — editorial limpio (regla fina, sin franjas)
% =====================================================================
\newcommand{\hdrinstituto}{Instituto Técnico Santo Tomás · Zapatoca}
\newcommand{\hdrcapitulo}{}            % se fija en el documento maestro
\pagestyle{fancy}\fancyhf{}
\renewcommand{\headrulewidth}{0pt}
\fancyhead[LE]{\footnotesize\sffamily\color{black!55}\hdrinstituto}
\fancyhead[RO]{\footnotesize\sffamily\color{subjectcolor}\hdrcapitulo}
\fancyhead[LO,RE]{}
\renewcommand{\headrule}{\vspace{2pt}{\color{black!18}\hrule height 0.4pt}}

% --- PIE: número exterior (cuadro de color) + crédito interior + filete ---
\renewcommand{\footrulewidth}{0pt}
\renewcommand{\footrule}{{\color{black!15}\hrule height 0.4pt}\vspace{3pt}}
\newcommand{\pagenumbox}{%
  \tikz[baseline=-3.2pt]{\node[fill=subjectcolor, text=white, inner xsep=7pt,
     inner ysep=3pt, font=\sffamily\bfseries\footnotesize]{\thepage};}}
\fancyfoot[RO,LE]{\pagenumbox}
\fancyfoot[LO,RE]{\footnotesize\sffamily\color{black!45}%
  \hdrcapitulo\ \textcolor{black!22}{·}\ Prof.\ Juan Diego A.\ Prada R.}

% =====================================================================
%  HELPERS
% =====================================================================
% Etiqueta de bloque:  TAG · título   (color del tag configurable)
\newcommand{\blocklabel}[3]{% #1=color  #2=TAG  #3=título
  {\sffamily\bfseries\footnotesize\color{#1}\MakeUppercase{#2}}%
  \ifblank{#3}{}{\,{\color{black!35}\textbf{·}}\,%
    {\sffamily\bfseries\footnotesize\color{black!78}#3}}%
  \par\nopagebreak\vspace{0.4em}}

% --- CAJAS AL MARGEN — personajes históricos e hitos -------------------
% Van en el margen exterior ancho (\marginpar respeta la geometría twoside
% y flota para no chocar). Son contenido semántico: colócalas en el cuerpo,
% cerca del párrafo que comentan, ENTRE párrafos. NO usarlas dentro de
% multicols, tcolorbox, tablas, ni en la misma línea que un \section.
%   \personaje{Nombre}{años · rol}{descripción breve}
%   \hito{año}{qué pasó (breve)}
\setlength{\marginparpush}{8pt}
\newcommand{\margenbox}[1]{\marginpar{%
  \begingroup\raggedright\footnotesize\sffamily
  \setlength{\parskip}{0.18em}\linespread{1.05}\selectfont
  #1\par\endgroup}}
\newcommand{\personaje}[3]{\margenbox{%
  {\color{subjectcolor}\rule{\linewidth}{1.1pt}}\par
  {\InterDisplaySB\scriptsize\color{subjectcolor}#1}\par
  {\fontsize{6.8pt}{8pt}\selectfont\itshape\color{black!50}#2}\par
  \vspace{0.15em}{\scriptsize\color{black!72}#3}}}
\newcommand{\hito}[2]{\margenbox{%
  {\sffamily\bfseries\scriptsize\colorbox{subjectcolor}{\textcolor{white}{\,#1\,}}}\par
  \vspace{0.15em}{\scriptsize\color{black!72}#2}}}

% =====================================================================
%  ENTORNOS — todos planos: filete lateral, sin marco, sin sombra
% =====================================================================

% --- DEFINICIÓN / concepto clave ---------------------------------------
\newtcolorbox{definicion}[1]{
  enhanced, breakable, sharp corners, boxrule=0pt, frame hidden,
  colback=subjectcolor!6!white,
  borderline west={2.2pt}{0pt}{subjectcolor},
  left=5mm, right=4mm, top=3mm, bottom=3mm,
  before skip=\bloqueskip, after skip=\bloqueskip,
  before upper={\blocklabel{subjectcolor}{Definición}{#1}}
}

% --- PROCEDIMIENTO / pasos ---------------------------------------------
\newtcolorbox{procedimiento}[1]{
  enhanced, breakable, sharp corners, boxrule=0pt, frame hidden,
  colback=black!3!white,
  borderline west={2.2pt}{0pt}{black!45},
  left=5mm, right=4mm, top=3mm, bottom=3mm,
  before skip=\bloqueskip, after skip=\bloqueskip,
  before upper={\blocklabel{black!60}{Procedimiento}{#1}}
}
\newcommand{\paso}[2]{%
  \par\vspace{0.45em}\noindent
  \tikz[baseline=-0.6ex]\node[circle, fill=subjectcolor, text=white,
     inner sep=0pt, minimum size=1.45em,
     font=\sffamily\bfseries\footnotesize]{#1};\hspace{0.5em}#2\par}

% Glosa "¿por qué?" bajo un paso — narración en lenguaje cercano
\newcommand{\porque}[1]{\par\vspace{0.1em}\noindent\hspace{2.0em}%
  {\sffamily\footnotesize\itshape\color{subjectcolor}¿por qué?\,}%
  {\sffamily\footnotesize\itshape\color{black!55}#1}\par}

% Operación glosada: cuenta a la izquierda + porqué al margen derecho
% (estilo "comentario al lado del código"). Úsala dentro de ejemplo.
% \opg{ecuación}{porqué}   — primer arg en modo display (mat, cap0)
% \opgt{etiqueta}{porqué}  — primer arg en texto sans bold (física, ciencias)
\newcommand{\opg}[2]{\par\vspace{0.25em}\noindent
  \makebox[0.46\linewidth][l]{$\displaystyle #1$}%
  {\sffamily\footnotesize\itshape\color{black!52}#2}\par}
\newcommand{\opgt}[2]{\par\vspace{0.25em}\noindent
  \makebox[0.46\linewidth][l]{\sffamily\bfseries\color{subjectcolor!80}#1}%
  {\sffamily\footnotesize\itshape\color{black!52}#2}\par}

% --- NOTA / aparte único (sabías · dato · video · margen) --------------
\newtcolorbox{nota}[1][]{
  enhanced, breakable, sharp corners, boxrule=0pt, frame hidden,
  colback=black!4!white,
  borderline west={2.2pt}{0pt}{black!35},
  left=5mm, right=4mm, top=3mm, bottom=3mm,
  before skip=\bloqueskip, after skip=\bloqueskip,
  before upper={\ifblank{#1}{}{\blocklabel{black!55}{Nota}{#1}}}
}

% --- ALERTA / advertencia (única en ámbar, uso escaso) -----------------
\newtcolorbox{alerta}[1]{
  enhanced, breakable, sharp corners, boxrule=0pt, frame hidden,
  colback=alertcolor!7!white,
  borderline west={2.2pt}{0pt}{alertcolor},
  left=5mm, right=4mm, top=3mm, bottom=3mm,
  before skip=\bloqueskip, after skip=\bloqueskip,
  before upper={\blocklabel{alertcolor}{Atención}{#1}}
}
% Sub-ítem de alerta (antes \caso)
\newcommand{\caso}[2]{\par\vspace{0.35em}\noindent
  {\sffamily\bfseries\color{alertcolor!85!black}#1.}\enspace #2\par}

% --- CONTRASTE (la virtud y su exceso → pregunta que cuestiona) --------
%   \contraste{Concepto}{La virtud / lo valioso}{Su exceso / su riesgo}{Pregunta}
\newtcolorbox{contrastebox}{
  enhanced, breakable, sharp corners, boxrule=0pt, frame hidden,
  colback=subjectcolor!5!white,
  borderline west={2.2pt}{0pt}{subjectcolor},
  left=5mm, right=4mm, top=3mm, bottom=3mm,
  before skip=\bloqueskip, after skip=\bloqueskip,
}
\newcommand{\contraste}[4]{%
\begin{contrastebox}
\blocklabel{subjectcolor}{Cuestiona esto}{\faBalanceScale\ #1}
\noindent\begin{minipage}[t]{0.45\linewidth}
  {\sffamily\footnotesize\bfseries\color{subjectcolor!85!black}La virtud}\par
  \vspace{0.2em}#2
\end{minipage}\hfill
{\color{subjectcolor!35}\vrule width 0.4pt}\hfill
\begin{minipage}[t]{0.45\linewidth}
  {\sffamily\footnotesize\bfseries\color{alertcolor!85!black}Su exceso}\par
  \vspace{0.2em}#3
\end{minipage}
\par\vspace{0.6em}
{\color{subjectcolor!25}\titlerule[0.6pt]}\par\vspace{0.4em}
\noindent{\sffamily\itshape\color{subjectcolor}\faArrowRight\ #4}
\end{contrastebox}}

% --- EJEMPLO (sin caja: regla fina + etiqueta numerada) ----------------
\newcounter{ejemplo}[section]
\newenvironment{ejemplo}[1]
  {\par\addvspace{\bloqueskip}\refstepcounter{ejemplo}%
   \noindent{\color{subjectcolor!28}\rule{\linewidth}{0.5pt}}\par\vspace{0.4em}%
   \noindent{\sffamily\bfseries\color{subjectcolor}Ejemplo \theejemplo.}%
   \enspace{\sffamily\bfseries\color{black!78}#1}\par\vspace{0.3em}}
  {\par\addvspace{\bloqueskip}}

% --- PRÁCTICA (ejercicios / hoja / ICFES) ------------------------------
\newtcolorbox{practica}[1]{
  enhanced, breakable, sharp corners, boxrule=0pt, frame hidden,
  colback=white,
  borderline north={1.4pt}{0pt}{subjectcolor},
  left=0mm, right=0mm, top=3mm, bottom=2mm,
  before skip=\bloqueskip, after skip=\bloqueskip,
  before upper={\blocklabel{subjectcolor}{Práctica}{#1}\raggedright}
}

% Ejercicio numerado con nivel opcional [básico|intermedio|reto]
\newcounter{ejer}[section]
\newcommand{\ejer}[2][]{%
  \stepcounter{ejer}\par\vspace{0.5em}\noindent
  {\sffamily\bfseries\color{subjectcolor}\theejer.}\enspace #2%
  \ifblank{#1}{}{\ \textcolor{black!45}{\footnotesize[\textit{#1}]}}\par}

% Niveles = un solo color a tres intensidades (NO nuevos colores)
\newcommand{\nivel}[2]{\par\smallskip\noindent
  \textcolor{subjectcolor!#1}{\small\faCircle}\enspace
  {\sffamily\bfseries\footnotesize\color{black!70}#2}\par}
\newcommand{\basico}{\nivel{40}{Básico}}
\newcommand{\intermedio}{\nivel{70}{Intermedio}}
\newcommand{\reto}{\nivel{100}{Reto}}

% Selección múltiple tipo ICFES — opciones con burbuja (estilo Saber)
\newcounter{pregicfes}[section]
\newcommand{\bubbleletter}[1]{\tikz[baseline=(bc.base)]\node[circle,
  draw=subjectcolor!55, line width=0.6pt, inner sep=1.5pt,
  font=\sffamily\bfseries\scriptsize, text=subjectcolor!90!black](bc){#1};}
\newcommand{\icfes}[5]{%
  \stepcounter{pregicfes}\par\vspace{0.6em}\noindent
  {\sffamily\bfseries\color{subjectcolor}\thepregicfes.}\enspace #1\par
  \begin{enumerate}[label={\protect\bubbleletter{\Alph*}}, leftmargin=1.5cm,
                    itemsep=3pt, topsep=3pt]
    \item #2 \item #3 \item #4 \item #5
  \end{enumerate}}
% Pregunta abierta de argumentación
\newcommand{\pregabierta}[1]{\par\vspace{0.6em}\noindent
  {\sffamily\bfseries\color{subjectcolor}\faComment\ Argumenta:}\enspace
  \textit{#1}\par\lineaescritura\lineaescritura\lineaescritura}

% --- LÍNEAS DE RESPUESTA -----------------------------------------------
\newcommand{\linea}[1][7cm]{\underline{\hspace{#1}}}
\newcommand{\lineaescritura}{\par\vspace{0.9em}\noindent
  {\color{black!30}\rule{\linewidth}{0.3pt}}}

% =====================================================================
%  EVALUACIÓN TIPO ICFES — cierre de unidad (banner + 2 columnas)
%  Uso:  \begin{evaluacion} \begin{multicols}{2} \icfes... \end{multicols}
%        \pregabierta{...} \end{evaluacion}
% =====================================================================
\newtcolorbox{evaluacion}[1][cierre de sección]{
  enhanced, breakable, sharp corners,
  colback=white, colframe=subjectcolor, boxrule=1pt,
  colbacktitle=subjectcolor, coltitle=white,
  fonttitle=\InterDisplaySB\normalsize,
  title={\faClipboardCheck\ \ EVALUACIÓN TIPO ICFES\ \ %
         \textnormal{\sffamily\small· #1}},
  toptitle=2mm, bottomtitle=2mm, left=4mm, right=4mm, top=3.5mm, bottom=3mm,
  before={\clearpage},
  before upper={\small\itshape\color{black!55}%
    Selección múltiple. Marca una sola opción por pregunta.\par\medskip\raggedright}
}

% =====================================================================
%  HOJA DE TRABAJO — recortable: borde punteado, escala de grises,
%  en página aparte (clearpage antes; el contenido que sigue continúa).
% =====================================================================
\newtcolorbox{hoja}[1]{
  enhanced, breakable, sharp corners,
  colback=black!2!white, boxrule=0pt, frame hidden,
  borderline={0.8pt}{0pt}{black!50, dash pattern=on 3pt off 2.5pt},
  before={\clearpage}, after={\par},
  left=6mm, right=6mm, top=6mm, bottom=6mm,
  before upper={%
    {\sffamily\bfseries\footnotesize\color{black!60}\faCut\ \ HOJA DE TRABAJO}%
    {\color{black!35}\ \textbf{·}\ }{\sffamily\bfseries\footnotesize\color{black!80}#1}%
    \hfill{\sffamily\footnotesize\color{black!55}%
      Nombre:\ \underline{\hspace{4.2cm}}\quad Fecha:\ \underline{\hspace{2cm}}}\par
    \vspace{2mm}{\color{black!30}\hrule height 0.4pt}\vspace{3mm}\raggedright}
}

% =====================================================================
%  TABLAHOJA — tabla con TODAS las líneas (rejilla) para hojas de trabajo.
%  Celdas altas para llenar a mano. Centrada. El encabezado se marca con
%  \rowcolor en la primera fila.
%  Uso:  \begin{tablahoja}{ccc} enc & enc & enc \\ a & b & c \\ \end{tablahoja}
% =====================================================================
\NewDocumentEnvironment{tablahoja}{m}{%
  \par\smallskip\centering\small\setlength{\tabcolsep}{4pt}%
  \begin{NiceTabular}{#1}[hvlines, cell-space-top-limit=5pt,
     cell-space-bottom-limit=9pt, colortbl-like]%
}{%
  \end{NiceTabular}\par\smallskip
}

% =====================================================================
%  FIGURAS — imagen real o placeholder automático (\IfFileExists)
% =====================================================================
\newcounter{figura}[section]
\newcommand{\figura}[3][0.7\linewidth]{% [ancho]{archivo}{pie}
  \par\addvspace{\bloqueskip}\refstepcounter{figura}%
  \begin{center}%
  \IfFileExists{assets/imgs/#2}{%
    \includegraphics[width=#1]{assets/imgs/#2}}{%
    \begin{tikzpicture}
      \node[draw=black!30, dashed, line width=0.6pt, fill=black!3,
            text width=0.66\linewidth, minimum height=2.6cm,
            align=center, inner sep=8pt,
            text=black!55, font=\sffamily\footnotesize]
        {{\Large\faImage}\\[4pt]\textbf{[ imagen pendiente ]}\\#3};
    \end{tikzpicture}}%
  \par\vspace{0.4em}%
  {\footnotesize\sffamily\color{black!60}%
    \textbf{Figura \thefigura.}\enspace #3}%
  \end{center}\par\addvspace{\bloqueskip}}

% =====================================================================
%  APERTURA DE SECCIÓN — marco branded de serie + título + "qué verás"
%  El "marco" es CONSTANTE en todas las secciones (= identidad de marca):
%  número grande de serie, kicker, cuadro-firma, franja de acento y sello
%  del instituto. El arte (imagen real o placeholder) va DENTRO del marco.
%  Uso (1ª línea de cada sNN.tex):
%    \aperturaseccion{ruta/img.png}{Título}{Una o dos frases.}
%  Personaliza por documento (opcional):  \renewcommand{\seccionkicker}{...}
% =====================================================================
\newcounter{secportada}
\newcounter{unidadnum}                       % nº de unidad actual (lo fija \aperturacapitulo)
% Número visible de sección = Unidad.Sección (ej. 1.2). Si no hay unidad
% definida (unidadnum=0), cae al número simple de sección.
\newcommand{\numseccion}{\ifnum\value{unidadnum}>0 \theunidadnum.\thesecportada\else\thesecportada\fi}
\newcommand{\seccionkicker}{\hdrcapitulo}   % por defecto, lo del encabezado

% El escenario branded (contenedor constante del arte)
\newtcolorbox{secstage}{%
  enhanced, sharp corners, boxrule=0pt, frame hidden,
  colback=subjectcolor!6, halign=center, valign=center,
  height=5.0cm, left=2mm, right=2mm, top=10mm, bottom=8mm,
  overlay={%
    % cuadro-firma + kicker (arriba-izquierda)
    \fill[subjectcolor] ([shift={(4mm,-4mm)}]frame.north west) rectangle ++(2.6mm,-2.6mm);
    \node[anchor=north west, font=\InterDisplaySB\footnotesize, text=subjectcolor!85]
      at ([shift={(8.5mm,-3.4mm)}]frame.north west) {\seccionkicker};
    % número de serie grande y tenue (arriba-derecha) = marca
    \node[anchor=north east, font=\InterDisplayBlk, text=subjectcolor!18, inner sep=0pt]
      at ([shift={(-4mm,-1.5mm)}]frame.north east)
      {\fontsize{34}{34}\selectfont \numseccion};
    % franja de acento (abajo-izquierda)
    \fill[subjectcolor] ([shift={(4mm,4mm)}]frame.south west) rectangle ++(2.8cm,2.2mm);
    % sello del instituto (abajo-derecha)
    \node[anchor=south east, font=\sffamily\scriptsize, text=subjectcolor!55]
      at ([shift={(-4mm,3.4mm)}]frame.south east) {ITSTZ · Zapatoca};
  }}

\newcommand{\aperturaseccion}[3]{% {archivo}{título}{intro}
  \clearpage\refstepcounter{secportada}%
  \begin{secstage}
  \IfFileExists{assets/imgs/#1}{%
    \includegraphics[width=0.9\linewidth, height=3.2cm, keepaspectratio]{assets/imgs/#1}}{%
    \begin{tikzpicture}
      \node[draw=subjectcolor!35, dashed, line width=0.6pt, fill=white,
            minimum width=0.86\linewidth, minimum height=2.9cm,
            align=center, text=subjectcolor!50, font=\sffamily\small]
        {{\fontsize{22}{22}\selectfont\faImage}\\[5pt]
         \textbf{[ imagen pendiente ]}\\[1pt]
         \scriptsize\ttfamily assets/imgs/#1};
    \end{tikzpicture}}%
  \end{secstage}
  \vspace{0.5em}
  {\color{subjectcolor}\rule{3.2cm}{2.5pt}}\par\vspace{0.3em}
  \section{{\color{subjectcolor}\numseccion}\enspace #2}%
  {\InterLight\fontsize{12.5}{17}\selectfont\color{black!58} #3\par}
  \vspace{0.4em}}

% =====================================================================
%  DESTACADO — "lo esencial": triángulo lateral + degradado muy sutil
% =====================================================================
\newtcolorbox{destacado}[1]{
  enhanced, breakable, sharp corners, boxrule=0pt, frame hidden,
  interior style={left color=subjectcolor!16, right color=subjectcolor!2},
  borderline west={2.4pt}{0pt}{subjectcolor},
  left=7mm, right=5mm, top=3.5mm, bottom=3.5mm,
  before skip=\bloqueskip, after skip=\bloqueskip,
  overlay unbroken and first={%
    \fill[subjectcolor] ([yshift=-4.5mm]frame.north west)
        -- ++(5.5pt,2.75pt) -- ++(0,-5.5pt) -- cycle;},
  before upper={\blocklabel{subjectcolor}{Lo esencial}{#1}}
}

% --- Resaltado inline suave ---
\newtcbox{\resaltar}{on line, nobeforeafter, boxsep=0pt,
  left=2.5pt, right=2.5pt, top=1pt, bottom=1pt, arc=2.5pt,
  colback=subjectcolor!15, colframe=subjectcolor!15}

% =====================================================================
%  EXTENSIONES SANTILLANA (opt-in) — cajas temáticas para dar variedad
%  visual y romper la monotonía. NO alteran los entornos base; úsalas
%  donde aporten. Mantienen el color de asignatura (coherencia de marca).
% =====================================================================

% --- ¿SABÍAS QUE? — curiosidad breve · tarjeta redondeada con bombillo --
\newtcolorbox{sabias}{
  enhanced, breakable, arc=2.5mm, boxrule=0.8pt,
  colback=subjectcolor!6!white, colframe=subjectcolor!30,
  left=4mm, right=4mm, top=3mm, bottom=3mm,
  before skip=\bloqueskip, after skip=\bloqueskip,
  before upper={{\sffamily\bfseries\footnotesize\color{subjectcolor}%
    \faLightbulb\enspace\MakeUppercase{¿Sabías que?}}%
    \par\nopagebreak\vspace{0.45em}}
}

% --- CIENCIA Y SOCIEDAD — impacto/ética · banner de título -------------
\newtcolorbox{cienciasociedad}[1]{
  enhanced, breakable, sharp corners,
  colback=subjectcolor!5!white, colframe=subjectcolor, boxrule=0pt,
  colbacktitle=subjectcolor, coltitle=white,
  fonttitle=\sffamily\bfseries\footnotesize,
  title={\faGlobeAmericas\ \ \MakeUppercase{Ciencia y sociedad}%
         \ifblank{#1}{}{\ \ \textnormal{\sffamily\small· #1}}},
  toptitle=1.6mm, bottomtitle=1.6mm,
  left=4mm, right=4mm, top=3mm, bottom=3mm,
  before skip=\bloqueskip, after skip=\bloqueskip,
}

% --- MANOS A LA OBRA — práctica experimental · marco sólido con matraz --
\newtcolorbox{laboratorio}[1]{
  enhanced, breakable, sharp corners, arc=0pt,
  colback=white, colframe=subjectcolor!55, boxrule=1pt,
  left=4mm, right=4mm, top=3mm, bottom=3mm,
  before skip=\bloqueskip, after skip=\bloqueskip,
  before upper={{\sffamily\bfseries\footnotesize\color{subjectcolor}%
    \faFlask\enspace\MakeUppercase{Manos a la obra}}%
    \ifblank{#1}{}{\ {\color{black!35}\textbf{·}}\ %
      {\sffamily\bfseries\footnotesize\color{black!78}#1}}%
    \par\nopagebreak\vspace{0.45em}}
}
% Lista de materiales dentro de \begin{laboratorio}
\newcommand{\materiales}[1]{\par\noindent
  {\sffamily\bfseries\footnotesize\color{black!70}Materiales:}\ %
  {\footnotesize #1}\par\vspace{0.3em}}

% --- IMAGEN INLINE con texto alrededor (Santillana) -------------------
%   \figinline[r|l]{ancho}{archivo}{pie}   — colócala al inicio de un párrafo
\newcommand{\figinline}[4][r]{%
  \begin{wrapfigure}{#1}{#2}
  \centering\vspace{-0.5\baselineskip}
  \IfFileExists{assets/imgs/#3}{%
    \includegraphics[width=\linewidth]{assets/imgs/#3}}{%
    \fbox{\parbox[c][3cm][c]{0.92\linewidth}{\centering\sffamily\scriptsize
      \color{black!55}{\Large\faImage}\\[3pt][imagen pendiente]\\\ttfamily\tiny #3}}}%
  \par\vspace{2pt}{\footnotesize\sffamily\color{black!60}\textbf{Fig.}\enspace #4}%
  \vspace{-0.4\baselineskip}
  \end{wrapfigure}}

% --- APERTURA CON HERO LATERAL (imagen grande al lado del título) ------
%   \aperturahero{archivo}{título}{intro}  — alternativa a \aperturaseccion
\newcommand{\aperturahero}[3]{%
  \clearpage\refstepcounter{secportada}\stepcounter{section}%
  \addcontentsline{toc}{section}{\numseccion\enspace #2}%
  \noindent\begin{minipage}[c]{0.62\linewidth}
    {\InterDisplaySB\footnotesize\color{subjectcolor!85}\MakeUppercase{\seccionkicker}}\par
    \vspace{1.8mm}
    {\LARGE\InterDisplaySB\color{subjectcolor}%
      {\color{subjectcolor!35}\numseccion}\,\,#2}\par
    \vspace{2mm}{\color{subjectcolor!30}\rule{\linewidth}{0.8pt}}\par\vspace{2mm}
    {\InterLight\fontsize{12}{16}\selectfont\color{black!58}#3}
  \end{minipage}\hfill
  \begin{minipage}[c]{0.34\linewidth}\centering
    \IfFileExists{assets/imgs/#1}{%
      \includegraphics[width=\linewidth]{assets/imgs/#1}}{%
      \fbox{\parbox[c][3.4cm][c]{0.95\linewidth}{\centering\sffamily\scriptsize
        \color{black!55}{\Large\faImage}\\[3pt][imagen pendiente]\\\ttfamily\tiny #1}}}%
  \end{minipage}\par\vspace{1.3em}}

% --- MAPAS (conceptual y mental) — estilos TikZ con texto nítido -------
\tikzset{
  croot/.style={draw=subjectcolor, fill=subjectcolor, text=white, rounded corners=3pt,
    font=\sffamily\bfseries\footnotesize, align=center, inner sep=5pt, minimum height=9mm},
  cnode/.style={draw=subjectcolor!55, fill=subjectcolor!8, rounded corners=2.5pt,
    font=\sffamily\footnotesize, align=center, inner sep=4pt, text width=2.5cm},
  cleaf/.style={draw=subjectcolor!45, fill=white, rounded corners=2.5pt,
    font=\sffamily\footnotesize, align=center, inner sep=4pt},
  clink/.style={-{Stealth[length=5pt]}, subjectcolor!65, semithick},
  cline/.style={subjectcolor!55, semithick},
}
% Caja contenedora para un mapa (título + regla superior)
\newtcolorbox{mapabox}[2][Mapa conceptual]{
  enhanced, sharp corners, boxrule=0pt, frame hidden, colback=white,
  borderline north={1.4pt}{0pt}{subjectcolor},
  left=1mm, right=1mm, top=3mm, bottom=2mm,
  before skip=\bloqueskip, after skip=\bloqueskip,
  before upper={\blocklabel{subjectcolor}{#1}{#2}}
}

% --- Papel cuadriculado para graficar a mano (hojas de trabajo) --------
%   \papelcuadricula[filas]{ancho en cm}   (cada celda = 0.6 cm)
\newcommand{\papelcuadricula}[2][8]{%
  \par\smallskip\begin{center}
  \begin{tikzpicture}
    \draw[step=0.6cm, black!16, line width=0.3pt] (0,0) grid (#2,{#1*0.6});
    \draw[black!45, line width=0.7pt] (0,0) -- (#2,0);
    \draw[black!45, line width=0.7pt] (0,0) -- (0,{#1*0.6});
  \end{tikzpicture}
  \end{center}\par\smallskip}

% =====================================================================
%  APERTURA DE CAPÍTULO — título grande, número marca de agua, aire
% =====================================================================
\newcommand{\aperturacapitulo}[3]{% {número}{título}{subtítulo}
  \setcounter{unidadnum}{#1}\setcounter{secportada}{0}% fija unidad y reinicia secciones
  \thispagestyle{empty}%
  \begin{tikzpicture}[remember picture, overlay]
    \node[anchor=north east, text=subjectcolor!9,
          font=\InterDisplayThin]
      at ([xshift=-0.9cm, yshift=-2.2cm]current page.north east)
      {\fontsize{210}{210}\selectfont #1};
  \end{tikzpicture}%
  \begingroup\raggedright\null\vspace{3.8cm}\par
  {\InterDisplaySB\fontsize{12}{14}\selectfont\color{subjectcolor}%
    CAP\'ITULO\ #1}\par
  \vspace{0.18cm}{\color{black!16}\rule{4cm}{1pt}}\par
  \vspace{0.9cm}
  {\InterDisplaySB\fontsize{46}{50}\selectfont\color{black!88}#2}\par
  \vspace{0.5cm}{\color{subjectcolor}\rule{3.2cm}{2.5pt}}\par
  \vspace{0.55cm}
  {\InterLight\fontsize{15}{21}\selectfont\color{black!55}#3}\par
  \vfill
  {\sffamily\footnotesize\color{black!45}\hdrinstituto\quad
    \textcolor{black!22}{·}\quad Prof.\ Juan Diego A.\ Prada R.}\par
  \endgroup\clearpage}

% =====================================================================
%  DESBLOQUEA — pie discreto, no caja
% =====================================================================
\newcommand{\desbloquea}[1]{%
  \par\addvspace{\bloqueskip}%
  \noindent{\color{subjectcolor!25}\rule{\linewidth}{0.5pt}}\par\vspace{0.4em}%
  {\footnotesize\sffamily\color{black!60}%
    \textcolor{subjectcolor}{\faUnlock}\ \textbf{Dominar esto te desbloquea:}\ #1}\par}
; opcional tema de asignatura):
%     % =====================================================================
%  DESIGN SYSTEM v5 — EDITORIAL ITSTZ
%  Filosofía: texto primero · un solo color de asignatura · sin sombras
%  · sin marcos pesados · tratamientos planos (filete lateral / regla).
%
%  Vocabulario semántico (esto es TODO lo disponible):
%    \section / \subsection      jerarquía
%    definicion{título}          concepto/regla clave (filete color)
%    ejemplo{título}             ejemplo resuelto (numerado, sin caja)
%    procedimiento{título}       pasos (filete neutro) + \paso{n}{texto}
%    nota[título]                aparte único gris (sabías/dato/video/margen)
%    alerta{título}              advertencia ámbar (usar rara vez)
%    practica{título}            ejercicios/hoja/ICFES (regla superior)
%    \figura{archivo}{pie}       imagen real o placeholder automático
%    \desbloquea{lista}          pie discreto de prerrequisitos
%
%  REGLA DE ORO: las secciones (sNN.tex) NO llevan \vspace, color ni
%  minipage de maquetación. Solo contenido semántico. El diseño vive aquí.
% =====================================================================
\documentclass[11pt, letterpaper, twoside]{article}

\usepackage{fontspec}
\usepackage{polyglossia}
\setdefaultlanguage{spanish}
\usepackage[inner=2.1cm, outer=3.7cm, top=2.6cm, bottom=2.3cm,
            marginparwidth=2.55cm, marginparsep=0.45cm,
            headheight=16pt, headsep=18pt, footskip=24pt]{geometry}
\usepackage{microtype}
\usepackage{multicol}
\usepackage{xcolor}
\usepackage{graphicx}
\usepackage{tikz}
\usetikzlibrary{calc, arrows.meta, positioning}
\usepackage{wrapfig}
\usepackage[most]{tcolorbox}
\usepackage{amsmath, amssymb}
\usepackage{siunitx}
\usepackage{fontawesome5}
\usepackage{fancyhdr}
\usepackage{titlesec}
\usepackage{enumitem}
\usepackage{etoolbox}
\usepackage{titletoc}
\usepackage{array}
\usepackage{booktabs}
\usepackage{colortbl}
\usepackage{nicematrix}   % rejilla completa (hvlines) para tablas de hojas

% --- TIPOGRAFÍA -------------------------------------------------------
% Cuerpo: Charter (serif cálida y muy legible). La matemática (Latin
% Modern, serif) combina con ella. Títulos/etiquetas/UI: Inter Display.
\setmainfont{Charter}[
  UprightFont    = {Charter},
  BoldFont       = {Charter Bold},
  ItalicFont     = {Charter Italic},
  BoldItalicFont = {Charter Bold Italic},
  Ligatures = TeX, Numbers = Lining ]
\setsansfont{Inter}[
  UprightFont = Inter-Regular, BoldFont = Inter-SemiBold,
  ItalicFont  = Inter-Italic, Ligatures = TeX, Numbers = Lining ]
\newfontfamily\InterLight{Inter-Light}[ItalicFont=Inter-LightItalic, Ligatures=TeX]
\newfontfamily\InterDisplay{Inter Display}[Ligatures=TeX]
\newfontfamily\InterDisplaySB{Inter Display SemiBold}[Ligatures=TeX]
\newfontfamily\InterDisplayThin{Inter Display Thin}[Ligatures=TeX]
\newfontfamily\InterDisplayBlk{Inter Display Black}[Ligatures=TeX]

% --- RITMO --------------------------------------------------------------
\linespread{1.22}
\setlength{\parindent}{0pt}
\setlength{\parskip}{0.5em}
\newlength{\bloqueskip}\setlength{\bloqueskip}{1.0em}   % aire entre bloques
\setlength{\columnsep}{1.6em}                            % medianil multicols

% =====================================================================
%  COLOR — UNA sola variable de asignatura + neutros. Sin arcoíris.
% =====================================================================
\definecolor{mat} {RGB}{ 30,  64, 120}   % azul profundo
\definecolor{cnat}{RGB}{ 16, 110,  80}   % esmeralda
\definecolor{csoc}{RGB}{165,  75,  45}   % terracota
\definecolor{eco} {RGB}{ 70,  55, 130}   % índigo
\definecolor{alertcolor}{RGB}{180, 95, 10}  % ámbar — solo advertencias

\colorlet{subjectcolor}{mat}             % por defecto
\newcommand{\setsubject}[1]{\colorlet{subjectcolor}{#1}}

% =====================================================================
%  JERARQUÍA TIPOGRÁFICA
% =====================================================================
\setcounter{secnumdepth}{0}
\titleformat{\section}
  {\LARGE\InterDisplaySB\color{subjectcolor}}{}{0pt}{}
  [\vspace{2pt}{\color{subjectcolor!30}\titlerule[0.8pt]}]
\titlespacing*{\section}{0pt}{3.2ex plus 0.6ex}{1.8ex plus 0.3ex}

\titleformat{\subsection}
  {\large\InterDisplaySB\color{black!82}}{}{0pt}{}
\titlespacing*{\subsection}{0pt}{2.0ex plus 0.4ex}{0.5ex}

% --- TABLA DE CONTENIDO -------------------------------------------------
\titlecontents{section}[0em]
  {\addvspace{0.45pc}\sffamily\bfseries\small\color{black!85}}
  {}{}{\titlerule*[0.35pc]{.}\contentspage}
\titlecontents{subsection}[1.4em]
  {\sffamily\footnotesize\color{black!60}}
  {}{}{\titlerule*[0.35pc]{.}\contentspage}

% =====================================================================
%  ENCABEZADO / PIE — editorial limpio (regla fina, sin franjas)
% =====================================================================
\newcommand{\hdrinstituto}{Instituto Técnico Santo Tomás · Zapatoca}
\newcommand{\hdrcapitulo}{}            % se fija en el documento maestro
\pagestyle{fancy}\fancyhf{}
\renewcommand{\headrulewidth}{0pt}
\fancyhead[LE]{\footnotesize\sffamily\color{black!55}\hdrinstituto}
\fancyhead[RO]{\footnotesize\sffamily\color{subjectcolor}\hdrcapitulo}
\fancyhead[LO,RE]{}
\renewcommand{\headrule}{\vspace{2pt}{\color{black!18}\hrule height 0.4pt}}

% --- PIE: número exterior (cuadro de color) + crédito interior + filete ---
\renewcommand{\footrulewidth}{0pt}
\renewcommand{\footrule}{{\color{black!15}\hrule height 0.4pt}\vspace{3pt}}
\newcommand{\pagenumbox}{%
  \tikz[baseline=-3.2pt]{\node[fill=subjectcolor, text=white, inner xsep=7pt,
     inner ysep=3pt, font=\sffamily\bfseries\footnotesize]{\thepage};}}
\fancyfoot[RO,LE]{\pagenumbox}
\fancyfoot[LO,RE]{\footnotesize\sffamily\color{black!45}%
  \hdrcapitulo\ \textcolor{black!22}{·}\ Prof.\ Juan Diego A.\ Prada R.}

% =====================================================================
%  HELPERS
% =====================================================================
% Etiqueta de bloque:  TAG · título   (color del tag configurable)
\newcommand{\blocklabel}[3]{% #1=color  #2=TAG  #3=título
  {\sffamily\bfseries\footnotesize\color{#1}\MakeUppercase{#2}}%
  \ifblank{#3}{}{\,{\color{black!35}\textbf{·}}\,%
    {\sffamily\bfseries\footnotesize\color{black!78}#3}}%
  \par\nopagebreak\vspace{0.4em}}

% --- CAJAS AL MARGEN — personajes históricos e hitos -------------------
% Van en el margen exterior ancho (\marginpar respeta la geometría twoside
% y flota para no chocar). Son contenido semántico: colócalas en el cuerpo,
% cerca del párrafo que comentan, ENTRE párrafos. NO usarlas dentro de
% multicols, tcolorbox, tablas, ni en la misma línea que un \section.
%   \personaje{Nombre}{años · rol}{descripción breve}
%   \hito{año}{qué pasó (breve)}
\setlength{\marginparpush}{8pt}
\newcommand{\margenbox}[1]{\marginpar{%
  \begingroup\raggedright\footnotesize\sffamily
  \setlength{\parskip}{0.18em}\linespread{1.05}\selectfont
  #1\par\endgroup}}
\newcommand{\personaje}[3]{\margenbox{%
  {\color{subjectcolor}\rule{\linewidth}{1.1pt}}\par
  {\InterDisplaySB\scriptsize\color{subjectcolor}#1}\par
  {\fontsize{6.8pt}{8pt}\selectfont\itshape\color{black!50}#2}\par
  \vspace{0.15em}{\scriptsize\color{black!72}#3}}}
\newcommand{\hito}[2]{\margenbox{%
  {\sffamily\bfseries\scriptsize\colorbox{subjectcolor}{\textcolor{white}{\,#1\,}}}\par
  \vspace{0.15em}{\scriptsize\color{black!72}#2}}}

% =====================================================================
%  ENTORNOS — todos planos: filete lateral, sin marco, sin sombra
% =====================================================================

% --- DEFINICIÓN / concepto clave ---------------------------------------
\newtcolorbox{definicion}[1]{
  enhanced, breakable, sharp corners, boxrule=0pt, frame hidden,
  colback=subjectcolor!6!white,
  borderline west={2.2pt}{0pt}{subjectcolor},
  left=5mm, right=4mm, top=3mm, bottom=3mm,
  before skip=\bloqueskip, after skip=\bloqueskip,
  before upper={\blocklabel{subjectcolor}{Definición}{#1}}
}

% --- PROCEDIMIENTO / pasos ---------------------------------------------
\newtcolorbox{procedimiento}[1]{
  enhanced, breakable, sharp corners, boxrule=0pt, frame hidden,
  colback=black!3!white,
  borderline west={2.2pt}{0pt}{black!45},
  left=5mm, right=4mm, top=3mm, bottom=3mm,
  before skip=\bloqueskip, after skip=\bloqueskip,
  before upper={\blocklabel{black!60}{Procedimiento}{#1}}
}
\newcommand{\paso}[2]{%
  \par\vspace{0.45em}\noindent
  \tikz[baseline=-0.6ex]\node[circle, fill=subjectcolor, text=white,
     inner sep=0pt, minimum size=1.45em,
     font=\sffamily\bfseries\footnotesize]{#1};\hspace{0.5em}#2\par}

% Glosa "¿por qué?" bajo un paso — narración en lenguaje cercano
\newcommand{\porque}[1]{\par\vspace{0.1em}\noindent\hspace{2.0em}%
  {\sffamily\footnotesize\itshape\color{subjectcolor}¿por qué?\,}%
  {\sffamily\footnotesize\itshape\color{black!55}#1}\par}

% Operación glosada: cuenta a la izquierda + porqué al margen derecho
% (estilo "comentario al lado del código"). Úsala dentro de ejemplo.
% \opg{ecuación}{porqué}   — primer arg en modo display (mat, cap0)
% \opgt{etiqueta}{porqué}  — primer arg en texto sans bold (física, ciencias)
\newcommand{\opg}[2]{\par\vspace{0.25em}\noindent
  \makebox[0.46\linewidth][l]{$\displaystyle #1$}%
  {\sffamily\footnotesize\itshape\color{black!52}#2}\par}
\newcommand{\opgt}[2]{\par\vspace{0.25em}\noindent
  \makebox[0.46\linewidth][l]{\sffamily\bfseries\color{subjectcolor!80}#1}%
  {\sffamily\footnotesize\itshape\color{black!52}#2}\par}

% --- NOTA / aparte único (sabías · dato · video · margen) --------------
\newtcolorbox{nota}[1][]{
  enhanced, breakable, sharp corners, boxrule=0pt, frame hidden,
  colback=black!4!white,
  borderline west={2.2pt}{0pt}{black!35},
  left=5mm, right=4mm, top=3mm, bottom=3mm,
  before skip=\bloqueskip, after skip=\bloqueskip,
  before upper={\ifblank{#1}{}{\blocklabel{black!55}{Nota}{#1}}}
}

% --- ALERTA / advertencia (única en ámbar, uso escaso) -----------------
\newtcolorbox{alerta}[1]{
  enhanced, breakable, sharp corners, boxrule=0pt, frame hidden,
  colback=alertcolor!7!white,
  borderline west={2.2pt}{0pt}{alertcolor},
  left=5mm, right=4mm, top=3mm, bottom=3mm,
  before skip=\bloqueskip, after skip=\bloqueskip,
  before upper={\blocklabel{alertcolor}{Atención}{#1}}
}
% Sub-ítem de alerta (antes \caso)
\newcommand{\caso}[2]{\par\vspace{0.35em}\noindent
  {\sffamily\bfseries\color{alertcolor!85!black}#1.}\enspace #2\par}

% --- CONTRASTE (la virtud y su exceso → pregunta que cuestiona) --------
%   \contraste{Concepto}{La virtud / lo valioso}{Su exceso / su riesgo}{Pregunta}
\newtcolorbox{contrastebox}{
  enhanced, breakable, sharp corners, boxrule=0pt, frame hidden,
  colback=subjectcolor!5!white,
  borderline west={2.2pt}{0pt}{subjectcolor},
  left=5mm, right=4mm, top=3mm, bottom=3mm,
  before skip=\bloqueskip, after skip=\bloqueskip,
}
\newcommand{\contraste}[4]{%
\begin{contrastebox}
\blocklabel{subjectcolor}{Cuestiona esto}{\faBalanceScale\ #1}
\noindent\begin{minipage}[t]{0.45\linewidth}
  {\sffamily\footnotesize\bfseries\color{subjectcolor!85!black}La virtud}\par
  \vspace{0.2em}#2
\end{minipage}\hfill
{\color{subjectcolor!35}\vrule width 0.4pt}\hfill
\begin{minipage}[t]{0.45\linewidth}
  {\sffamily\footnotesize\bfseries\color{alertcolor!85!black}Su exceso}\par
  \vspace{0.2em}#3
\end{minipage}
\par\vspace{0.6em}
{\color{subjectcolor!25}\titlerule[0.6pt]}\par\vspace{0.4em}
\noindent{\sffamily\itshape\color{subjectcolor}\faArrowRight\ #4}
\end{contrastebox}}

% --- EJEMPLO (sin caja: regla fina + etiqueta numerada) ----------------
\newcounter{ejemplo}[section]
\newenvironment{ejemplo}[1]
  {\par\addvspace{\bloqueskip}\refstepcounter{ejemplo}%
   \noindent{\color{subjectcolor!28}\rule{\linewidth}{0.5pt}}\par\vspace{0.4em}%
   \noindent{\sffamily\bfseries\color{subjectcolor}Ejemplo \theejemplo.}%
   \enspace{\sffamily\bfseries\color{black!78}#1}\par\vspace{0.3em}}
  {\par\addvspace{\bloqueskip}}

% --- PRÁCTICA (ejercicios / hoja / ICFES) ------------------------------
\newtcolorbox{practica}[1]{
  enhanced, breakable, sharp corners, boxrule=0pt, frame hidden,
  colback=white,
  borderline north={1.4pt}{0pt}{subjectcolor},
  left=0mm, right=0mm, top=3mm, bottom=2mm,
  before skip=\bloqueskip, after skip=\bloqueskip,
  before upper={\blocklabel{subjectcolor}{Práctica}{#1}\raggedright}
}

% Ejercicio numerado con nivel opcional [básico|intermedio|reto]
\newcounter{ejer}[section]
\newcommand{\ejer}[2][]{%
  \stepcounter{ejer}\par\vspace{0.5em}\noindent
  {\sffamily\bfseries\color{subjectcolor}\theejer.}\enspace #2%
  \ifblank{#1}{}{\ \textcolor{black!45}{\footnotesize[\textit{#1}]}}\par}

% Niveles = un solo color a tres intensidades (NO nuevos colores)
\newcommand{\nivel}[2]{\par\smallskip\noindent
  \textcolor{subjectcolor!#1}{\small\faCircle}\enspace
  {\sffamily\bfseries\footnotesize\color{black!70}#2}\par}
\newcommand{\basico}{\nivel{40}{Básico}}
\newcommand{\intermedio}{\nivel{70}{Intermedio}}
\newcommand{\reto}{\nivel{100}{Reto}}

% Selección múltiple tipo ICFES — opciones con burbuja (estilo Saber)
\newcounter{pregicfes}[section]
\newcommand{\bubbleletter}[1]{\tikz[baseline=(bc.base)]\node[circle,
  draw=subjectcolor!55, line width=0.6pt, inner sep=1.5pt,
  font=\sffamily\bfseries\scriptsize, text=subjectcolor!90!black](bc){#1};}
\newcommand{\icfes}[5]{%
  \stepcounter{pregicfes}\par\vspace{0.6em}\noindent
  {\sffamily\bfseries\color{subjectcolor}\thepregicfes.}\enspace #1\par
  \begin{enumerate}[label={\protect\bubbleletter{\Alph*}}, leftmargin=1.5cm,
                    itemsep=3pt, topsep=3pt]
    \item #2 \item #3 \item #4 \item #5
  \end{enumerate}}
% Pregunta abierta de argumentación
\newcommand{\pregabierta}[1]{\par\vspace{0.6em}\noindent
  {\sffamily\bfseries\color{subjectcolor}\faComment\ Argumenta:}\enspace
  \textit{#1}\par\lineaescritura\lineaescritura\lineaescritura}

% --- LÍNEAS DE RESPUESTA -----------------------------------------------
\newcommand{\linea}[1][7cm]{\underline{\hspace{#1}}}
\newcommand{\lineaescritura}{\par\vspace{0.9em}\noindent
  {\color{black!30}\rule{\linewidth}{0.3pt}}}

% =====================================================================
%  EVALUACIÓN TIPO ICFES — cierre de unidad (banner + 2 columnas)
%  Uso:  \begin{evaluacion} \begin{multicols}{2} \icfes... \end{multicols}
%        \pregabierta{...} \end{evaluacion}
% =====================================================================
\newtcolorbox{evaluacion}[1][cierre de sección]{
  enhanced, breakable, sharp corners,
  colback=white, colframe=subjectcolor, boxrule=1pt,
  colbacktitle=subjectcolor, coltitle=white,
  fonttitle=\InterDisplaySB\normalsize,
  title={\faClipboardCheck\ \ EVALUACIÓN TIPO ICFES\ \ %
         \textnormal{\sffamily\small· #1}},
  toptitle=2mm, bottomtitle=2mm, left=4mm, right=4mm, top=3.5mm, bottom=3mm,
  before={\clearpage},
  before upper={\small\itshape\color{black!55}%
    Selección múltiple. Marca una sola opción por pregunta.\par\medskip\raggedright}
}

% =====================================================================
%  HOJA DE TRABAJO — recortable: borde punteado, escala de grises,
%  en página aparte (clearpage antes; el contenido que sigue continúa).
% =====================================================================
\newtcolorbox{hoja}[1]{
  enhanced, breakable, sharp corners,
  colback=black!2!white, boxrule=0pt, frame hidden,
  borderline={0.8pt}{0pt}{black!50, dash pattern=on 3pt off 2.5pt},
  before={\clearpage}, after={\par},
  left=6mm, right=6mm, top=6mm, bottom=6mm,
  before upper={%
    {\sffamily\bfseries\footnotesize\color{black!60}\faCut\ \ HOJA DE TRABAJO}%
    {\color{black!35}\ \textbf{·}\ }{\sffamily\bfseries\footnotesize\color{black!80}#1}%
    \hfill{\sffamily\footnotesize\color{black!55}%
      Nombre:\ \underline{\hspace{4.2cm}}\quad Fecha:\ \underline{\hspace{2cm}}}\par
    \vspace{2mm}{\color{black!30}\hrule height 0.4pt}\vspace{3mm}\raggedright}
}

% =====================================================================
%  TABLAHOJA — tabla con TODAS las líneas (rejilla) para hojas de trabajo.
%  Celdas altas para llenar a mano. Centrada. El encabezado se marca con
%  \rowcolor en la primera fila.
%  Uso:  \begin{tablahoja}{ccc} enc & enc & enc \\ a & b & c \\ \end{tablahoja}
% =====================================================================
\NewDocumentEnvironment{tablahoja}{m}{%
  \par\smallskip\centering\small\setlength{\tabcolsep}{4pt}%
  \begin{NiceTabular}{#1}[hvlines, cell-space-top-limit=5pt,
     cell-space-bottom-limit=9pt, colortbl-like]%
}{%
  \end{NiceTabular}\par\smallskip
}

% =====================================================================
%  FIGURAS — imagen real o placeholder automático (\IfFileExists)
% =====================================================================
\newcounter{figura}[section]
\newcommand{\figura}[3][0.7\linewidth]{% [ancho]{archivo}{pie}
  \par\addvspace{\bloqueskip}\refstepcounter{figura}%
  \begin{center}%
  \IfFileExists{assets/imgs/#2}{%
    \includegraphics[width=#1]{assets/imgs/#2}}{%
    \begin{tikzpicture}
      \node[draw=black!30, dashed, line width=0.6pt, fill=black!3,
            text width=0.66\linewidth, minimum height=2.6cm,
            align=center, inner sep=8pt,
            text=black!55, font=\sffamily\footnotesize]
        {{\Large\faImage}\\[4pt]\textbf{[ imagen pendiente ]}\\#3};
    \end{tikzpicture}}%
  \par\vspace{0.4em}%
  {\footnotesize\sffamily\color{black!60}%
    \textbf{Figura \thefigura.}\enspace #3}%
  \end{center}\par\addvspace{\bloqueskip}}

% =====================================================================
%  APERTURA DE SECCIÓN — marco branded de serie + título + "qué verás"
%  El "marco" es CONSTANTE en todas las secciones (= identidad de marca):
%  número grande de serie, kicker, cuadro-firma, franja de acento y sello
%  del instituto. El arte (imagen real o placeholder) va DENTRO del marco.
%  Uso (1ª línea de cada sNN.tex):
%    \aperturaseccion{ruta/img.png}{Título}{Una o dos frases.}
%  Personaliza por documento (opcional):  \renewcommand{\seccionkicker}{...}
% =====================================================================
\newcounter{secportada}
\newcounter{unidadnum}                       % nº de unidad actual (lo fija \aperturacapitulo)
% Número visible de sección = Unidad.Sección (ej. 1.2). Si no hay unidad
% definida (unidadnum=0), cae al número simple de sección.
\newcommand{\numseccion}{\ifnum\value{unidadnum}>0 \theunidadnum.\thesecportada\else\thesecportada\fi}
\newcommand{\seccionkicker}{\hdrcapitulo}   % por defecto, lo del encabezado

% El escenario branded (contenedor constante del arte)
\newtcolorbox{secstage}{%
  enhanced, sharp corners, boxrule=0pt, frame hidden,
  colback=subjectcolor!6, halign=center, valign=center,
  height=5.0cm, left=2mm, right=2mm, top=10mm, bottom=8mm,
  overlay={%
    % cuadro-firma + kicker (arriba-izquierda)
    \fill[subjectcolor] ([shift={(4mm,-4mm)}]frame.north west) rectangle ++(2.6mm,-2.6mm);
    \node[anchor=north west, font=\InterDisplaySB\footnotesize, text=subjectcolor!85]
      at ([shift={(8.5mm,-3.4mm)}]frame.north west) {\seccionkicker};
    % número de serie grande y tenue (arriba-derecha) = marca
    \node[anchor=north east, font=\InterDisplayBlk, text=subjectcolor!18, inner sep=0pt]
      at ([shift={(-4mm,-1.5mm)}]frame.north east)
      {\fontsize{34}{34}\selectfont \numseccion};
    % franja de acento (abajo-izquierda)
    \fill[subjectcolor] ([shift={(4mm,4mm)}]frame.south west) rectangle ++(2.8cm,2.2mm);
    % sello del instituto (abajo-derecha)
    \node[anchor=south east, font=\sffamily\scriptsize, text=subjectcolor!55]
      at ([shift={(-4mm,3.4mm)}]frame.south east) {ITSTZ · Zapatoca};
  }}

\newcommand{\aperturaseccion}[3]{% {archivo}{título}{intro}
  \clearpage\refstepcounter{secportada}%
  \begin{secstage}
  \IfFileExists{assets/imgs/#1}{%
    \includegraphics[width=0.9\linewidth, height=3.2cm, keepaspectratio]{assets/imgs/#1}}{%
    \begin{tikzpicture}
      \node[draw=subjectcolor!35, dashed, line width=0.6pt, fill=white,
            minimum width=0.86\linewidth, minimum height=2.9cm,
            align=center, text=subjectcolor!50, font=\sffamily\small]
        {{\fontsize{22}{22}\selectfont\faImage}\\[5pt]
         \textbf{[ imagen pendiente ]}\\[1pt]
         \scriptsize\ttfamily assets/imgs/#1};
    \end{tikzpicture}}%
  \end{secstage}
  \vspace{0.5em}
  {\color{subjectcolor}\rule{3.2cm}{2.5pt}}\par\vspace{0.3em}
  \section{{\color{subjectcolor}\numseccion}\enspace #2}%
  {\InterLight\fontsize{12.5}{17}\selectfont\color{black!58} #3\par}
  \vspace{0.4em}}

% =====================================================================
%  DESTACADO — "lo esencial": triángulo lateral + degradado muy sutil
% =====================================================================
\newtcolorbox{destacado}[1]{
  enhanced, breakable, sharp corners, boxrule=0pt, frame hidden,
  interior style={left color=subjectcolor!16, right color=subjectcolor!2},
  borderline west={2.4pt}{0pt}{subjectcolor},
  left=7mm, right=5mm, top=3.5mm, bottom=3.5mm,
  before skip=\bloqueskip, after skip=\bloqueskip,
  overlay unbroken and first={%
    \fill[subjectcolor] ([yshift=-4.5mm]frame.north west)
        -- ++(5.5pt,2.75pt) -- ++(0,-5.5pt) -- cycle;},
  before upper={\blocklabel{subjectcolor}{Lo esencial}{#1}}
}

% --- Resaltado inline suave ---
\newtcbox{\resaltar}{on line, nobeforeafter, boxsep=0pt,
  left=2.5pt, right=2.5pt, top=1pt, bottom=1pt, arc=2.5pt,
  colback=subjectcolor!15, colframe=subjectcolor!15}

% =====================================================================
%  EXTENSIONES SANTILLANA (opt-in) — cajas temáticas para dar variedad
%  visual y romper la monotonía. NO alteran los entornos base; úsalas
%  donde aporten. Mantienen el color de asignatura (coherencia de marca).
% =====================================================================

% --- ¿SABÍAS QUE? — curiosidad breve · tarjeta redondeada con bombillo --
\newtcolorbox{sabias}{
  enhanced, breakable, arc=2.5mm, boxrule=0.8pt,
  colback=subjectcolor!6!white, colframe=subjectcolor!30,
  left=4mm, right=4mm, top=3mm, bottom=3mm,
  before skip=\bloqueskip, after skip=\bloqueskip,
  before upper={{\sffamily\bfseries\footnotesize\color{subjectcolor}%
    \faLightbulb\enspace\MakeUppercase{¿Sabías que?}}%
    \par\nopagebreak\vspace{0.45em}}
}

% --- CIENCIA Y SOCIEDAD — impacto/ética · banner de título -------------
\newtcolorbox{cienciasociedad}[1]{
  enhanced, breakable, sharp corners,
  colback=subjectcolor!5!white, colframe=subjectcolor, boxrule=0pt,
  colbacktitle=subjectcolor, coltitle=white,
  fonttitle=\sffamily\bfseries\footnotesize,
  title={\faGlobeAmericas\ \ \MakeUppercase{Ciencia y sociedad}%
         \ifblank{#1}{}{\ \ \textnormal{\sffamily\small· #1}}},
  toptitle=1.6mm, bottomtitle=1.6mm,
  left=4mm, right=4mm, top=3mm, bottom=3mm,
  before skip=\bloqueskip, after skip=\bloqueskip,
}

% --- MANOS A LA OBRA — práctica experimental · marco sólido con matraz --
\newtcolorbox{laboratorio}[1]{
  enhanced, breakable, sharp corners, arc=0pt,
  colback=white, colframe=subjectcolor!55, boxrule=1pt,
  left=4mm, right=4mm, top=3mm, bottom=3mm,
  before skip=\bloqueskip, after skip=\bloqueskip,
  before upper={{\sffamily\bfseries\footnotesize\color{subjectcolor}%
    \faFlask\enspace\MakeUppercase{Manos a la obra}}%
    \ifblank{#1}{}{\ {\color{black!35}\textbf{·}}\ %
      {\sffamily\bfseries\footnotesize\color{black!78}#1}}%
    \par\nopagebreak\vspace{0.45em}}
}
% Lista de materiales dentro de \begin{laboratorio}
\newcommand{\materiales}[1]{\par\noindent
  {\sffamily\bfseries\footnotesize\color{black!70}Materiales:}\ %
  {\footnotesize #1}\par\vspace{0.3em}}

% --- IMAGEN INLINE con texto alrededor (Santillana) -------------------
%   \figinline[r|l]{ancho}{archivo}{pie}   — colócala al inicio de un párrafo
\newcommand{\figinline}[4][r]{%
  \begin{wrapfigure}{#1}{#2}
  \centering\vspace{-0.5\baselineskip}
  \IfFileExists{assets/imgs/#3}{%
    \includegraphics[width=\linewidth]{assets/imgs/#3}}{%
    \fbox{\parbox[c][3cm][c]{0.92\linewidth}{\centering\sffamily\scriptsize
      \color{black!55}{\Large\faImage}\\[3pt][imagen pendiente]\\\ttfamily\tiny #3}}}%
  \par\vspace{2pt}{\footnotesize\sffamily\color{black!60}\textbf{Fig.}\enspace #4}%
  \vspace{-0.4\baselineskip}
  \end{wrapfigure}}

% --- APERTURA CON HERO LATERAL (imagen grande al lado del título) ------
%   \aperturahero{archivo}{título}{intro}  — alternativa a \aperturaseccion
\newcommand{\aperturahero}[3]{%
  \clearpage\refstepcounter{secportada}\stepcounter{section}%
  \addcontentsline{toc}{section}{\numseccion\enspace #2}%
  \noindent\begin{minipage}[c]{0.62\linewidth}
    {\InterDisplaySB\footnotesize\color{subjectcolor!85}\MakeUppercase{\seccionkicker}}\par
    \vspace{1.8mm}
    {\LARGE\InterDisplaySB\color{subjectcolor}%
      {\color{subjectcolor!35}\numseccion}\,\,#2}\par
    \vspace{2mm}{\color{subjectcolor!30}\rule{\linewidth}{0.8pt}}\par\vspace{2mm}
    {\InterLight\fontsize{12}{16}\selectfont\color{black!58}#3}
  \end{minipage}\hfill
  \begin{minipage}[c]{0.34\linewidth}\centering
    \IfFileExists{assets/imgs/#1}{%
      \includegraphics[width=\linewidth]{assets/imgs/#1}}{%
      \fbox{\parbox[c][3.4cm][c]{0.95\linewidth}{\centering\sffamily\scriptsize
        \color{black!55}{\Large\faImage}\\[3pt][imagen pendiente]\\\ttfamily\tiny #1}}}%
  \end{minipage}\par\vspace{1.3em}}

% --- MAPAS (conceptual y mental) — estilos TikZ con texto nítido -------
\tikzset{
  croot/.style={draw=subjectcolor, fill=subjectcolor, text=white, rounded corners=3pt,
    font=\sffamily\bfseries\footnotesize, align=center, inner sep=5pt, minimum height=9mm},
  cnode/.style={draw=subjectcolor!55, fill=subjectcolor!8, rounded corners=2.5pt,
    font=\sffamily\footnotesize, align=center, inner sep=4pt, text width=2.5cm},
  cleaf/.style={draw=subjectcolor!45, fill=white, rounded corners=2.5pt,
    font=\sffamily\footnotesize, align=center, inner sep=4pt},
  clink/.style={-{Stealth[length=5pt]}, subjectcolor!65, semithick},
  cline/.style={subjectcolor!55, semithick},
}
% Caja contenedora para un mapa (título + regla superior)
\newtcolorbox{mapabox}[2][Mapa conceptual]{
  enhanced, sharp corners, boxrule=0pt, frame hidden, colback=white,
  borderline north={1.4pt}{0pt}{subjectcolor},
  left=1mm, right=1mm, top=3mm, bottom=2mm,
  before skip=\bloqueskip, after skip=\bloqueskip,
  before upper={\blocklabel{subjectcolor}{#1}{#2}}
}

% --- Papel cuadriculado para graficar a mano (hojas de trabajo) --------
%   \papelcuadricula[filas]{ancho en cm}   (cada celda = 0.6 cm)
\newcommand{\papelcuadricula}[2][8]{%
  \par\smallskip\begin{center}
  \begin{tikzpicture}
    \draw[step=0.6cm, black!16, line width=0.3pt] (0,0) grid (#2,{#1*0.6});
    \draw[black!45, line width=0.7pt] (0,0) -- (#2,0);
    \draw[black!45, line width=0.7pt] (0,0) -- (0,{#1*0.6});
  \end{tikzpicture}
  \end{center}\par\smallskip}

% =====================================================================
%  APERTURA DE CAPÍTULO — título grande, número marca de agua, aire
% =====================================================================
\newcommand{\aperturacapitulo}[3]{% {número}{título}{subtítulo}
  \setcounter{unidadnum}{#1}\setcounter{secportada}{0}% fija unidad y reinicia secciones
  \thispagestyle{empty}%
  \begin{tikzpicture}[remember picture, overlay]
    \node[anchor=north east, text=subjectcolor!9,
          font=\InterDisplayThin]
      at ([xshift=-0.9cm, yshift=-2.2cm]current page.north east)
      {\fontsize{210}{210}\selectfont #1};
  \end{tikzpicture}%
  \begingroup\raggedright\null\vspace{3.8cm}\par
  {\InterDisplaySB\fontsize{12}{14}\selectfont\color{subjectcolor}%
    CAP\'ITULO\ #1}\par
  \vspace{0.18cm}{\color{black!16}\rule{4cm}{1pt}}\par
  \vspace{0.9cm}
  {\InterDisplaySB\fontsize{46}{50}\selectfont\color{black!88}#2}\par
  \vspace{0.5cm}{\color{subjectcolor}\rule{3.2cm}{2.5pt}}\par
  \vspace{0.55cm}
  {\InterLight\fontsize{15}{21}\selectfont\color{black!55}#3}\par
  \vfill
  {\sffamily\footnotesize\color{black!45}\hdrinstituto\quad
    \textcolor{black!22}{·}\quad Prof.\ Juan Diego A.\ Prada R.}\par
  \endgroup\clearpage}

% =====================================================================
%  DESBLOQUEA — pie discreto, no caja
% =====================================================================
\newcommand{\desbloquea}[1]{%
  \par\addvspace{\bloqueskip}%
  \noindent{\color{subjectcolor!25}\rule{\linewidth}{0.5pt}}\par\vspace{0.4em}%
  {\footnotesize\sffamily\color{black!60}%
    \textcolor{subjectcolor}{\faUnlock}\ \textbf{Dominar esto te desbloquea:}\ #1}\par}

%     % =====================================================================
%  TEMA · CIENCIAS NATURALES — variante del design system (GUÍAS)
%  Identidad visual propia, distinta de Matemáticas (azul/pizarra):
%    · color base ESMERALDA (\setsubject{cnat})
%    · acento LIMA orgánico solo para ornamentos (no rompe la regla de
%      "un solo color de asignatura")
%    · iconografía de ciencias (hoja · ADN · matraz · átomo)
%    · reutiliza las cajas de ciencia que ya trae el preamble:
%      \begin{laboratorio}, \begin{sabias}, \begin{cienciasociedad}
%
%  Sirve a futuro para BIOLOGÍA, QUÍMICA y FÍSICA EXPERIMENTAL.
%
%  INVOCACIÓN (una sola línea, tras \input{design/preamble.tex}):
%     \input{design/tema-cnat.tex}
%  El doc puede seguir personalizando \hdrcapitulo, \seccionkicker, etc.
% =====================================================================

% --- Color base de la asignatura --------------------------------------
\setsubject{cnat}                 % esmeralda RGB(16,110,80)

% Acento LIMA orgánico — SOLO ornamentos (hoja, filete). Nunca contenido:
% la regla "un solo color de asignatura" se mantiene intacta.
\definecolor{cnatlima}{RGB}{124,169,60}

% --- Iconografía semántica de ciencias (fontawesome5, ya cargado) -----
\newcommand{\icobio}{{\color{subjectcolor}\faLeaf}}     % biología
\newcommand{\icoquim}{{\color{subjectcolor}\faFlask}}   % química
\newcommand{\icofis}{{\color{subjectcolor}\faAtom}}     % física
\newcommand{\icogen}{{\color{subjectcolor}\faDna}}      % genética

% --- Kicker temático por defecto (el documento puede renovarlo) -------
\renewcommand{\seccionkicker}{\faLeaf\enspace Ciencias Naturales}

% --- Nombre científico: género + especie en itálica + nombre común ----
%   \especie{Pisum}{sativum}{arveja}  ->  *Pisum sativum* (arveja)
\newcommand{\especie}[3]{\textit{#1 #2}\ (#3)}

% --- Genotipo en línea: preserva mayúsculas/minúsculas de los alelos ---
%   \genotipo{Aa}  (NUNCA aplicar \MakeUppercase: borraría dominante/recesivo)
\newcommand{\genotipo}[1]{{\sffamily\bfseries\color{subjectcolor}#1}}

% --- Filete orgánico con hoja: cierre decorativo opcional de bloque ----
\newcommand{\filetehoja}{%
  \par\vspace{0.4em}\noindent
  {\color{cnatlima}\faLeaf}\hspace{0.3em}%
  {\color{subjectcolor!35}\rule[0.35ex]{0.86\linewidth}{0.7pt}}\par}
     % si es ciencias (verde); mate va en azul
%     % =====================================================================
%  TEMA · MODO TABLERO — "clave del profe" a TODO EL ANCHO
%  Rompe a propósito el layout de libro (twoside angosto): página ancha
%  de una sola columna, banner moderno, imágenes libres al lado con el
%  texto rodeándolas (\figinline), gráficos flat TikZ/SVG y TODO RESUELTO.
%
%  Pensado para material que el PROFE resuelve en el tablero (taller/quiz
%  con solución). NO es el estilo Santillana del kit: es intencionalmente
%  más "revista/póster", más aire, más ancho.
%
%  INVOCACIÓN (tras \input{design/preamble.tex}; opcional tema de asignatura):
%     \input{design/preamble.tex}
%     \input{design/tema-cnat.tex}     % si es ciencias (verde); mate va en azul
%     \input{design/tema-tablero.tex}
% =====================================================================

% --- Geometría a TODO EL ANCHO (simétrica, una cara) -------------------
\geometry{hmargin=1.5cm, top=1.4cm, bottom=1.6cm,
  marginparwidth=0pt, marginparsep=0pt, headsep=12pt, footskip=20pt}

% --- Encabezado / pie limpios (iguales en ambas caras) -----------------
\fancyhf{}
\fancyhead[L]{\footnotesize\sffamily\color{black!50}\hdrinstituto}
\fancyhead[R]{\footnotesize\sffamily\bfseries\color{subjectcolor}\hdrcapitulo}
\renewcommand{\headrule}{\vspace{2pt}{\color{black!15}\hrule height 0.4pt}}
\fancyfoot[R]{\pagenumbox}
\fancyfoot[L]{\footnotesize\sffamily\color{black!45}%
  {\color{subjectcolor}\faCheckCircle}\ \textbf{Clave del profe}\ %
  \textcolor{black!22}{·}\ Prof.\ Juan Diego A.\ Prada R.}

% =====================================================================
%  BANNER de título full-width (barra de color + badge "CLAVE DEL PROFE")
%    \tablerobanner{eyebrow}{Título grande}
% =====================================================================
\newcommand{\tablerobanner}[2]{%
\noindent\begin{tcolorbox}[enhanced, sharp corners, arc=2mm, boxrule=0pt,
   frame hidden, colback=subjectcolor, coltext=white, width=\textwidth,
   left=6mm, right=6mm, top=5mm, bottom=5.5mm,
   overlay={%
     \node[anchor=north east, font=\sffamily\bfseries\footnotesize, text=subjectcolor,
       fill=white, rounded corners=2pt, inner xsep=7pt, inner ysep=3.5pt]
       at ([shift={(-5mm,-4.2mm)}]frame.north east)
       {\faChalkboardTeacher\ \ CLAVE DEL PROFE};}]
  {\sffamily\bfseries\footnotesize\color{white}\MakeUppercase{#1}}\\[3pt]
  {\InterDisplaySB\fontsize{22}{25}\selectfont\color{white} #2}
\end{tcolorbox}\par\vspace{0.4em}}

% =====================================================================
%  ENCABEZADO DE BLOQUE full-width (cuadro de color + título + filete)
%    \bloque{Título del bloque}
% =====================================================================
\newcommand{\bloque}[1]{\par\addvspace{1.15em}\noindent
  {\color{subjectcolor}\rule[-0.02em]{0.85em}{0.85em}}\hspace{0.55em}%
  {\InterDisplaySB\large\color{black!85}#1}\par\vspace{0.25em}%
  {\color{subjectcolor!30}\hrule height 1.1pt}\par\vspace{0.55em}}

% =====================================================================
%  SOLUCIONES (todo resuelto)
% =====================================================================
% Chip de RESULTADO resaltado:  \sol{$(x+3)(x-3)$}
\newtcbox{\sol}{on line, nobeforeafter, boxsep=1pt,
  left=3.5pt, right=3.5pt, top=1.5pt, bottom=1.5pt, arc=3pt,
  colback=subjectcolor!12, colframe=subjectcolor!55,
  coltext=subjectcolor!95!black, fontupper=\sffamily\bfseries}

% Renglón "ejercicio resuelto":  \resuelto{cuerpo math}{caso en gris}
\newcommand{\resuelto}[2]{\par\smallskip\noindent
  {\color{subjectcolor}\faCheckCircle}~ $#1$%
  \ {\footnotesize\sffamily\color{black!45}(#2)}\par}

% Caja de solución paso a paso (filete verde/azul + etiqueta)
\newtcolorbox{solucion}[1][]{
  enhanced, breakable, sharp corners, boxrule=0pt, frame hidden,
  colback=subjectcolor!6!white,
  borderline west={2.4pt}{0pt}{subjectcolor},
  left=5mm, right=4mm, top=3mm, bottom=3mm,
  before skip=\bloqueskip, after skip=\bloqueskip,
  before upper={\blocklabel{subjectcolor}{Solución}{#1}}
}

% Pista / tip corto al margen del razonamiento
\newcommand{\pista}[1]{{\sffamily\footnotesize\itshape\color{black!52}#1}}

% Encabezado de un ejercicio RESUELTO paso a paso:
%   \ejerc{expresión math}{caso/etiqueta}   (luego van los \opg de cada paso)
\newcommand{\ejerc}[2]{\par\addvspace{0.6em}\noindent
  {\color{subjectcolor}\faCheckCircle}\ \ $\displaystyle #1$%
  \ {\footnotesize\sffamily\color{black!45}[#2]}\par\vspace{0.1em}}
% Variante con encabezado de TEXTO (no math): \ejercT{texto}{etiqueta}
\newcommand{\ejercT}[2]{\par\addvspace{0.6em}\noindent
  {\color{subjectcolor}\faCheckCircle}\ \ {\sffamily\bfseries #1}%
  \ {\footnotesize\sffamily\color{black!45}[#2]}\par\vspace{0.1em}}

% =====================================================================

% --- Geometría a TODO EL ANCHO (simétrica, una cara) -------------------
\geometry{hmargin=1.5cm, top=1.4cm, bottom=1.6cm,
  marginparwidth=0pt, marginparsep=0pt, headsep=12pt, footskip=20pt}

% --- Encabezado / pie limpios (iguales en ambas caras) -----------------
\fancyhf{}
\fancyhead[L]{\footnotesize\sffamily\color{black!50}\hdrinstituto}
\fancyhead[R]{\footnotesize\sffamily\bfseries\color{subjectcolor}\hdrcapitulo}
\renewcommand{\headrule}{\vspace{2pt}{\color{black!15}\hrule height 0.4pt}}
\fancyfoot[R]{\pagenumbox}
\fancyfoot[L]{\footnotesize\sffamily\color{black!45}%
  {\color{subjectcolor}\faCheckCircle}\ \textbf{Clave del profe}\ %
  \textcolor{black!22}{·}\ Prof.\ Juan Diego A.\ Prada R.}

% =====================================================================
%  BANNER de título full-width (barra de color + badge "CLAVE DEL PROFE")
%    \tablerobanner{eyebrow}{Título grande}
% =====================================================================
\newcommand{\tablerobanner}[2]{%
\noindent\begin{tcolorbox}[enhanced, sharp corners, arc=2mm, boxrule=0pt,
   frame hidden, colback=subjectcolor, coltext=white, width=\textwidth,
   left=6mm, right=6mm, top=5mm, bottom=5.5mm,
   overlay={%
     \node[anchor=north east, font=\sffamily\bfseries\footnotesize, text=subjectcolor,
       fill=white, rounded corners=2pt, inner xsep=7pt, inner ysep=3.5pt]
       at ([shift={(-5mm,-4.2mm)}]frame.north east)
       {\faChalkboardTeacher\ \ CLAVE DEL PROFE};}]
  {\sffamily\bfseries\footnotesize\color{white}\MakeUppercase{#1}}\\[3pt]
  {\InterDisplaySB\fontsize{22}{25}\selectfont\color{white} #2}
\end{tcolorbox}\par\vspace{0.4em}}

% =====================================================================
%  ENCABEZADO DE BLOQUE full-width (cuadro de color + título + filete)
%    \bloque{Título del bloque}
% =====================================================================
\newcommand{\bloque}[1]{\par\addvspace{1.15em}\noindent
  {\color{subjectcolor}\rule[-0.02em]{0.85em}{0.85em}}\hspace{0.55em}%
  {\InterDisplaySB\large\color{black!85}#1}\par\vspace{0.25em}%
  {\color{subjectcolor!30}\hrule height 1.1pt}\par\vspace{0.55em}}

% =====================================================================
%  SOLUCIONES (todo resuelto)
% =====================================================================
% Chip de RESULTADO resaltado:  \sol{$(x+3)(x-3)$}
\newtcbox{\sol}{on line, nobeforeafter, boxsep=1pt,
  left=3.5pt, right=3.5pt, top=1.5pt, bottom=1.5pt, arc=3pt,
  colback=subjectcolor!12, colframe=subjectcolor!55,
  coltext=subjectcolor!95!black, fontupper=\sffamily\bfseries}

% Renglón "ejercicio resuelto":  \resuelto{cuerpo math}{caso en gris}
\newcommand{\resuelto}[2]{\par\smallskip\noindent
  {\color{subjectcolor}\faCheckCircle}~ $#1$%
  \ {\footnotesize\sffamily\color{black!45}(#2)}\par}

% Caja de solución paso a paso (filete verde/azul + etiqueta)
\newtcolorbox{solucion}[1][]{
  enhanced, breakable, sharp corners, boxrule=0pt, frame hidden,
  colback=subjectcolor!6!white,
  borderline west={2.4pt}{0pt}{subjectcolor},
  left=5mm, right=4mm, top=3mm, bottom=3mm,
  before skip=\bloqueskip, after skip=\bloqueskip,
  before upper={\blocklabel{subjectcolor}{Solución}{#1}}
}

% Pista / tip corto al margen del razonamiento
\newcommand{\pista}[1]{{\sffamily\footnotesize\itshape\color{black!52}#1}}

% Encabezado de un ejercicio RESUELTO paso a paso:
%   \ejerc{expresión math}{caso/etiqueta}   (luego van los \opg de cada paso)
\newcommand{\ejerc}[2]{\par\addvspace{0.6em}\noindent
  {\color{subjectcolor}\faCheckCircle}\ \ $\displaystyle #1$%
  \ {\footnotesize\sffamily\color{black!45}[#2]}\par\vspace{0.1em}}
% Variante con encabezado de TEXTO (no math): \ejercT{texto}{etiqueta}
\newcommand{\ejercT}[2]{\par\addvspace{0.6em}\noindent
  {\color{subjectcolor}\faCheckCircle}\ \ {\sffamily\bfseries #1}%
  \ {\footnotesize\sffamily\color{black!45}[#2]}\par\vspace{0.1em}}

% =====================================================================

% --- Geometría a TODO EL ANCHO (simétrica, una cara) -------------------
\geometry{hmargin=1.5cm, top=1.4cm, bottom=1.6cm,
  marginparwidth=0pt, marginparsep=0pt, headsep=12pt, footskip=20pt}

% --- Encabezado / pie limpios (iguales en ambas caras) -----------------
\fancyhf{}
\fancyhead[L]{\footnotesize\sffamily\color{black!50}\hdrinstituto}
\fancyhead[R]{\footnotesize\sffamily\bfseries\color{subjectcolor}\hdrcapitulo}
\renewcommand{\headrule}{\vspace{2pt}{\color{black!15}\hrule height 0.4pt}}
\fancyfoot[R]{\pagenumbox}
\fancyfoot[L]{\footnotesize\sffamily\color{black!45}%
  {\color{subjectcolor}\faCheckCircle}\ \textbf{Clave del profe}\ %
  \textcolor{black!22}{·}\ Prof.\ Juan Diego A.\ Prada R.}

% =====================================================================
%  BANNER de título full-width (barra de color + badge "CLAVE DEL PROFE")
%    \tablerobanner{eyebrow}{Título grande}
% =====================================================================
\newcommand{\tablerobanner}[2]{%
\noindent\begin{tcolorbox}[enhanced, sharp corners, arc=2mm, boxrule=0pt,
   frame hidden, colback=subjectcolor, coltext=white, width=\textwidth,
   left=6mm, right=6mm, top=5mm, bottom=5.5mm,
   overlay={%
     \node[anchor=north east, font=\sffamily\bfseries\footnotesize, text=subjectcolor,
       fill=white, rounded corners=2pt, inner xsep=7pt, inner ysep=3.5pt]
       at ([shift={(-5mm,-4.2mm)}]frame.north east)
       {\faChalkboardTeacher\ \ CLAVE DEL PROFE};}]
  {\sffamily\bfseries\footnotesize\color{white}\MakeUppercase{#1}}\\[3pt]
  {\InterDisplaySB\fontsize{22}{25}\selectfont\color{white} #2}
\end{tcolorbox}\par\vspace{0.4em}}

% =====================================================================
%  ENCABEZADO DE BLOQUE full-width (cuadro de color + título + filete)
%    \bloque{Título del bloque}
% =====================================================================
\newcommand{\bloque}[1]{\par\addvspace{1.15em}\noindent
  {\color{subjectcolor}\rule[-0.02em]{0.85em}{0.85em}}\hspace{0.55em}%
  {\InterDisplaySB\large\color{black!85}#1}\par\vspace{0.25em}%
  {\color{subjectcolor!30}\hrule height 1.1pt}\par\vspace{0.55em}}

% =====================================================================
%  SOLUCIONES (todo resuelto)
% =====================================================================
% Chip de RESULTADO resaltado:  \sol{$(x+3)(x-3)$}
\newtcbox{\sol}{on line, nobeforeafter, boxsep=1pt,
  left=3.5pt, right=3.5pt, top=1.5pt, bottom=1.5pt, arc=3pt,
  colback=subjectcolor!12, colframe=subjectcolor!55,
  coltext=subjectcolor!95!black, fontupper=\sffamily\bfseries}

% Renglón "ejercicio resuelto":  \resuelto{cuerpo math}{caso en gris}
\newcommand{\resuelto}[2]{\par\smallskip\noindent
  {\color{subjectcolor}\faCheckCircle}~ $#1$%
  \ {\footnotesize\sffamily\color{black!45}(#2)}\par}

% Caja de solución paso a paso (filete verde/azul + etiqueta)
\newtcolorbox{solucion}[1][]{
  enhanced, breakable, sharp corners, boxrule=0pt, frame hidden,
  colback=subjectcolor!6!white,
  borderline west={2.4pt}{0pt}{subjectcolor},
  left=5mm, right=4mm, top=3mm, bottom=3mm,
  before skip=\bloqueskip, after skip=\bloqueskip,
  before upper={\blocklabel{subjectcolor}{Solución}{#1}}
}

% Pista / tip corto al margen del razonamiento
\newcommand{\pista}[1]{{\sffamily\footnotesize\itshape\color{black!52}#1}}

% Encabezado de un ejercicio RESUELTO paso a paso:
%   \ejerc{expresión math}{caso/etiqueta}   (luego van los \opg de cada paso)
\newcommand{\ejerc}[2]{\par\addvspace{0.6em}\noindent
  {\color{subjectcolor}\faCheckCircle}\ \ $\displaystyle #1$%
  \ {\footnotesize\sffamily\color{black!45}[#2]}\par\vspace{0.1em}}
% Variante con encabezado de TEXTO (no math): \ejercT{texto}{etiqueta}
\newcommand{\ejercT}[2]{\par\addvspace{0.6em}\noindent
  {\color{subjectcolor}\faCheckCircle}\ \ {\sffamily\bfseries #1}%
  \ {\footnotesize\sffamily\color{black!45}[#2]}\par\vspace{0.1em}}

% =====================================================================

% --- Geometría a TODO EL ANCHO (simétrica, una cara) -------------------
\geometry{hmargin=1.5cm, top=1.4cm, bottom=1.6cm,
  marginparwidth=0pt, marginparsep=0pt, headsep=12pt, footskip=20pt}

% --- Encabezado / pie limpios (iguales en ambas caras) -----------------
\fancyhf{}
\fancyhead[L]{\footnotesize\sffamily\color{black!50}\hdrinstituto}
\fancyhead[R]{\footnotesize\sffamily\bfseries\color{subjectcolor}\hdrcapitulo}
\renewcommand{\headrule}{\vspace{2pt}{\color{black!15}\hrule height 0.4pt}}
\fancyfoot[R]{\pagenumbox}
\fancyfoot[L]{\footnotesize\sffamily\color{black!45}%
  {\color{subjectcolor}\faCheckCircle}\ \textbf{Clave del profe}\ %
  \textcolor{black!22}{·}\ Prof.\ Juan Diego A.\ Prada R.}

% =====================================================================
%  BANNER de título full-width (barra de color + badge "CLAVE DEL PROFE")
%    \tablerobanner{eyebrow}{Título grande}
% =====================================================================
\newcommand{\tablerobanner}[2]{%
\noindent\begin{tcolorbox}[enhanced, sharp corners, arc=2mm, boxrule=0pt,
   frame hidden, colback=subjectcolor, coltext=white, width=\textwidth,
   left=6mm, right=6mm, top=5mm, bottom=5.5mm,
   overlay={%
     \node[anchor=north east, font=\sffamily\bfseries\footnotesize, text=subjectcolor,
       fill=white, rounded corners=2pt, inner xsep=7pt, inner ysep=3.5pt]
       at ([shift={(-5mm,-4.2mm)}]frame.north east)
       {\faChalkboardTeacher\ \ CLAVE DEL PROFE};}]
  {\sffamily\bfseries\footnotesize\color{white}\MakeUppercase{#1}}\\[3pt]
  {\InterDisplaySB\fontsize{22}{25}\selectfont\color{white} #2}
\end{tcolorbox}\par\vspace{0.4em}}

% =====================================================================
%  ENCABEZADO DE BLOQUE full-width (cuadro de color + título + filete)
%    \bloque{Título del bloque}
% =====================================================================
\newcommand{\bloque}[1]{\par\addvspace{1.15em}\noindent
  {\color{subjectcolor}\rule[-0.02em]{0.85em}{0.85em}}\hspace{0.55em}%
  {\InterDisplaySB\large\color{black!85}#1}\par\vspace{0.25em}%
  {\color{subjectcolor!30}\hrule height 1.1pt}\par\vspace{0.55em}}

% =====================================================================
%  SOLUCIONES (todo resuelto)
% =====================================================================
% Chip de RESULTADO resaltado:  \sol{$(x+3)(x-3)$}
\newtcbox{\sol}{on line, nobeforeafter, boxsep=1pt,
  left=3.5pt, right=3.5pt, top=1.5pt, bottom=1.5pt, arc=3pt,
  colback=subjectcolor!12, colframe=subjectcolor!55,
  coltext=subjectcolor!95!black, fontupper=\sffamily\bfseries}

% Renglón "ejercicio resuelto":  \resuelto{cuerpo math}{caso en gris}
\newcommand{\resuelto}[2]{\par\smallskip\noindent
  {\color{subjectcolor}\faCheckCircle}~ $#1$%
  \ {\footnotesize\sffamily\color{black!45}(#2)}\par}

% Caja de solución paso a paso (filete verde/azul + etiqueta)
\newtcolorbox{solucion}[1][]{
  enhanced, breakable, sharp corners, boxrule=0pt, frame hidden,
  colback=subjectcolor!6!white,
  borderline west={2.4pt}{0pt}{subjectcolor},
  left=5mm, right=4mm, top=3mm, bottom=3mm,
  before skip=\bloqueskip, after skip=\bloqueskip,
  before upper={\blocklabel{subjectcolor}{Solución}{#1}}
}

% Pista / tip corto al margen del razonamiento
\newcommand{\pista}[1]{{\sffamily\footnotesize\itshape\color{black!52}#1}}

% Encabezado de un ejercicio RESUELTO paso a paso:
%   \ejerc{expresión math}{caso/etiqueta}   (luego van los \opg de cada paso)
\newcommand{\ejerc}[2]{\par\addvspace{0.6em}\noindent
  {\color{subjectcolor}\faCheckCircle}\ \ $\displaystyle #1$%
  \ {\footnotesize\sffamily\color{black!45}[#2]}\par\vspace{0.1em}}
% Variante con encabezado de TEXTO (no math): \ejercT{texto}{etiqueta}
\newcommand{\ejercT}[2]{\par\addvspace{0.6em}\noindent
  {\color{subjectcolor}\faCheckCircle}\ \ {\sffamily\bfseries #1}%
  \ {\footnotesize\sffamily\color{black!45}[#2]}\par\vspace{0.1em}}
